% Leçon 30 — Vocabulaire et compréhension de texte : publicité — Chinois Tle A — Cours signé OnBuch+
% Programme officiel MINESEC. Compiler avec : tectonic ZH30-vocabulaire-publicite.tex
\newcommand{\DOCMATIERE}{Chinois}
\newcommand{\DOCNIVEAU}{Tle A}
\newcommand{\DOCMODULE}{Module 4 : Activités économiques et monde du travail}
\newcommand{\DOCLECON}{ZH30}
\newcommand{\DOCTITRE}{Leçon 30 — Vocabulaire et compréhension de texte : publicité}
% OnBuch+ — préambule « cours signés » ENRICHI (compilé avec tectonic / XeLaTeX)
% Système visuel « Pop » d'OnBuch+ : fond crème (jamais blanc), bordures encre
% épaisses, ombres « dures » décalées, titres Archivo Black en capitales.
% Version MATHÉMATIQUES : graphes (pgfplots/tikz), boîtes pédagogiques
% (définition, propriété, méthode, attention, le savais-tu, exercice résolu,
% exercice, TP, bilan), page de garde et signature OnBuch+ ancrée en bas.
%
% Macros attendues dans main.tex AVANT \input{preamble} :
%   \DOCMATIERE  \DOCNIVEAU  \DOCMODULE  \DOCLECON (numéro)  \DOCTITRE
\documentclass[11pt]{article}

\usepackage{fontspec}
\usepackage[french]{babel} % français = langue principale ; le chinois via \zh{..} / textebox (xeCJK)
\usepackage{xeCJK}
\usepackage{pifont} % \ding{51} \ding{55} utilisés par les modèles
\usepackage[a4paper,margin=2cm,top=3.4cm,bottom=2.8cm,headheight=1.4cm,footskip=1.3cm]{geometry}
\usepackage{amsmath,amssymb,amsfonts}
\usepackage{mathtools} % \xmapsto, \coloneqq... souvent utilisés par les modèles
\usepackage{cancel} % simplifications barrées (bilans de forces)
% \abs{z} (module, valeur absolue) et \norm{u} (norme), utilisés par les modèles
\providecommand{\abs}{}\let\abs\relax\DeclarePairedDelimiter{\abs}{\lvert}{\rvert}
\providecommand{\norm}{}\let\norm\relax\DeclarePairedDelimiter{\norm}{\lVert}{\rVert}
\usepackage[table]{xcolor} % [table] : fournit aussi \rowcolors (alternance de lignes)
\usepackage{pagecolor}
\usepackage{tikz}
\usetikzlibrary{arrows.meta,positioning,calc,shapes.geometric,shapes.misc,shapes.arrows,decorations.pathmorphing,decorations.pathreplacing,decorations.markings,patterns,shadows,mindmap,trees,fit,backgrounds,matrix,intersections,angles,quotes}
\usepackage{pgfplots}
\usepackage[version=4]{mhchem} % équations nucléaires \ce{^{235}_{92}U}
\usepackage[european,siunitx]{circuitikz} % schémas électriques
\pgfplotsset{compat=1.18}
\usepgfplotslibrary{fillbetween,groupplots} % aires entre courbes (calcul intégral)
\usepackage{enumitem}
\usepackage{fancyhdr}
% [most] charge toutes les bibliothèques tcolorbox (listings, minted, raster,
% documentation...) et gonfle inutilement la mémoire TeX sur un document long
% (~300 boîtes) : on ne charge que ce qui est réellement utilisé.
\usepackage[skins,breakable]{tcolorbox}
\usepackage{graphicx}
\usepackage{colortbl}
\usepackage{booktabs}
\usepackage{tabularx}
\usepackage{diagbox} % case d'en-tête diagonale des tableaux à double entrée
% Colonnes extensibles centrée / gauche / droite (C, L, R), souvent utilisées par les modèles
\newcolumntype{C}{>{\centering\arraybackslash}X}
\newcolumntype{L}{>{\raggedright\arraybackslash}X}
\newcolumntype{R}{>{\raggedleft\arraybackslash}X}
\usepackage{array}
\usepackage{multicol}
\usepackage{multirow}
\usepackage{siunitx}
% parse-numbers=false : accepte aussi \SI{2,5 \times 10^{-3}}{...}
\sisetup{locale=FR,output-decimal-marker={,},group-minimum-digits=4,parse-numbers=false}
\DeclareSIUnit\ohm{\ensuremath{\symup{\Omega}}} % U+2126 absent de la police
\DeclareSIUnit\division{div} % réglages d'oscilloscope (ms/div, V/div)
\DeclareSIUnit\tr{tr} % tours (vitesse de rotation, ex. \SI{1500}{\tr\per\minute})
\DeclareSIUnit\molar{M} % concentration molaire (ex. \SI{3}{\molar}, \SI{1}{\micro\molar})
\DeclareSIUnit\an{an} % année (le modèle confond parfois avec \year, un compteur TeX) — ex. \SI{5}{\kilo\meter\per\an}
\DeclareSIUnit\FCFA{FCFA} % franc CFA (ex. \SI{500}{\FCFA\per\kilo\gram})
\DeclareSIUnit\degreeBrix{\textdegree Brix} % teneur en sucre d'un sirop/jus (réfractométrie)
\DeclareSIUnit\month{mois} % le modèle utilise parfois \month, un compteur TeX, comme unité de durée
\usepackage{qrcode}
% \degree : commande fréquemment utilisée par les modèles pour un angle ou une
% température en dehors de \SI{}{} (non fournie par les paquets chargés).
\providecommand{\degree}{^{\circ}}
% \qed (fin de démonstration), utilisé par les modèles sans amsthm
\providecommand{\qed}{\hfill$\square$}
% \barre : double barre des valeurs interdites dans un tableau de variations (tkz-tab)
\providecommand{\barre}{\ensuremath{\|}}
% environnement proof (amsthm non chargé) utilisé par les modèles
\ifdefined\proof\else\newenvironment{proof}[1][Démonstration]{\par\noindent\textit{#1.}\ }{\qed\par}\fi
\usepackage{lastpage}
\usepackage{hyperref}
\hypersetup{hidelinks}

% ============================================================
% Palette « Pop » OnBuch+ — mêmes hex que la classe P (plus_ui.dart)
% ============================================================
\definecolor{poporange}{HTML}{F59321}
\definecolor{popdark}{HTML}{E07A0C}
\definecolor{popcream}{HTML}{FFF6E8}
\definecolor{popink}{HTML}{1C1714}
\definecolor{poppurple}{HTML}{7A5AE0}
\definecolor{popgreen}{HTML}{2E9E6A}
\definecolor{poppink}{HTML}{E0457B}
\definecolor{popblue}{HTML}{2F7DE1}
\definecolor{popgold}{HTML}{C98A12}
\definecolor{popmuted}{HTML}{8A7F76}
\definecolor{popline}{HTML}{EADFD2}
\definecolor{popgray}{HTML}{8A7F76}
\colorlet{popgrey}{popgray}
% Alias tolérants (le modèle nomme parfois les couleurs différemment)
\colorlet{popwhite}{white}
\colorlet{popblack}{popink}
\colorlet{popred}{poppink}
\colorlet{popyellow}{popgold}
% Teintes claires (fonds de tableaux/encadrés) — le modèle nomme parfois
% les couleurs avec un suffixe L (ex. popblueL) sans les avoir définies.
\colorlet{poporangeL}{poporange!20}
\colorlet{popdarkL}{popdark!20}
\colorlet{poppurpleL}{poppurple!20}
\colorlet{popgreenL}{popgreen!20}
\colorlet{poppinkL}{poppink!20}
\colorlet{popblueL}{popblue!20}
\colorlet{popgoldL}{popgold!20}
\colorlet{popredL}{popred!20}
\colorlet{popgrayL}{popgray!20}
% Teintes claires pour fonds de tableaux / zones
\colorlet{popblueL}{popblue!10!white}
\colorlet{poporangeL}{poporange!14!white}
\colorlet{popgreenL}{popgreen!12!white}
\colorlet{poppurpleL}{poppurple!12!white}
\colorlet{poppinkL}{poppink!12!white}
\colorlet{popgoldL}{popgold!14!white}

\pagecolor{popcream}

% ============================================================
% Typographie
% ============================================================
\defaultfontfeatures{Ligatures=TeX}
\setmainfont{PlusJakartaSans-Medium.ttf}[Path=../../../fonts/,BoldFont=PlusJakartaSans-Bold.ttf]
% Symboles, API et emojis absents de la police principale : repli sur DejaVu Sans (caractères actifs)
\newfontface\symfont{DejaVuSans.ttf}[Path=../../../fonts/]
\makeatletter
\newcommand\symactive[1]{\catcode#1=13 \begingroup\lccode`\~=#1 \lowercase{\endgroup\def~}{{\symfont\char#1}}}
\newcommand\symblank[1]{\catcode#1=13 \begingroup\lccode`\~=#1 \lowercase{\endgroup\def~}{}}
\newcommand\symhyph[1]{\catcode#1=13 \begingroup\lccode`\~=#1 \lowercase{\endgroup\def~}{\mbox{-}}}
\newcount\symc
\newcommand\symrange[2]{\symc=#1 \@whilenum\symc<#2 \do{\expandafter\symactive\expandafter{\the\symc}\advance\symc by 1 }}
\newcommand\symblankrange[2]{\symc=#1 \@whilenum\symc<#2 \do{\expandafter\symblank\expandafter{\the\symc}\advance\symc by 1 }}
\makeatother
\symrange{"0250}{"02FF}\symactive{"2713}\symactive{"2714}\symactive{"2717}\symactive{"2718}\symactive{"2610}\symactive{"2611}\symactive{"2612}
\symrange{"2600}{"27BF}
\catcode"2705=13 \begingroup\lccode`\~="2705 \lowercase{\endgroup\def~}{{\symfont\char"2713}}\symblank{"26BD}\symblank{"26A1}
\symrange{"2460}{"257F}\symactive{"223C}\symactive{"2248}\symactive{"2260}
\symactive{"01F8}\symactive{"01F9}\symactive{"1E3E}\symactive{"1E3F}
\symhyph{"2011}\symhyph{"2E1A}\symblankrange{"1F300}{"1FAFF}\symblank{"FE0F}
% Caractères chinois : WenQuanYi Zen Hei (sous-ensemble GB2312 + pinyin), gras/italique simulés
\setCJKmainfont{WenQuanYiZenHei-subset.ttf}[Path=../../../fonts/,AutoFakeBold=3,AutoFakeSlant=0.2]
\xeCJKsetup{CJKspace=false}
% Pinyin : voyelles à caron (ǎ ǐ ǒ ǔ ǖ ǘ ǚ ǜ) absentes de la police principale -> caractères actifs
\newfontface\pyfont{WenQuanYiZenHei-subset.ttf}[Path=../../../fonts/]
\newcommand\pyactive[1]{\catcode#1=13 \begingroup\lccode`\~=#1 \lowercase{\endgroup\def~}{{\pyfont\char#1}}}
\pyactive{"01CD}\pyactive{"01CE}\pyactive{"01CF}\pyactive{"01D0}\pyactive{"01D1}\pyactive{"01D2}\pyactive{"01D3}\pyactive{"01D4}\pyactive{"01D5}\pyactive{"01D6}\pyactive{"01D7}\pyactive{"01D8}\pyactive{"01D9}\pyactive{"01DA}\pyactive{"01DB}\pyactive{"01DC}
% Maths en sans-serif (Fira Math) avec chiffres et lettres droites de Plus
% Jakarta, pour que les formules se fondent dans le texte.
\usepackage[math-style=ISO,bold-style=ISO]{unicode-math}
\setmathfont{FiraMath-Regular.otf}[Path=../../../fonts/]
\setmathfont{PlusJakartaSans-Medium.ttf}[Path=../../../fonts/,range={up/num,up/{Latin,latin},\mathup}]
\usepackage{icomma} % virgule décimale sans espace en mode math
% Caractères mathématiques Unicode absents des polices de texte (Plus Jakarta,
% Archivo Black) : rendus par leur équivalent mathématique, en texte comme en
% formule — sinon ils s'affichent en carré vide (ex. « Arithmétique dans ℤ »).
\usepackage{newunicodechar}
\newunicodechar{ℝ}{\ensuremath{\mathbb{R}}}
\newunicodechar{ℤ}{\ensuremath{\mathbb{Z}}}
\newunicodechar{ℂ}{\ensuremath{\mathbb{C}}}
\newunicodechar{ℕ}{\ensuremath{\mathbb{N}}}
\newunicodechar{ℚ}{\ensuremath{\mathbb{Q}}}
\newunicodechar{𝔻}{\ensuremath{\mathbb{D}}}
\newunicodechar{α}{\ensuremath{\alpha}}
\newunicodechar{β}{\ensuremath{\beta}}
\newunicodechar{γ}{\ensuremath{\gamma}}
\newunicodechar{δ}{\ensuremath{\delta}}
\newunicodechar{ε}{\ensuremath{\varepsilon}}
\newunicodechar{θ}{\ensuremath{\theta}}
\newunicodechar{λ}{\ensuremath{\lambda}}
\newunicodechar{μ}{\ensuremath{\mu}}
\newunicodechar{σ}{\ensuremath{\sigma}}
\newunicodechar{τ}{\ensuremath{\tau}}
\newunicodechar{φ}{\ensuremath{\varphi}}
\newunicodechar{ψ}{\ensuremath{\psi}}
\newunicodechar{ω}{\ensuremath{\omega}}
\newunicodechar{Σ}{\ensuremath{\Sigma}}
% majuscules produites par \MakeUppercase dans les titres (α-aminés -> Α)
\newunicodechar{Α}{\ensuremath{\alpha}}
\newunicodechar{Β}{\ensuremath{\beta}}
\newunicodechar{Γ}{\ensuremath{\Gamma}}
\newunicodechar{Δ}{\ensuremath{\Delta}}
\newunicodechar{Ε}{\ensuremath{\varepsilon}}
\newunicodechar{Θ}{\ensuremath{\Theta}}
\newunicodechar{Λ}{\ensuremath{\Lambda}}
\newunicodechar{Μ}{\ensuremath{\mu}}
\newunicodechar{Τ}{\ensuremath{\tau}}
\newunicodechar{Φ}{\ensuremath{\Phi}}
\newunicodechar{Ψ}{\ensuremath{\Psi}}
\newunicodechar{Ω}{\ensuremath{\Omega}}
\newunicodechar{↦}{\ensuremath{\mapsto}}
\newunicodechar{∈}{\ensuremath{\in}}
\newunicodechar{∉}{\ensuremath{\notin}}
\newunicodechar{∧}{\ensuremath{\wedge}}
\newunicodechar{⇔}{\ensuremath{\Leftrightarrow}}
\newunicodechar{⟺}{\ensuremath{\Longleftrightarrow}}
\newunicodechar{⇒}{\ensuremath{\Rightarrow}}
\newunicodechar{⟹}{\ensuremath{\Longrightarrow}}
\newunicodechar{∘}{\ensuremath{\circ}}
\newunicodechar{∪}{\ensuremath{\cup}}
\newunicodechar{∩}{\ensuremath{\cap}}
\newunicodechar{≡}{\ensuremath{\equiv}}
\newunicodechar{⊂}{\ensuremath{\subset}}
\newunicodechar{⊥}{\ensuremath{\perp}}
\newunicodechar{∃}{\ensuremath{\exists}}
\newunicodechar{∀}{\ensuremath{\forall}}
\newunicodechar{∛}{\ensuremath{\sqrt[3]{\,}}}
\newunicodechar{∜}{\ensuremath{\sqrt[4]{\,}}}
\newunicodechar{⃗}{\ensuremath{{}^{\to}}}
\newunicodechar{ⁿ}{\ifmmode^{n}\else\textsuperscript{\ensuremath{n}}\fi}
\newunicodechar{ˣ}{\ifmmode^{x}\else\textsuperscript{\ensuremath{x}}\fi}
\newunicodechar{ᵇ}{\ifmmode^{b}\else\textsuperscript{\ensuremath{b}}\fi}
\newunicodechar{ʳ}{\ifmmode^{r}\else\textsuperscript{\ensuremath{r}}\fi}
\newunicodechar{ᵘ}{\ifmmode^{u}\else\textsuperscript{\ensuremath{u}}\fi}
\newunicodechar{ʸ}{\ifmmode^{y}\else\textsuperscript{\ensuremath{y}}\fi}
\newunicodechar{ᵃ}{\ifmmode^{a}\else\textsuperscript{\ensuremath{a}}\fi}
\newunicodechar{ᵉ}{\ifmmode^{e}\else\textsuperscript{\ensuremath{e}}\fi}
\newunicodechar{ᵏ}{\ifmmode^{k}\else\textsuperscript{\ensuremath{k}}\fi}
\newunicodechar{ᵖ}{\ifmmode^{p}\else\textsuperscript{\ensuremath{p}}\fi}
\newunicodechar{ᴺ}{\ifmmode^{N}\else\textsuperscript{\ensuremath{N}}\fi}
\newunicodechar{ᶜ}{\ifmmode^{c}\else\textsuperscript{\ensuremath{c}}\fi}
\newunicodechar{ʰ}{\ifmmode^{h}\else\textsuperscript{\ensuremath{h}}\fi}
\newunicodechar{ᵠ}{\ifmmode^{q}\else\textsuperscript{\ensuremath{q}}\fi}
\newunicodechar{ₙ}{\ifmmode_{n}\else\textsubscript{\ensuremath{n}}\fi}
\newunicodechar{ₐ}{\ifmmode_{a}\else\textsubscript{\ensuremath{a}}\fi}
\newunicodechar{ᵢ}{\ifmmode_{i}\else\textsubscript{\ensuremath{i}}\fi}
\newunicodechar{ᵣ}{\ifmmode_{r}\else\textsubscript{\ensuremath{r}}\fi}
\newunicodechar{ₖ}{\ifmmode_{k}\else\textsubscript{\ensuremath{k}}\fi}
\newunicodechar{ᵦ}{\ifmmode_{\beta}\else\textsubscript{\ensuremath{\beta}}\fi}
\newunicodechar{ₓ}{\ifmmode_{x}\else\textsubscript{\ensuremath{x}}\fi}
\newfontfamily\popdisplay{ArchivoBlack-Regular.ttf}[Path=../../../fonts/]
\newfontfamily\poplogo{SpaceGrotesk-Bold.ttf}[Path=../../../fonts/]
\newfontfamily\popsemi{PlusJakartaSans-SemiBold.ttf}[Path=../../../fonts/]
\newfontfamily\popxbold{PlusJakartaSans-ExtraBold.ttf}[Path=../../../fonts/]
\newfontfamily\popxboldls{PlusJakartaSans-ExtraBold.ttf}[Path=../../../fonts/,LetterSpace=3.0]

\setlist[enumerate]{leftmargin=1.7em,itemsep=5pt,topsep=4pt}
\setlist[itemize]{leftmargin=1.7em,itemsep=4pt,topsep=4pt,label={\color{poporange}\textbullet}}
\setlength{\parindent}{0pt}
\setlength{\parskip}{5pt}
\renewcommand{\arraystretch}{1.25}

% pgfplots : style Pop par défaut
\pgfplotsset{
  every axis/.append style={axis line style={popink,line width=1pt},tick style={popink},
    grid style={popline,line width=0.5pt},label style={font=\small},tick label style={font=\footnotesize}},
  popcurve/.style={poporange,line width=1.8pt,no markers,smooth},
}
% Même style disponible dans un \draw TikZ simple (hors pgfplots)
\tikzset{popcurve/.style={poporange,line width=1.8pt}}
\tikzset{popgrid/.style={popline,very thin}}
\colorlet{popcurve}{poporange} % le nom du style est parfois utilisé comme couleur

% ============================================================
% Pill / squiggle
% ============================================================
\newcommand{\poppill}[3]{%
  \tikz[baseline=-0.55ex]\node[draw=popink,line width=1.2pt,rounded corners=2.6pt,%
    fill=#2,inner xsep=6.5pt,inner ysep=3pt]{%
    \popxboldls\fontsize{7}{7}\selectfont\color{#3}\MakeUppercase{#1}};%
}
\newcommand{\popsquiggle}[1]{%
  \tikz[baseline=-0.4ex]{
    \draw[#1,line width=2.6pt,line cap=round]
      plot [smooth] coordinates {(0,0.09) (0.34,0.01) (0.68,0.21) (1.02,0.01) (1.36,0.10)};
  }%
}
% Mot-clé surligné (marqueur orange sous le texte)
\newcommand{\cle}[1]{{\popxbold\color{popdark}#1}}

% ============================================================
% En-tête / pied
% ============================================================
\pagestyle{fancy}
\fancyhf{}
\renewcommand{\headrulewidth}{0pt}
\renewcommand{\footrulewidth}{0pt}
\lhead{\raisebox{-0.15ex}{\poplogo\fontsize{13}{13}\selectfont\color{popink}OnBuch\color{poporange}+}}
\rhead{\poppill{\DOCMATIERE\ \textbullet\ \DOCNIVEAU\ \textbullet\ Leçon \DOCLECON}{white}{popink}}
\lfoot{\popxbold\fontsize{7.6}{7.6}\selectfont\color{popmuted}\MakeUppercase{Cours signé OnBuch+}\ \textbullet\ {\popsemi\DOCTITRE}}
\rfoot{\tikz[baseline=-0.6ex]\node[rounded corners=8.5pt,fill=popink,minimum height=17pt,minimum width=17pt,inner xsep=5pt,inner sep=0]{\popxbold\fontsize{8}{8}\selectfont\color{popcream}\thepage\,/\,\pageref*{LastPage}};}

% ============================================================
% Titres
% ============================================================
\newcommand{\chaptitre}[2]{%
  \begin{tcolorbox}[enhanced,colback=poporange,colframe=popink,boxrule=2.6pt,arc=11pt,
      drop shadow=popink,left=15pt,right=15pt,top=12pt,bottom=11pt,before skip=3pt,after skip=18pt]
    \poppill{#2}{popink}{popcream}\par\vspace{9pt}
    {\raggedright\hyphenpenalty=10000\exhyphenpenalty=10000\popdisplay\fontsize{22}{24}\selectfont\color{popcream}\MakeUppercase{#1}\par}
  \end{tcolorbox}
}
% Section numérotée : pastille orange avec le numéro + titre Archivo
\newcounter{popsec}
\newcommand{\coursec}[1]{%
  \stepcounter{popsec}\setcounter{popsub}{0}%
  \par\vspace{16pt}\noindent
  \begin{minipage}{\linewidth}
  \tikz[baseline=-0.7ex]\node[draw=popink,line width=1.6pt,fill=poporange,rounded corners=4pt,minimum size=22pt,inner sep=0]{\popdisplay\fontsize{12}{12}\selectfont\color{popcream}\thepopsec};%
  \hspace{8pt}{\popdisplay\fontsize{15}{17}\selectfont\color{popink}\MakeUppercase{#1}}\par
  \vspace{4pt}\noindent{\color{poporange}\rule{36pt}{3.6pt}}
  \end{minipage}\par\nopagebreak\vspace{8pt}
}
\newcounter{popsub}
\newcommand{\courssub}[1]{%
  \stepcounter{popsub}%
  \par\vspace{9pt}\noindent{\popxbold\fontsize{11.5}{13}\selectfont\color{popink}{\color{poporange}\thepopsec.\thepopsub}\hspace{6pt}#1}\par\nopagebreak\vspace{4pt}
}

% ============================================================
% Boîtes pédagogiques — même recette : fond blanc, bordure épaisse colorée,
% ombre dure encre, badge plein. Titre optionnel [#1].
% ============================================================
\tcbset{popbase/.style={breakable,enhanced,colback=white,boxrule=2.3pt,arc=9pt,
  shadow={3pt}{-3pt}{0pt}{popink},left=14pt,right=14pt,top=11pt,bottom=11pt,before skip=11pt,after skip=15pt,
  fontupper=\popsemi\fontsize{10.6}{14.5}\selectfont\color{popink},
  fontlower=\popsemi\fontsize{10.6}{14.5}\selectfont\color{popink}}}
% Coche dessinée (le glyphe ✓ n'existe pas dans les polices du document)
\newcommand{\popcheck}[1]{\tikz[baseline=-0.5ex]\draw[#1,line width=1.6pt,line cap=round,line join=round](-0.13,0.0)--(-0.04,-0.09)--(0.13,0.1);}
\providecommand{\coche}{\popcheck{popgreen}}
\providecommand{\resultat}[1]{\par\smallskip\noindent\fcolorbox{popgreen}{popgreenL}{\parbox{\dimexpr\linewidth-2\fboxsep-2\fboxrule\relax}{\textbf{Résultat.} #1}}\par\smallskip}
\newcommand{\popboxhead}[3]{\poppill{#1}{#2}{white}\ifx\relax#3\relax\else\hspace{6pt}{\popxbold\fontsize{10.6}{12}\selectfont\color{popink}#3}\fi\par\vspace{7pt}}

\newtcolorbox{aretenir}[1][]{popbase,colframe=popgreen,before upper={\popboxhead{À retenir}{popgreen}{#1}}}
% Texte en chinois (document de lecture, dialogue, extrait) : cadre
% d'étude. Un titre optionnel (source, auteur) s'affiche dans l'en-tête.
\newtcolorbox{textebox}[1][]{popbase,colframe=popblue,before upper={\popboxhead{文本 · Texte}{popblue}{#1}},after upper={}}
% Marques ✓ / ✗ (forme correcte / fautive) souvent utilisées par les modèles
\providecommand{\cross}{\textcolor{poppink}{\ensuremath{\times}}}
\providecommand{\xmark}{\cross}\providecommand{\crossmark}{\cross}\providecommand{\cmark}{\ensuremath{\checkmark}}
% Ligne de réponse / justification à compléter (inventée par les modèles)
\providecommand{\justification}[1]{\par\noindent\textit{Justification~:} \dotfill\par}
% \textipa (package tipa non chargé) : la police ne couvre pas l'API -> simple passage
\providecommand{\textipa}[1]{#1}
% Soulignement (ulem non chargé) : \uline, \ul
\providecommand{\uline}[1]{\underline{#1}}\providecommand{\ul}[1]{\underline{#1}}
% \glossaire{\item ...} inventée par les modèles : liste de glossaire
\providecommand{\glossaire}[1]{\begin{itemize}#1\end{itemize}}
% \alert (beamer) inventée par les modèles : texte mis en évidence
\providecommand{\alert}[1]{\textbf{\textcolor{poppink}{#1}}}
% variantes ASCII d'ordinaux écrites en macro
\providecommand{\iere}{\textsuperscript{re}}\providecommand{\ieme}{\textsuperscript{e}}\providecommand{\ere}{\textsuperscript{re}}\providecommand{\eme}{\textsuperscript{e}}\providecommand{\er}{\textsuperscript{er}}\providecommand{\ie}{\textsuperscript{e}}
% Ordinaux français écrits en macro par les modèles : 1\ière, 3\ième
\newcommand{\ière}{\textsuperscript{re}}\newcommand{\ième}{\textsuperscript{e}}
% \exemple (inventée par les modèles) : étiquette « Exemple : »
\providecommand{\exemple}{\textbf{Exemple~:}\ }
% \square utilisable aussi hors mode math (cases à cocher)
\AtBeginDocument{\let\squareMath\square\renewcommand{\square}{\ensuremath{\squareMath}}}
% Mot ou phrase en chinois à l'intérieur d'un paragraphe français
\newcommand{\zh}[1]{\ifmmode\text{#1}\else{#1}\fi}
\let\eng\zh \let\esp\zh \let\all\zh % alias tolérés
% flèches d'intonation inventées par les modèles
\providecommand{\textsearrow}{}\renewcommand{\textsearrow}{\ensuremath{\searrow}}\providecommand{\textnearrow}{}\renewcommand{\textnearrow}{\ensuremath{\nearrow}}
% \fr{..} : mot français (inventé par les modèles)
\providecommand{\fr}[1]{\textit{#1}}
% zhbox : encadré de texte chinois inventé par les modèles
\newenvironment{zhbox}{\begin{quote}}{\end{quote}}
% \courssubsub : titre de niveau 3 inventé par les modèles
\providecommand{\courssubsub}[1]{\par\medskip\noindent\textbf{#1}\par\smallskip}
% \zhbf : texte chinois en gras inventé par les modèles
\providecommand{\zhbf}[1]{\textbf{#1}}
% \vs : « versus » inventé par les modèles
\providecommand{\vs}{\textit{vs}\ }
% \blank : case à compléter inventée par les modèles
\providecommand{\blank}[1][]{\underline{\hspace{1.6em}}}
\newtcolorbox{exemplebox}[1][]{popbase,colframe=poporange,before upper={\popboxhead{Exemple}{poporange}{#1}}}
\newtcolorbox{definition}[1][]{popbase,colframe=popblue,before upper={\popboxhead{Définition}{popblue}{#1}}}
\newtcolorbox{propriete}[1][]{popbase,colframe=poppurple,before upper={\popboxhead{Propriété}{poppurple}{#1}}}
\newtcolorbox{methode}[1][]{popbase,colframe=popgold,before upper={\popboxhead{Méthode}{popgold}{#1}}}
\newtcolorbox{attention}[1][]{popbase,colframe=poppink,before upper={\popboxhead{Attention}{poppink}{#1}}}
\newtcolorbox{experience}[1][]{popbase,colframe=popblue,colback=popblueL,before upper={\popboxhead{Expérience}{popblue}{#1}}}
% « Le savais-tu ? » : carte violette pleine (contexte, culture, Cameroun)
\newtcolorbox{savaistu}[1][]{popbase,colback=poppurple,colframe=popink,
  fontupper=\popsemi\fontsize{10.4}{14.2}\selectfont\color{white},
  before upper={\poppill{Le savais-tu\,?}{white}{poppurple}\ifx\relax#1\relax\else\hspace{6pt}{\popxbold\color{white}#1}\fi\par\vspace{7pt}}}
% Exercice résolu : énoncé en haut, \tcblower puis la solution sur fond crème
\newcounter{popexo}\newcounter{popexr}
\newtcolorbox{exoresolu}[1][]{popbase,colframe=poporange,colbacklower=popcream,
  segmentation style={popink,line width=1.2pt,dashed},
  before upper={\stepcounter{popexr}\popboxhead{Exercice résolu \thepopexr}{poporange}{#1}},
  before lower={\poppill{Solution}{popink}{popcream}\par\vspace{6pt}}}
% Exercice (non résolu en place) — la correction est dans la partie Corrigés
\newtcolorbox{exercice}[1][]{popbase,colframe=popink,
  before upper={\stepcounter{popexo}\popboxhead{Exercice \thepopexo}{popink}{#1}}}
% Corrigé
\newtcolorbox{corrige}[1][]{popbase,colframe=popgreen,colback=popgreenL,
  before upper={\popboxhead{Corrigé}{popgreen}{#1}}}
% Objectifs / prérequis / bilan
\newtcolorbox{objectifs}{popbase,colframe=popink,colback=poporangeL,
  before upper={\popboxhead{Objectifs de la leçon}{popink}{}}}
\newtcolorbox{prerequis}{popbase,colframe=popink,colback=popblueL,
  before upper={\popboxhead{Prérequis}{popblue}{}}}
\newtcolorbox{bilan}{popbase,colframe=popink,colback=white,boxrule=2.8pt,
  before upper={\popboxhead{Fiche bilan}{poporange}{L'essentiel à connaître pour l'examen}}}
% Figure encadrée avec légende
\newtcolorbox{popfigure}[1][]{enhanced,breakable=false,colback=white,colframe=popline,boxrule=1.4pt,arc=8pt,
  left=10pt,right=10pt,top=10pt,bottom=8pt,before skip=10pt,after skip=12pt,halign=center,
  lower separated=false,halign lower=center,
  fontlower=\popsemi\fontsize{9}{11.5}\selectfont\color{popmuted},
  #1}
\newcommand{\legende}[1]{\par\vspace{6pt}{\popsemi\fontsize{9}{11.5}\selectfont\color{popmuted}#1}}

% ============================================================
% Signature OnBuch+
% ============================================================
\newcommand{\poplogoinline}[1]{{\poplogo\fontsize{#1}{#1}\selectfont\color{popink}OnBuch\color{poporange}+}}
\newcommand{\signatureonbuch}{%
  \begin{tcolorbox}[enhanced,colback=popink,colframe=popink,boxrule=2.2pt,arc=10pt,
      drop shadow=poporange,left=14pt,right=14pt,top=10pt,bottom=10pt,before skip=0pt,after skip=0pt]
    \tikz[baseline=-0.7ex]\node[circle,fill=popgreen,draw=popcream,line width=1.3pt,minimum size=22pt,inner sep=0]{\popcheck{white}};%
    \hspace{9pt}\begin{minipage}[c]{0.62\linewidth}
      {\popxbold\fontsize{12}{13}\selectfont\color{popcream}Signé\ \textbullet\ {\poplogo OnBuch\color{poporange}+}}\par\vspace{2pt}
      {\raggedright\popsemi\fontsize{8}{10.5}\selectfont\color{popline}\MakeUppercase{Équipe pédagogique OnBuch+}\\\MakeUppercase{Conforme au programme officiel MINESEC}\par}
    \end{minipage}\hfill
    {\poplogo\fontsize{22}{22}\selectfont\color{popcream}OnBuch\color{poporange}+}
  \end{tcolorbox}%
}

% ============================================================
% Page de garde
% #1 titre · #2 matière · #3 niveau · #4 module · #5 n° leçon · #6 durée ·
% #7 description / objectifs (texte ou itemize)
% ============================================================
\newcommand{\pagegarde}[7]{%
  \thispagestyle{empty}
  \newgeometry{a4paper,margin=1.7cm,top=1.6cm,bottom=1.4cm}%
  \begin{tikzpicture}[remember picture,overlay]
    \fill[poporange!18] ([xshift=-3.2cm,yshift=-3.6cm]current page.north east) circle (4.6cm);
    \fill[poppurple!14] ([xshift=2.4cm,yshift=7.6cm]current page.south west) circle (3.8cm);
  \end{tikzpicture}%
  {\poplogo\fontsize{30}{30}\selectfont\color{popink}OnBuch\color{poporange}+}\hfill
  \raisebox{0.6ex}{\poppill{Cours signé \textbullet\ Premium}{popink}{popcream}}\par
  \vspace{1pt}\popsquiggle{poppurple}\par
  \vspace{8pt}
  \begin{tcolorbox}[enhanced,colback=poporange,colframe=popink,boxrule=2.8pt,arc=13pt,
      drop shadow=popink,left=18pt,right=18pt,top=12pt,bottom=13pt,after skip=10pt]
    \poppill{#2 \textbullet\ #3}{popink}{popcream}\hspace{5pt}\poppill{#4}{white}{popink}\par\vspace{8pt}
    {\popdisplay\fontsize{14}{14}\selectfont\color{popink}LEÇON #5}\par\vspace{4pt}
    {\raggedright\hyphenpenalty=10000\exhyphenpenalty=10000\popdisplay\fontsize{25}{27}\selectfont\color{popcream}\MakeUppercase{#1}\par}
    \vspace{8pt}
    {\popxbold\fontsize{9.5}{9.5}\selectfont\color{popink}\MakeUppercase{Durée conseillée : #6}}
  \end{tcolorbox}
  \begin{tcolorbox}[enhanced,colback=white,colframe=popink,boxrule=2.2pt,arc=10pt,
      drop shadow=popink,left=16pt,right=16pt,top=10pt,bottom=10pt,after skip=8pt]
    \poppill{Dans cette leçon}{poppurple}{white}\par\vspace{5pt}
    \popsemi\fontsize{9.3}{12.6}\selectfont\color{popink}#7
  \end{tcolorbox}
  \noindent\begin{minipage}[t]{0.487\linewidth}
    \begin{tcolorbox}[enhanced,colback=popgreen,colframe=popink,boxrule=2.2pt,arc=10pt,
        drop shadow=popink,left=11pt,right=11pt,top=10pt,bottom=10pt,height=3.1cm,valign=center]
      \tcbox[colback=white,colframe=popink,boxrule=1.3pt,arc=5pt,boxsep=0pt,nobeforeafter,
        left=3pt,right=3pt,top=3pt,bottom=3pt,valign=center]{\qrcode[height=1.9cm]{https://play.google.com/store/apps/details?id=com.luvvix.onbuch&hl=fr}}%
      \hspace{8pt}\begin{minipage}[c]{0.5\linewidth}
        {\popdisplay\fontsize{11}{12.5}\selectfont\color{white}\MakeUppercase{Télécharge OnBuch}}\par\vspace{4pt}
        {\raggedright\popsemi\fontsize{8.4}{11}\selectfont\color{white}Gratuit \textbullet\ Google Play\\Tuteur IA, annales, concours, résultats\par}
      \end{minipage}
    \end{tcolorbox}
  \end{minipage}\hfill
  \begin{minipage}[t]{0.487\linewidth}
    \begin{tcolorbox}[enhanced,colback=poppurple,colframe=popink,boxrule=2.2pt,arc=10pt,
        drop shadow=popink,left=11pt,right=11pt,top=10pt,bottom=10pt,height=3.1cm,valign=center]
      \tikz[baseline=-0.7ex]\node[circle,fill=white,draw=popink,line width=1.3pt,minimum size=1.6cm,inner sep=0]{\popdisplay\fontsize{24}{24}\selectfont\color{poppurple}+};%
      \hspace{8pt}\begin{minipage}[c]{0.6\linewidth}
        {\popdisplay\fontsize{11}{12.5}\selectfont\color{white}\MakeUppercase{Passe à OnBuch+}}\par\vspace{4pt}
        {\raggedright\popsemi\fontsize{8.4}{11}\selectfont\color{white}Tous les cours signés, Tuteur IA illimité, fiches PDF — onglet OnBuch+ de l'app\par}
      \end{minipage}
    \end{tcolorbox}
  \end{minipage}
  \vspace{8pt}
  \popcontenu
  \vfill
  \signatureonbuch
  \restoregeometry
  \newpage
}

% Rangée « Ce cours contient » (page de garde)
\newcommand{\poptile}[3]{% #1 couleur #2 gros texte #3 légende
  \begin{tcolorbox}[enhanced,colback=#1,colframe=popink,boxrule=2pt,arc=9pt,shadow={2.5pt}{-2.5pt}{0pt}{popink},
      left=6pt,right=6pt,top=9pt,bottom=9pt,halign=center,valign=center,height=1.75cm,width=\linewidth,nobeforeafter]
    {\popdisplay\fontsize{12.5}{13}\selectfont\color{white}\MakeUppercase{#2}}\par\vspace{3pt}
    {\centering\popsemi\fontsize{7.8}{9.5}\selectfont\color{white}#3\par}
  \end{tcolorbox}}
\newcommand{\popcontenu}{%
  \par\noindent{\popxboldls\fontsize{7.5}{7.5}\selectfont\color{popmuted}CE COURS SIGNÉ CONTIENT}\par\vspace{7pt}
  \noindent\begin{minipage}[t]{0.235\linewidth}\poptile{popblue}{Cours}{complet, illustré, pas à pas}\end{minipage}\hfill
  \begin{minipage}[t]{0.235\linewidth}\poptile{poporange}{Méthodes}{et exercices résolus}\end{minipage}\hfill
  \begin{minipage}[t]{0.235\linewidth}\poptile{poppink}{Exercices}{type examen + corrigés}\end{minipage}\hfill
  \begin{minipage}[t]{0.235\linewidth}\poptile{popgreen}{Fiche bilan}{l'essentiel à retenir}\end{minipage}\par}

% Sommaire
\newtcolorbox{sommaire}{enhanced,colback=white,colframe=popink,boxrule=2.2pt,arc=10pt,
  drop shadow=popink,left=18pt,right=18pt,top=13pt,bottom=13pt,before skip=6pt,after skip=20pt,
  before upper={\poppill{Sommaire}{poppurple}{white}\par\vspace{9pt}\popsemi\fontsize{10.6}{14.5}\selectfont\color{popink}}}

% Signature de fin de cours — poussée tout en bas de la dernière page
\newcommand{\findecours}{%
  \par\vfill\nopagebreak
  \begin{center}\popsquiggle{poporange}\end{center}\vspace{4pt}
  \signatureonbuch
}
\AtBeginDocument{\def\llbracket{\mathopen{[\mkern-3.5mu[}}\def\rrbracket{\mathclose{]\mkern-3.5mu]}}} % intervalles d'entiers (glyphe ⟦ absent de la police)
% Glyphes absents de Fira Math : redéfinis à partir de glyphes disponibles
\AtBeginDocument{%
  \def\setminus{\mathbin{\backslash}}\let\smallsetminus\setminus
  \def\vdots{\mathord{\vbox{\baselineskip3.2pt\lineskiplimit0pt\kern4pt\hbox{.}\hbox{.}\hbox{.}}}}%
  \def\ddots{\mathinner{\mkern1mu\raise6.5pt\hbox{.}\mkern2mu\raise3.6pt\hbox{.}\mkern2mu\raise0.7pt\hbox{.}\mkern1mu}}%
  \def\triangle{\mathord{\scalebox{0.9}{$\symup{\Delta}$}}}%
  \def\nabla{\mathord{\raisebox{\depth}{\scalebox{1}[-1]{$\symup{\Delta}$}}}}%
  \def\checkmark{\popcheck{popdark}}%
  \def\star{\mathbin{\ast}}\def\bigstar{\mathord{\ast}}%
  \def\diamond{\mathbin{\scalebox{0.6}{$\Diamond$}}}%
  \def\bigcup{\mathop{\mathchoice{\raisebox{-0.3em}{\scalebox{1.7}{$\cup$}}}{\scalebox{1.25}{$\cup$}}{\cup}{\cup}}\displaylimits}%
  \def\bigcap{\mathop{\mathchoice{\raisebox{-0.3em}{\scalebox{1.7}{$\cap$}}}{\scalebox{1.25}{$\cap$}}{\cap}{\cap}}\displaylimits}%
  \def\lesseqqgtr{\mathrel{\lessgtr}}\def\bigoplus{\mathop{\oplus}\displaylimits}%
}
\providecommand{\ssi}{si et seulement si} % abréviation fréquente
\AtBeginDocument{\providecommand{\wideparen}{\overparen}} % arc de cercle

%% FIGFIX-BEGIN v5 — correctifs globaux des figures (docs/onbuch-generation/tools/figfix)
\usepackage{adjustbox}\usepackage{etoolbox}\tcbuselibrary{raster}
% toute figure TikZ/pgfplots plus large que la ligne est réduite à la largeur disponible
\BeforeBeginEnvironment{tikzpicture}{\begin{adjustbox}{max width=\linewidth}}
\AfterEndEnvironment{tikzpicture}{\end{adjustbox}}
% légendes pgfplots lisibles par-dessus les courbes
\pgfplotsset{every axis/.append style={legend style={fill=white,fill opacity=0.92,text opacity=1,draw=black!15,inner sep=3pt}}}
% étiquettes d'arêtes (syntaxe edge["..."]) avec fond blanc
\tikzset{every edge quotes/.append style={fill=white,inner sep=1.2pt}}
%% FIGFIX-END
\begin{document}

\pagegarde{Leçon 30 — Vocabulaire et compréhension de texte : publicité}{Chinois}{Tle A}{Module 4 : Activités économiques et monde du travail}{ZH30}{2 h}{Cette leçon du module « Activités économiques et monde du travail » entraîne les élèves de Terminale A à lire et comprendre un texte publicitaire chinois authentique et simplifié. Elle consolide le vocabulaire de la consommation et de la promotion, l'étude des radicaux et la réponse à des questions de compréhension de type examen.
\vspace{4pt}
\begin{multicols}{2}\small
\begin{itemize}
\item À la fin de cette leçon, je sais lire à voix haute un court texte publicitaire chinois en respectant le pinyin et les tons.
\item Je sais reconnaître, prononcer et traduire au moins 15 mots de vocabulaire liés à la publicité et à la consommation.
\item Je sais identifier les clés (radicaux) et tracer l'ordre des traits des caractères clés du thème.
\item Je sais repérer dans un texte publicitaire le produit, le prix, la réduction et les arguments de vente.
\item Je sais répondre par écrit à des questions de compréhension en phrases complètes.
\item Je sais employer les collocations fréquentes du domaine (faire de la publicité, soldes, qualité, prix).
\end{itemize}\end{multicols}}

\begin{sommaire}
\begin{multicols}{2}
\begin{enumerate}
\item Mise en route et découverte du thème
\item Texte support : lecture et écoute
\item Glossaire du thème
\item Clés (radicaux) et ordre des traits
\item Compréhension du texte et collocations
\item Production guidée : ma publicité
\item Activité d'intégration
\item Méthodes et exercices résolus
\item Exercices
\item Corrigés des exercices
\item Fiche bilan
\end{enumerate}
\end{multicols}
\end{sommaire}

\begin{objectifs}
\begin{itemize}
\item À la fin de cette leçon, je sais lire à voix haute un court texte publicitaire chinois en respectant le pinyin et les tons.
\item Je sais reconnaître, prononcer et traduire au moins 15 mots de vocabulaire liés à la publicité et à la consommation.
\item Je sais identifier les clés (radicaux) et tracer l'ordre des traits des caractères clés du thème.
\item Je sais repérer dans un texte publicitaire le produit, le prix, la réduction et les arguments de vente.
\item Je sais répondre par écrit à des questions de compréhension en phrases complètes.
\item Je sais employer les collocations fréquentes du domaine (faire de la publicité, soldes, qualité, prix).
\item Je sais rédiger une mini-annonce publicitaire guidée de 4 à 6 phrases en chinois.
\item Je sais comparer brièvement les pratiques publicitaires au Cameroun et en Chine.
\end{itemize}
\end{objectifs}

\begin{prerequis}
\begin{itemize}
\item Vocabulaire des achats et des prix (leçon sur les courses : 买, 卖, 钱, 块, 便宜, 贵)
\item Nombres, mesures et expression des prix en chinois
\item Structure de base sujet-verbe-complément et particule 的
\item Lecture du pinyin et maîtrise des quatre tons
\item Notions sur les magasins et le commerce vues en Première
\end{itemize}
\end{prerequis}

\newpage

\coursec{Mise en route et découverte du thème}

\courssub{Situation d'accroche et problématique}

En marchant dans les rues de Douala ou de Yaoundé, tu es entouré de publicités : affiches pour le crédit téléphonique, spots radio pour une marque de savon, panneaux pour les universités, banderoles au marché Mokolo ou à Deïdo. La publicité fait partie de notre quotidien, au Cameroun comme partout dans le monde.

Imagine maintenant qu'une entreprise camerounaise veuille exporter son jus de bissap en Chine : comment rédiger une petite annonce en chinois ? Quels mots utiliser pour vanter un produit ? Comment annoncer un prix, une réduction, une qualité ?

\textbf{Problématique :} comment lire, comprendre et rédiger une publicité simple en chinois ?

Cette leçon te donne les clés pour y parvenir. À la fin du parcours, tu seras capable de repérer dans une annonce chinoise le produit, le prix, la réduction et le lieu de vente, puis de rédiger ta propre mini-publicité.

\courssub{Brainstorming : la publicité autour de nous}

\begin{experience}[Brainstorming collectif]
En classe entière, réponds à ces questions en français, puis essaie de deviner comment on dirait ces mots en chinois :
\begin{enumerate}
\item Quelles publicités vois-tu chaque jour au Cameroun ? (affiches, spots radio, messages sur le téléphone mobile, panneaux publicitaires, réseaux sociaux)
\item Quels produits sont le plus souvent annoncés ? (crédit téléphonique, boissons, savon, écoles et universités, banques)
\item Quelles informations trouve-t-on toujours dans une publicité ? (le produit, le prix, les avantages, où acheter)
\end{enumerate}
Le professeur note les réponses au tableau et fait deviner les mots chinois correspondants : \zh{广告} (la publicité), \zh{产品} (le produit), \zh{价格} (le prix), \zh{打折} (faire une réduction), \zh{买} (acheter).
\end{experience}

\courssub{Le mot clé : 广告 (guǎnggào)}

\begin{textebox}[Le mot du jour]
广告 guǎnggào — la publicité.\\
广 signifie « vaste, large » ; 告 signifie « annoncer, informer ».\\
La publicité, c'est littéralement « annoncer largement », c'est-à-dire informer le plus grand nombre de personnes possible !
\end{textebox}

Observons la composition de ce mot, car elle aide à le mémoriser :

\begin{center}
\begin{tabularx}{\linewidth}{|X|X|X|X|}
\hline
\rowcolor{popblueL} Caractère & Pinyin & Sens propre & Rôle dans le mot \\
\hline
广 & guǎng & vaste, large & idée de diffusion large \\
\hline
告 & gào & annoncer, informer & action de communiquer \\
\hline
广告 & guǎnggào & publicité, annonce & annoncer largement \\
\hline
\end{tabularx}
\end{center}

\begin{exemplebox}[Premiers exemples]
\zh{这是一个广告。} \emph{Zhè shì yí ge guǎnggào.} « C'est une publicité. »\\
\zh{我在手机上看到了广告。} \emph{Wǒ zài shǒujī shang kàndào le guǎnggào.} « J'ai vu une publicité sur mon téléphone. »\\
\zh{这个广告很有意思。} \emph{Zhège guǎnggào hěn yǒu yìsi.} « Cette publicité est très intéressante. »
\end{exemplebox}

\begin{attention}[Piège : le classificateur]
En français on dit « une publicité » ; en chinois, on ne peut pas dire \zh{一广告}. Il faut un \cle{classificateur} entre le nombre et le nom. Pour 广告, on utilise le classificateur général 个 : \zh{一个广告} \emph{yí ge guǎnggào} « une publicité ». Erreur fréquente des francophones : oublier 个, car le français n'a pas de classificateurs.
\end{attention}

\courssub{Carte mentale du vocabulaire de la publicité}

Avant de lire le texte, organisons les mots du thème autour du mot clé 广告. Cette carte mentale sera complétée tout au long de la leçon.

\begin{popfigure}
\begin{center}\tcbox[colback=popink,colframe=popink,coltext=white,arc=9pt,boxsep=4pt,left=6pt,right=6pt]{\bfseries\small 广告 \emph{guǎnggào} publicité}\end{center}
\vspace{-2pt}
\begin{tcbraster}[raster columns=3,raster equal height=rows,raster column skip=6pt,raster row skip=6pt,enhanced,blankest]
\begin{tcolorbox}[enhanced,colback=poporange!80,colframe=poporange!80,coltext=white,boxrule=0.9pt,arc=3pt,left=4pt,right=4pt,top=3pt,bottom=3pt,boxsep=1pt,halign=left]\bfseries\scriptsize 产品 \emph{chǎnpǐn} produit\end{tcolorbox}
\begin{tcolorbox}[enhanced,colback=popgreen!80,colframe=popgreen!80,coltext=white,boxrule=0.9pt,arc=3pt,left=4pt,right=4pt,top=3pt,bottom=3pt,boxsep=1pt,halign=left]\bfseries\scriptsize 价格 \emph{jiàgé} prix\end{tcolorbox}
\begin{tcolorbox}[enhanced,colback=poppink!80,colframe=poppink!80,coltext=white,boxrule=0.9pt,arc=3pt,left=4pt,right=4pt,top=3pt,bottom=3pt,boxsep=1pt,halign=left]\bfseries\scriptsize 打折 \emph{dǎzhé} réduction\end{tcolorbox}
\begin{tcolorbox}[enhanced,colback=poppurple!80,colframe=poppurple!80,coltext=white,boxrule=0.9pt,arc=3pt,left=4pt,right=4pt,top=3pt,bottom=3pt,boxsep=1pt,halign=left]\bfseries\scriptsize 质量 \emph{zhìliàng} qualité\end{tcolorbox}
\begin{tcolorbox}[enhanced,colback=popgold!90,colframe=popgold!90,coltext=white,boxrule=0.9pt,arc=3pt,left=4pt,right=4pt,top=3pt,bottom=3pt,boxsep=1pt,halign=left]\bfseries\scriptsize 买 \emph{mǎi} acheter\end{tcolorbox}
\end{tcbraster}

\legende{Carte mentale : le vocabulaire de la publicité autour du mot 广告.}
\end{popfigure}

\courssub{Anticipation : deviner le contenu du texte}

\begin{methode}[Lire un texte publicitaire en quatre étapes]
Avant même de connaître tous les mots, on peut comprendre l'essentiel d'une publicité chinoise en suivant une démarche méthodique :
\begin{enumerate}
\item \textbf{Lire le titre} : il contient souvent le nom du produit ou la promesse principale.
\item \textbf{Repérer les chiffres et les prix} : cherche les nombres suivis de 块 (kuai, unité monétaire) ou de 折 (zhé, réduction).
\item \textbf{Chercher le nom du produit} : c'est le nom qui revient le plus souvent dans le texte.
\item \textbf{Lire la phrase d'accroche finale} : elle invite à acheter, avec des mots comme 快来 (kuài lái, « venez vite ») ou 欢迎 (huānyíng, « bienvenue »).
\end{enumerate}
\end{methode}

\begin{experience}[Anticipation guidée]
Le professeur écrit au tableau le titre du texte support de la leçon : \zh{新店开张，果汁大减价！} \emph{Xīn diàn kāizhāng, guǒzhī dà jiǎnjià !} « Ouverture d'un nouveau magasin, grande réduction sur les jus ! »
\begin{enumerate}
\item À ton avis, de quoi parle ce texte ?
\item Quelles informations t'attends-tu à trouver ?
\item Quels mots du brainstorming retrouveras-tu probablement ?
\end{enumerate}
Note tes hypothèses : tu les vérifieras pendant la lecture.
\end{experience}

\courssub{Objectifs de lecture}

\begin{aretenir}[Ce que tu dois trouver dans le texte]
En lisant la publicité de la section suivante, ta mission est de repérer quatre informations précises :
\begin{itemize}
\item \textbf{Le produit} (产品 chǎnpǐn) : que vend-on ?
\item \textbf{Le prix} (价格 jiàgé) : combien coûte le produit ?
\item \textbf{La réduction} (打折 dǎzhé) : y a-t-il une promotion ? De combien ?
\item \textbf{Le lieu} (地方 dìfang) : où peut-on acheter le produit ?
\end{itemize}
\end{aretenir}

\courssub{Rappel : le vocabulaire des prix déjà connu}

Tu as déjà appris à parler des prix. Révisons rapidement ces mots indispensables pour comprendre une publicité.

\begin{center}
\begin{tabularx}{\linewidth}{|X|X|X|X|}
\hline
\rowcolor{popblueL} Caractères & Pinyin & Nature & Traduction \\
\hline
块 & kuài & classificateur / nom & unité monétaire (« yuan » familier) \\
\hline
钱 & qián & nom & argent \\
\hline
便宜 & piányi & adjectif & bon marché, pas cher \\
\hline
贵 & guì & adjectif & cher, coûteux \\
\hline
多少 & duōshao & pronom interrogatif & combien \\
\hline
多少钱 & duōshao qián & expression & combien (ça coûte) ? \\
\hline
\end{tabularx}
\end{center}

\begin{exemplebox}[Révision des structures de prix]
\zh{这个多少钱？} \emph{Zhège duōshao qián ?} « Combien ça coûte ? »\\
\zh{五块钱。} \emph{Wǔ kuài qián.} « Cinq yuans. »\\
\zh{太贵了！} \emph{Tài guì le !} « C'est trop cher ! »\\
\zh{现在很便宜。} \emph{Xiànzài hěn piányi.} « Maintenant c'est très bon marché. »
\end{exemplebox}

\begin{propriete}[La structure du prix]
Structure : \cle{nom du produit + (nombre) + 块 + (钱)}.\\
\zh{果汁十块钱。} \emph{Guǒzhī shí kuài qián.} « Le jus coûte dix yuans. »\\
Remarque : le chinois place le prix directement après le produit, sans verbe « coûter ». En français on dit « le jus \emph{coûte} dix yuans » ; en chinois, la phrase est plus économique : littéralement « jus dix yuans ».
\end{propriete}

\begin{exoresolu}[Repérer les informations clés]
Énoncé : lis cette mini-annonce et trouve le produit, le prix et la réduction.\\
\zh{苹果汁，原价十块，现在八块！} \emph{Píngguǒzhī, yuánjià shí kuài, xiànzài bā kuài !}
\tcblower
Solution détaillée :
\begin{enumerate}
\item \textbf{Le produit} : 苹果汁 \emph{píngguǒzhī} « jus de pomme » (苹果 pomme + 汁 jus).
\item \textbf{Le prix} : 原价 \emph{yuánjià} « prix d'origine » = 十块 « dix yuans » ; 现在 \emph{xiànzài} « maintenant » = 八块 « huit yuans ».
\item \textbf{La réduction} : le prix passe de 10 à 8 yuans, soit 2 yuans de réduction.
\end{enumerate}
Traduction complète : « Jus de pomme, prix d'origine dix yuans, maintenant huit yuans ! »
\end{exoresolu}

\begin{attention}[Piège : 便宜 se prononce piányi]
Beaucoup de francophones prononcent \zh{便宜} « biànyí » en lisant les caractères un par un. Or ici, 便 se lit \emph{pián} (et non biàn) et 宜 se lit \emph{yi} (ton léger). Retiens : \emph{piányi} « bon marché ». Autre piège : ne confonds pas 贵 (guì, cher) avec 鬼 (guǐ, fantôme) — les tons changent tout !
\end{attention}

\begin{savaistu}[La publicité mobile en Chine et au Cameroun]
En Chine, la publicité passe énormément par le téléphone mobile : les consommateurs découvrent les produits sur WeChat (微信 Wēixìn) et sur Douyin (抖音 Dǒuyīn), la version chinoise de TikTok, où des vendeurs présentent leurs produits en direct. Au Cameroun, c'est le même phénomène avec WhatsApp et Facebook : de nombreux commerçants de Douala et de Yaoundé vendent leurs produits (ndolé, tissus, jus de bissap) en publiant des annonces dans des groupes. La publicité mobile est donc un point commun fascinant entre les deux pays !
\end{savaistu}

\begin{exercice}[À toi de jouer]
\begin{enumerate}
\item Complète avec le bon mot : 块, 便宜, 贵, 广告 :\\
a) \zh{这个太＿＿了，我不买。} (C'est trop cher, je n'achète pas.)\\
b) \zh{果汁三＿＿钱。} (Le jus coûte trois yuans.)\\
c) \zh{我在手机上看到一个＿＿。} (J'ai vu une publicité sur mon téléphone.)
\item Traduis en chinois : « Ce jus est très bon marché, seulement cinq yuans ! »
\item D'après le titre \zh{新店开张，果汁大减价！}, quelles seront selon toi les quatre informations clés du texte ?
\end{enumerate}
\end{exercice}

\begin{corrige}[Exercice « À toi de jouer »]
\begin{enumerate}
\item a) 贵 (guì) ; b) 块 (kuài) ; c) 广告 (guǎnggào).
\item \zh{这个果汁很便宜，只要五块钱！} \emph{Zhège guǒzhī hěn piányi, zhǐ yào wǔ kuài qián !}
\item Le produit : du jus (果汁) ; le prix et la réduction : une grande baisse de prix (大减价) ; le lieu : un nouveau magasin (新店).
\end{enumerate}
\end{corrige}

Tu es maintenant prêt à découvrir le texte support : une véritable petite publicité chinoise. Garde en tête tes quatre objectifs de lecture : le produit, le prix, la réduction et le lieu !

\coursec{Leçon 30 — Vocabulaire et compréhension de texte : publicité}

\courssub{Mise en route et découverte du thème}

Dans la séquence précédente, nous avons abordé le monde de l'entreprise et du commerce ; vous savez désormais présenter une entreprise et décrire un poste de travail. Nous entrons maintenant dans le concret de la vie commerciale : la \cle{publicité} (\zh{广告} \emph{guǎnggào}). Que ce soit sur les panneaux lumineux de \zh{北京} (\emph{Běijīng}), les vitrines du quartier \zh{商业街} (\emph{shāngyèjiē}) ou les haut-parleurs du marché \zh{莫科洛} (\emph{Mòkēluó}) à \zh{雅温得} (\emph{Yǎwēndé}), le discours publicitaire suit des codes précis. Cette leçon vous donne les clés pour le décrypter, en comprendre la structure, en maîtriser le vocabulaire et, in fine, en produire un vous-même.

\courssub{Texte support : lecture et écoute}

\begin{textebox}[大减价！— Grande réduction !]
大减价！欢迎光临！\\
Dà jiǎnjià! Huānyíng guānglín!\\
Grande réduction ! Bienvenue !\\[0.4em]
本店出售各种商品：衣服、鞋子、水果和饮料。\\
Běn diàn chūshòu gè zhǒng shāngpǐn: yīfu, xiézi, shuǐguǒ hé yǐnliào.\\
Notre magasin vend toutes sortes de produits : vêtements, chaussures, fruits et boissons.\\[0.4em]
今天所有商品打八折！\\
Jīntiān suǒyǒu shāngpǐn dǎ bā zhé!\\
Aujourd'hui, tous les produits sont à moins 20 \% !\\[0.4em]
我们的东西又好又便宜。\\
Wǒmen de dōngxi yòu hǎo yòu piányi.\\
Nos articles sont à la fois bons et bon marché.\\[0.4em]
快来买吧！\\
Kuài lái mǎi ba!\\
Venez vite acheter !
\end{textebox}

\textbf{Observation.} Avant toute explication, lisez le texte en silence et répondez mentalement à trois questions : de quoi parle-t-on ? Qui parle ? À qui s'adresse-t-on ? Vous remarquez immédiatement les points d'exclamation nombreux (！) : c'est la marque typique du discours publicitaire, qui cherche à capter l'attention et à créer l'enthousiasme. Notez aussi l'usage de \zh{本店} (\emph{běn diàn}, « notre magasin ») plutôt que \zh{我们的店} : une marque de professionnalisme et de politesse commerciale que nous analyserons plus loin.

\courssub{Lecture et écoute : les trois étapes}

\begin{methode}[Comment travailler la lecture d'un texte chinois]
\begin{enumerate}
\item \textbf{Écoute d'abord.} L'enseignant lit le texte une première fois, à vitesse naturelle, livre fermé ou regard baissé : vous vous concentrez uniquement sur la mélodie des tons et le rythme des groupes de mots.
\item \textbf{Lecture chorale.} Toute la classe lit ensemble, en suivant le pinyin, en imitant le modèle de l'enseignant. On répète phrase par phrase, puis le texte entier.
\item \textbf{Lecture individuelle.} Chaque rangée lit une phrase à voix haute ; l'enseignant corrige les tons à chaud. On veille particulièrement :
\begin{itemize}
\item Au changement de ton sur \zh{水果} (\emph{shuǐguǒ} $\rightarrow$ \emph{shuíguǒ}) : deux 3\textsuperscript{e} tons de suite.
\item Au 3\textsuperscript{e} ton de \zh{减} (\emph{jiǎn}) dans \zh{减价} (\emph{jiǎnjià}, 3-4) et de \zh{我} (\emph{wǒ}) dans \zh{我们} (\emph{wǒmen}, 3-neutre).
\item Au ton neutre sur \zh{服} (\emph{fu}) dans \zh{衣服} (\emph{yīfu}) : ne dites pas \emph{yīfú} !
\end{itemize}
\end{enumerate}
\end{methode}

\begin{attention}[Les tons difficiles du texte : rappel de la règle du 3\textsuperscript{e} ton]
Deux successions de troisièmes tons demandent une attention spéciale :
\begin{itemize}
\item \zh{水果} (\emph{shuǐguǒ}, « fruit ») : deux 3\textsuperscript{e} tons de suite. Par la règle du changement de ton (sandhi), le premier devient un 2\textsuperscript{e} ton à l'oral : on prononce donc \emph{shuíguǒ}.
\item \zh{很好} (\emph{hěn hǎo}) serait pareil, mais ici le texte dit \zh{又好} (\emph{yòu hǎo}) : pas de changement, car \zh{又} (\emph{yòu}) porte un 4\textsuperscript{e} ton.
\end{itemize}
Autre piège fréquent : \zh{衣服} (\emph{yīfu}, « vêtement ») se prononce avec un \zh{fu} atone (ton neutre) ; ne dites pas \emph{yīfú} (ce qui signifierait « vêtements » avec une intonation incorrecte, voire un autre sens).
\end{attention}

\courssub{Analyse de la structure publicitaire : les cinq étapes}

Une publicité, en chinois comme en français, obéit à une architecture très régulière. Repérons dans notre texte les cinq étapes du discours commercial.

\begin{popfigure}
\begin{tikzpicture}[node distance=0.4cm, every node/.style={align=center, minimum width=2.8cm, minimum height=1.1cm, rounded corners, font=\small}]
\node[draw=popblue, fill=popblueL, thick, align=center] (a){\textbf{1. Accroche}\\大减价！};
\node[draw=poppurple, fill=poppurpleL, thick, right=of a, align=center] (b){\textbf{2. Produits}\\衣服、鞋子…};
\node[draw=poporange, fill=poporangeL, thick, right=of b, align=center] (c){\textbf{3. Offre}\\打八折！};
\node[draw=popgreen, fill=popgreenL, thick, right=of c, align=center] (d){\textbf{4. Arguments}\\又好又便宜};
\node[draw=poppink, fill=poppinkL, thick, right=of d, align=center] (e){\textbf{5. Action}\\快来买吧！};
\draw[-{Stealth[length=3mm]}, thick, popdark] (a) -- (b);
\draw[-{Stealth[length=3mm]}, thick, popdark] (b) -- (c);
\draw[-{Stealth[length=3mm]}, thick, popdark] (c) -- (d);
\draw[-{Stealth[length=3mm]}, thick, popdark] (d) -- (e);
\end{tikzpicture}
\legende{La structure d'une publicité chinoise : accroche, produits, offre, arguments, appel à l'action.}
\end{popfigure}

\begin{center}
\begin{tabularx}{\linewidth}{|X|X|X|}
\hline
\rowcolor{popblueL} \textbf{Étape} & \textbf{Formule du texte (caractères + pinyin)} & \textbf{Fonction communicative} \\
\hline
1. Accroche & 大减价！ \emph{Dà jiǎnjià!} & Attirer l'attention par l'annonce d'un événement exceptionnel. \\
\hline
2. Produits & 衣服、鞋子、水果和饮料 \emph{yīfu, xiézi, shuǐguǒ hé yǐnliào} & Présenter l'étendue de l'offre (liste non exhaustive signalée par \zh{各种}). \\
\hline
3. Offre / Prix & 打八折 \emph{dǎ bā zhé} & Annoncer la réduction selon la méthode chinoise (part payée). \\
\hline
4. Arguments & 又好又便宜 \emph{yòu hǎo yòu piányi} & Convaincre par la double qualité (qualité + prix). Structure \zh{又…又…}. \\
\hline
5. Appel à l'action & 快来买吧！ \emph{Kuài lái mǎi ba!} & Inciter à l'acte d'achat immédiat (impératif doux + particule \zh{吧}). \\
\hline
\end{tabularx}
\end{center}

\courssub{Points de langue essentiels}

\begin{definition}[打折 — Le système chinois des réductions]
Le verbe \zh{打折} (\emph{dǎ zhé}) signifie « appliquer une réduction / faire un rabais ». Il se construit avec un chiffre entre \zh{打} et \zh{折} : \zh{打} + \textbf{chiffre} + \zh{折}. Ce chiffre indique la \cle{part du prix que le client paie}, exprimée en dixièmes (10 = prix initial).
\begin{itemize}
\item \zh{打九折} (\emph{dǎ jiǔ zhé}) = payer 9/10 = 90 \% du prix $\rightarrow$ \textbf{réduction de 10 \%}.
\item \zh{打八折} (\emph{dǎ bā zhé}) = payer 8/10 = 80 \% du prix $\rightarrow$ \textbf{réduction de 20 \%}.
\item \zh{打五折} (\emph{dǎ wǔ zhé}) = payer 5/10 = 50 \% du prix $\rightarrow$ \textbf{moitié prix (réduction 50 \%)}.
\item \zh{打折扣} (\emph{dǎ zhékoù}) = faire une réduction (terme général).
\end{itemize}
\end{definition}

\begin{exemplebox}[Les réductions à la chinoise : exemples chiffrés]
\zh{原价一百块，打八折，现价八十块。} \emph{Yuánjià yìbǎi kuài, dǎ bā zhé, xiànjià bāshí kuài.} « Prix initial 100 yuans, -20 \%, prix actuel 80 yuans. »\\
\zh{这双鞋子打五折，真便宜！} \emph{Zhè shuāng xiézi dǎ wǔ zhé, zhēn piányi!} « Ces chaussures sont à moitié prix, vraiment pas cher ! »\\
\zh{全场打九折，会员再减十块。} \emph{Quánchǎng dǎ jiǔ zhé, huìyuán zài jiǎn shí kuài.} « Tout le magasin à -10 \%, 10 yuans de plus pour les membres. »
\end{exemplebox}

\begin{attention}[Le piège numéro 1 des francophones : \textbf{ne jamais inverser le calcul !}]
En France, on annonce la part \emph{enlevée} : « -20 \% », « -50 \% ».
En Chine, on annonce la part \emph{payée} : \zh{打八折} (paie 80 \%), \zh{打五折} (paie 50 \%).
\begin{center}
\textbf{\zh{打八折} = « -20 \% » (et non « -80 \% »)} \\
\textbf{\zh{打五折} = « -50 \% » (et non « -5 \% »)}
\end{center}
Une affiche \zh{打八折} correspond à notre « réduction de 20 \% ». Traduire \zh{打八折} par « moins 80 \% » ferait payer 20 yuans au lieu de 80 : erreur fatale pour le commerçant ! Au marché de Douala comme à Shanghai, vérifiez toujours le sens du chiffre affiché.
\end{attention}

\begin{propriete}[La structure 又…又… : deux qualités simultanées]
\textbf{Structure :} \cle{又 + Adj₁ + 又 + Adj₂} (ou Verbe psychologique).
\textbf{Emploi :} Attribue deux qualités simultanées à un même sujet. Équivalent français : « à la fois… et… », « à la fois… et… ».
\textbf{Contrainte sémantique forte :} Les deux adjectifs doivent avoir la \textbf{même polarité} (tous deux positifs \emph{ou} tous deux négatifs). On ne mélange pas.
\end{propriete}

\begin{exemplebox}[又…又… en action (polarité respectée)]
\zh{我们的东西又好又便宜。} \emph{Wǒmen de dōngxi yòu hǎo yòu piányi.} (Positif + Positif) « Nos articles sont à la fois bons et bon marché. »\\
\zh{这个西瓜又大又甜。} \emph{Zhège xīguā yòu dà yòu tián.} (Positif + Positif) « Ce pastèque est à la fois gros et sucré. »\\
\zh{雅温得的雨季又热又湿。} \emph{Yǎwēndé de yǔjì yòu rè yòu shī.} (Négatif/Négatif ou neutre) « La saison des pluies à Yaoundé est à la fois chaude et humide. »\\
\zh{这部电影又长又无聊。} \emph{Zhè bù diànyǐng yòu cháng yòu wúliáo.} (Négatif + Négatif) « Ce film est à la fois long et ennuyeux. »
\end{exemplebox}

\begin{attention}[又…又… : l'erreur de polarité]
On ne dit \textbf{PAS} : \zh{*这件衣服又好又贵} (\emph{*Zhè jiàn yīfu yòu hǎo yòu guì}).
Pour exprimer une concession (qualité + défaut), on utilise \zh{虽然…但是…} (\emph{suīrán… dànshì…}) ou \zh{可是} (\emph{kěshì}) :
\zh{虽然贵，但是很好。} \emph{Suīrán guì, dànshì hěn hǎo.} « Bien que cher, c'est très bon. »
\zh{这件衣服很好，可是太贵了。} \emph{Zhè jiàn yīfu hěn hǎo, kěshì tài guì le.} « Ce vêtement est très bien, mais trop cher. »
\end{attention}

\begin{propriete}[本 + Nom : la politesse institutionnelle et commerciale]
Le préfixe \cle{本} (\emph{běn}, « ce…-ci, le présent, notre ») placé devant un nom d'organisation, de lieu ou de collectivité marque l'appartenance et la responsabilité de l'énonciateur. C'est le registre formel / professionnel.
\begin{center}
\begin{tabularx}{\linewidth}{|X|X|X|}
\hline
\rowcolor{poppurpleL} \textbf{Formule formelle (本 + N)} & \textbf{Pinyin} & \textbf{Équivalent courant (我们的 + N)} \\
\hline
本店 \emph{běn diàn} & notre magasin (le commerçant parle) & 我们的店 \emph{wǒmen de diàn} \\
\hline
本校 \emph{běn xiào} & notre établissement (direction/profs) & 我们的学校 \emph{wǒmen de xuéxiào} \\
\hline
本公司 \emph{běn gōngsī} & notre entreprise (direction/commercial) & 我们的公司 \emph{wǒmen de gōngsī} \\
\hline
本国 \emph{běn guó} & notre pays (officiel, médias) & 我们的国家 \emph{wǒmen de guójiā} \\
\hline
本人 \emph{běn rén} & moi-même (formel, administratif) & 我 \emph{wǒ} \\
\hline
\end{tabularx}
\end{center}
Dans le texte : \zh{本店出售…} « \textbf{Notre magasin} vend… » (le gérant s'adresse au public). Un client dirait : \zh{你们的店卖…} « \textbf{Votre magasin} vend… ».
\end{propriete}

\courssub{Dialogue d'approfondissement : en magasin}

\begin{textebox}[在商店里 — Dans le magasin]
顾客：老板，这件衬衫多少钱？\\
Gùkè: Lǎobǎn, zhè jiàn chènshān duōshao qián?\\
Le client : Patron, combien coûte cette chemise ?\\[0.4em]
老板：一百块。今天打八折，只要八十块！\\
Lǎobǎn: Yìbǎi kuài. Jīntiān dǎ bā zhé, zhǐ yào bāshí kuài!\\
Le patron : Cent yuans. Aujourd'hui, moins 20 \%, seulement quatre-vingts yuans !\\[0.4em]
顾客：太好了！你们的衣服又漂亮又便宜。\\
Gùkè: Tài hǎo le! Nǐmen de yīfu yòu piàoliang yòu piányi.\\
Le client : Formidable ! Vos vêtements sont à la fois beaux et bon marché.\\[0.4em]
老板：谢谢！欢迎下次再来！\\
Lǎobǎn: Xièxie! Huānyíng xià cì zài lái!\\
Le patron : Merci ! Revenez une prochaine fois, soyez le bienvenu !
\end{textebox}

\begin{center}
\begin{tabularx}{\linewidth}{|X|X|X|X|}
\hline
\rowcolor{poporangeL} \textbf{Caractères} & \textbf{Pinyin} & \textbf{Nature} & \textbf{Traduction} \\
\hline
顾客 & gùkè & n. & client \\
\hline
老板 & lǎobǎn & n. & patron, boss \\
\hline
衬衫 & chènshān & n. & chemise \\
\hline
件 & jiàn & class. & classificateur pour vêtements (haut), événements \\
\hline
只要 & zhǐ yào & adv. & seulement, il suffit que ; seulement (prix) \\
\hline
漂亮 & piàoliang & adj. & beau, joli, élégant \\
\hline
下次 & xià cì & n./adv. & la prochaine fois \\
\hline
再来 & zài lái & v. & revenir, venir à nouveau \\
\hline
\end{tabularx}
\end{center}

\courssub{Glossaire du thème : vocabulaire de la publicité et du commerce}

\begin{center}
\begin{tabularx}{\linewidth}{|X|X|X|X|}
\hline
\rowcolor{popblueL} \textbf{Caractères} & \textbf{Pinyin} & \textbf{Nature} & \textbf{Traduction} \\
\hline
广告 & guǎnggào & n. & publicité, annonce \\
\hline
减价 & jiǎnjià & v./n. & baisser le prix / réduction, solde \\
\hline
打折 & dǎ zhé & v. & faire une réduction (structure : 打+chiffre+折) \\
\hline
八折 / 九折 / 五折 & bā zhé / jiǔ zhé / wǔ zhé & n. & 80\% / 90\% / 50\% du prix (voir attention) \\
\hline
半价 & bànjià & n. & demi-prix ( = 打五折) \\
\hline
优惠 & yōuhuì & v./adj./n. & favoriser, avantageux / avantage, promotion \\
\hline
促销 & cùxiāo & v./n. & promotion des ventes, soldes \\
\hline
商品 & shāngpǐn & n. & produit, marchandise \\
\hline
出售 & chūshòu & v. & vendre (formel, écrit) \\
\hline
卖 & mài & v. & vendre (courant) \\
\hline
买 & mǎi & v. & acheter \\
\hline
衣服 & yīfu & n. & vêtements (fu = ton neutre) \\
\hline
鞋子 & xiézi & n. & chaussures \\
\hline
水果 & shuǐguǒ & n. & fruits (shuǐ 3\textsuperscript{e} $\rightarrow$ 2\textsuperscript{e} ton oral) \\
\hline
饮料 & yǐnliào & n. & boissons \\
\hline
便宜 & piányi & adj. & bon marché, pas cher \\
\hline
贵 & guì & adj. & cher \\
\hline
好 & hǎo & adj. & bon, bien \\
\hline
坏 & huài & adj. & mauvais, cassé \\
\hline
漂亮 & piàoliang & adj. & beau, joli \\
\hline
欢迎光临 & huānyíng guānglín & expr. & Bienvenue ! (formule d'accueil client) \\
\hline
顾客 & gùkè & n. & client \\
\hline
老板 & lǎobǎn & n. & patron, commerçant \\
\hline
块 & kuài & class./n. & yuan (unité monétaire courante, oral) \\
\hline
元 & yuán & n. & yuan (unité monétaire formelle, écrit) \\
\hline
\end{tabularx}
\end{center}

\courssub{Clés (radicaux) et ordre des traits : focus sur 6 caractères clés}

\begin{definition}[Analyse graphique : radicaux et ordre des traits]
Maîtriser l'écriture, c'est respecter l'ordre des traits (haut $\rightarrow$ bas, gauche $\rightarrow$ droite, horizontal avant vertical, trait central avant ailes, cadre avant contenu, fermer le cadre en dernier) et identifier la clé (radical) pour le dictionnaire.
\end{definition}

\begin{center}
\begin{tabularx}{\linewidth}{|c|X|X|X|X|}
\hline
\rowcolor{popgreenL} \textbf{Car.} & \textbf{Clé (Radical) + Sens} & \textbf{Nb traits} & \textbf{Ordre des traits (description)} & \textbf{Exemple de mot} \\
\hline
\zh{广} & \zh{广} (yǎn, « toit, abri ») — large, vaste & 3 & 1. Point en haut à gauche 2. Horizontal long 3. Crochet vertical descendant & 广告 \emph{guǎnggào} (publicité) \\
\hline
\zh{折} & \zh{扌} (shǒu, « main ») — action de plier/casser & 7 & 1. Horizontal 2. Crochet vertical 3. Trait montant (main) 4. Horizontal 5. Vertical 6. Horizontal 7. Crochet final & 打折 \emph{dǎ zhé} (réduction) \\
\hline
\zh{售} & \zh{卖} (mài, « vendre ») — phonétique \zh{尸} & 11 & Clé \zh{卖} (7 traits) + \zh{尸} (3 traits) + point. Écrire \zh{卖} puis \zh{尸} + point à droite. & 出售 \emph{chūshòu} (vendre) \\
\hline
\zh{果} & \zh{木} (mù, « arbre, bois ») — fruit sur l'arbre & 8 & Clé \zh{木} (4 traits) + \zh{田} (5 traits : vertical, pli, 3 horizontaux). & 水果 \emph{shuǐguǒ} (fruit) \\
\hline
\zh{饮} & \zh{饣} (yǐn, « manger/boire », forme de \zh{食}) & 7 & Clé \zh{饣} (3 traits : point, horizontal, crochet) + \zh{欠} (4 traits). & 饮料 \emph{yǐnliào} (boisson) \\
\hline
\zh{促} & \zh{亻} (rén, « homme ») — action humaine & 9 & Clé \zh{亻} (2 traits) + \zh{足} (zú, « pied », 7 traits : vertical, horizontal, etc.). & 促销 \emph{cùxiāo} (promotion) \\
\hline
\end{tabularx}
\end{center}

\courssub{Compréhension du texte et collocations}

\begin{exoresolu}[Comprendre l'offre et calculer]
\textbf{Énoncé :} Un pantalon (\zh{裤子} \emph{kùzi}) coûte normalement 150 yuans (\zh{元} \emph{yuán}). Le magasin annonce \zh{打八折}. Combien le client paie-t-il ? Écrivez la phrase chinoise annonçant le prix final et traduisez l'annonce publicitaire en français correct.
\tcblower
\textbf{Solution détaillée :}
\begin{enumerate}
\item \textbf{Analyse de \zh{打八折}} : Le client paie 8/10 du prix initial (80 \%). La réduction est de 20 \%.
\item \textbf{Calcul} : 150 $\times$ 0,8 = 120 yuans.
\item \textbf{Phrase chinoise} : \zh{这条裤子原价一百五十元，打八折，现价一百二十元。} (\emph{Zhè tiáo kùzi yuánjià yìbǎi wǔshí yuán, dǎ bā zhé, xiànjià yìbǎi èrshí yuán.})
\item \textbf{Traduction française correcte} : « Ce pantalon coûte 150 yuans à l'origine, avec une réduction de 20 \%, le prix actuel est de 120 yuans. » Ou version affiche : « \textbf{-20 \% sur tous les pantalons : 120 yuans au lieu de 150.} »
\item \textbf{Piège évité} : Ne pas écrire « -80 \% » (ce qui donnerait 30 yuans).
\end{enumerate}
\end{exoresolu}

\begin{exemplebox}[Collocations incontournables du thème (à mémoriser par paires)]
\begin{tabularx}{\linewidth}{|X|X|X|}
\hline
\rowcolor{poppinkL} \textbf{Chinois (Car. + Pinyin)} & \textbf{Structure} & \textbf{Français} \\
\hline
大减价 \emph{dà jiǎnjià} & Adj + V/N & Grande braderie / Soldes monstres \\
\hline
欢迎光临 \emph{huānyíng guānglín} & V + V & Bienvenue (formule fixe commerce) \\
\hline
出售商品 \emph{chūshòu shāngpǐn} & V + N & Vendre des produits (écrit/formel) \\
\hline
打八折 \emph{dǎ bā zhé} & V + Num + N & Faire -20 \% / Soldes à 80\% \\
\hline
又好又便宜 \emph{yòu hǎo yòu piányi} & 又 Adj 又 Adj & Bon marché et de qualité \\
\hline
快来买 \emph{kuài lái mǎi} & Adv + V + V & Venez vite acheter \\
\hline
只要…块 \emph{zhǐ yào… kuài} & Adv + Prix & Seulement … yuans \\
\hline
下次再来 \emph{xià cì zài lái} & Adv + Adv + V & Revenez la prochaine fois \\
\hline
\end{tabularx}
\end{exemplebox}

\begin{exercice}[Compréhension fine du texte support]
Relisez le texte support et répondez en chinois (caractères + pinyin) aux questions suivantes :
\begin{enumerate}
\item 本店出售什么商品？ (\emph{Běn diàn chūshòu shénme shāngpǐn?})
\item 今天所有商品打几折？ (\emph{Jīntiān suǒyǒu shāngpǐn dǎ jǐ zhé?})
\item 这个店的东西怎么样？ (\emph{Zhège diàn de dōngxi zěnmeyàng?})
\item 广告最后怎么说？ (\emph{Guǎnggào zuìhòu zěnme shuō?})
\item Traduisez en français : \zh{欢迎光临！}
\end{enumerate}
\end{exercice}

\begin{corrige}[Exercice « Compréhension fine »]
1. \zh{本店出售衣服、鞋子、水果和饮料。} \emph{Běn diàn chūshòu yīfu, xiézi, shuǐguǒ hé yǐnliào.} (Le magasin vend des vêtements, chaussures, fruits et boissons.)\\
2. \zh{今天所有商品打八折。} \emph{Jīntiān suǒyǒu shāngpǐn dǎ bā zhé.} (Aujourd'hui tous les produits sont à -20\%).\\
3. \zh{他们的东西又好又便宜。} \emph{Tāmen de dōngxi yòu hǎo yòu piányi.} (Leurs articles sont à la fois bons et bon marché.)\\
4. \zh{广告最后说：「快来买吧！」} \emph{Guǎnggào zuìhòu shuō: « Kuài lái mǎi ba! »} (La pub finit par : « Venez vite acheter ! »).\\
5. « Bienvenue ! » (Formule d'accueil standard dans tout commerce chinois).
\end{corrige}

\courssub{Production guidée : ma publicité}

\begin{experience}[Atelier rédaction : « Spécial Rentrée / Fête de la Jeunesse »]
\textbf{Contexte :} Vous gérez une boutique à \zh{杜阿拉} (\emph{Dùālā}) ou \zh{雅温得} (\emph{Yǎwēndé}). C'est la rentrée scolaire (septembre) ou la \zh{青年节} (\emph{Qīngniánjié}, 11 février). Vous faites une affiche publicitaire.
\textbf{Consignes :}
\begin{enumerate}
\item Choisissez 3 à 4 produits (vocabulaire leçon + connus : \zh{书包} \emph{shūbāo}, \zh{笔记本} \emph{bǐjìběn}, \zh{钢笔} \emph{gāngbǐ}, \zh{T恤} \emph{Txù}, \zh{运动鞋} \emph{yùndòngxié}).
\item Choisissez une réduction : \zh{打九折}, \zh{打八折} ou \zh{打五折}.
\item Rédigez l'affiche en respectant les \textbf{5 étapes} (Accroche, Produits, Offre, Arguments \zh{又…又…}, Action).
\item Écrivez en \textbf{caractères} + \textbf{pinyin} (avec tons). Traduisez en français.
\end{enumerate}
\textbf{Modèle de canevas :}
\begin{itemize}
\item [Accroche] \zh{大减价！} / \zh{开学大促销！} (\emph{Kāixué dà cùxiāo!})
\item [Produits] \zh{本店出售……} (\emph{Běn diàn chūshòu…})
\item [Offre] \zh{今天全部打……折！} (\emph{Jīntiān quánbù dǎ… zhé!})
\item [Arguments] \zh{又……又……} (\emph{Yòu… yòu…})
\item [Action] \zh{快来……吧！} (\emph{Kuài lái… ba!})
\end{itemize}
\textbf{Production orale :} Présentez votre affiche à la classe (1 min). Les camarades notent : produits, réduction, arguments.
\end{experience}

\begin{aretenir}[L'essentiel de la Leçon 30 — Publicité]
\begin{itemize}
\item \textbf{Structure type d'une pub :} Accroche (\zh{大减价}) $\rightarrow$ Produits (\zh{出售…}) $\rightarrow$ Offre (\zh{打八折}) $\rightarrow$ Arguments (\zh{又好又便宜}) $\rightarrow$ Action (\zh{快来买吧}).
\item \textbf{Calcul des réductions :} \zh{打} + Chiffre + \zh{折} = Part \textbf{payée} en dixièmes. \zh{打八折} = 80\% payé = \textbf{-20\%}. \textbf{Jamais} « -80\% ».
\item \textbf{又…又…} : Deux adjectifs de \textbf{même polarité} (++ ou --). « À la fois… et… ». Interdit : *\zh{又好又贵}.
\item \textbf{本店 / 本公司 / 本校} : Formules formelles pour « notre magasin / entreprise / école » (registre pro). Courant : \zh{我们的店}.
\item \textbf{Vocabulaire clé :} \zh{广告, 减价, 打折, 优惠, 促销, 商品, 出售, 欢迎光临, 顾客, 老板, 块/元, 衬衫, 鞋子, 水果, 饮料, 漂亮, 便宜}.
\item \textbf{Écriture :} Radicaux \zh{广} (toit), \zh{扌} (main), \zh{卖} (vendre), \zh{木} (arbre), \zh{饣} (manger), \zh{亻} (homme). Respecter l'ordre des traits.
\end{itemize}
\end{aretenir}

\coursec{Glossaire du thème}

Après avoir découvert le texte publicitaire à l'oral et à l'écrit, nous allons maintenant fixer le lexique essentiel. Cette section constitue votre \textbf{répertoire de référence} pour toute la leçon : chaque mot y est présenté avec son caractère, sa prononciation standard (pinyin aux tons), sa nature grammaticale et son sens précis. Nous les organisons par familles sémantiques pour en faciliter la mémorisation, conformément à la méthode des \emph{champs lexicaux} préconisée au programme.

\courssub{Vocabulaire de base : la publicité et l'acte de vendre}

Commençons par les mots « moteurs » du thème : le nom de l'objet publicitaire, l'objet vendu, et les verbes d'action commerciale. Observez le tableau ci‑dessous, puis lisez les exemples qui suivent à voix haute en imitant les tons.

\begin{popfigure}
\legende{Vocabulaire de base : publicité et vente}
\begin{center}
\begin{tikzpicture}
\node[inner sep=0pt, align=center] (tab){
\begin{tabularx}{\linewidth}{|X|X|X|X|}
\hline
\rowcolor{popblueL} \textbf{Caractères} & \textbf{Pinyin} & \textbf{Nature} & \textbf{Traduction} \\ \hline
广告 & guǎnggào & n. & publicité, annonce \\ \hline
商品 & shāngpǐn & n. & produit, marchandise, article \\ \hline
出售 & chūshòu & v. & vendre, mettre en vente \\ \hline
减价 & jiǎnjià & v./n. & baisser le prix ; solde, réduction de prix \\ \hline
打折 & dǎzhé & v. & faire une réduction, accorder un rabais \\ \hline
折扣 & zhékòu & n. & rabais, remise, taux de réduction \\ \hline
\end{tabularx}
};
\end{tikzpicture}
\end{center}
\end{popfigure}

\begin{exemplebox}[Mise en contexte : phrases types]
Chaque mot s'insère dans une structure simple (SVO) que vous maîtrisez déjà.
\begin{itemize}
    \item \zh{这张广告很漂亮。} \emph{Zhè zhāng guǎnggào hěn piàoliang.} — Cette publicité est très jolie.
    \item \zh{我们的商品质量好。} \emph{Wǒmen de shāngpǐn zhìliàng hǎo.} — Nos produits sont de bonne qualité.
    \item \zh{这家店出售新鲜水果。} \emph{Zhè jiā diàn chūshòu xīnxiān shuǐguǒ.} — Ce magasin vend des fruits frais.
    \item \zh{冬天衣服减价。} \emph{Dōngtiān yīfu jiǎnjià.} — En hiver, les vêtements sont en solde.
    \item \zh{全场打八折。} \emph{Quánchǎng dǎ bā zhé.} — Tout le magasin est à 20\% de réduction (litt. « fait 8/10 du prix »).
    \item \zh{这个折扣很大。} \emph{Zhège zhékòu hěn dà.} — Cette remise est importante.
\end{itemize}
\end{exemplebox}

\begin{methode}[Mémoriser par familles de caractères (clé \textbf{商})]
Le chinois construit son lexique par \textbf{composition morphologique} : un même caractère « racine » entre dans la fabrication de dizaines de mots. Repérer la racine \textbf{商} (commerce, marchand) démultiplie votre vocabulaire sans effort supplémentaire.

\begin{center}
\begin{tabularx}{\linewidth}{|X|X|X|}
\hline
\rowcolor{poporangeL} \textbf{Mot} & \textbf{Pinyin} & \textbf{Sens} \\ \hline
商品 & shāngpǐn & produit, marchandise \\ \hline
商店 & shāngdiàn & magasin, boutique \\ \hline
商人 & shāngrén & commerçant, marchand \\ \hline
商业 & shāngyè & commerce, activité commerciale \\ \hline
商量 & shāngliang & discuter, négocier (au sens de « mettre en commun ») \\ \hline
\end{tabularx}
\end{center}

\emph{Procédé :} Apprenez \textbf{商} une fois (clé \emph{bouche} 口 + \emph{huit} 八 + \emph{un} 一, ordre des traits : haut → bas, gauche → droite) ; chaque nouveau mot ne vous coûtera plus que l'apprentissage du second caractère.
\end{methode}

\courssub{Vocabulaire complémentaire : qualité, prix et relation client}

Une publicité efficace vante la qualité, justifie le prix et accueille le client. Voici le second bloc lexical, indispensable pour comprendre et rédiger une annonce.

\begin{popfigure}
\legende{Vocabulaire : qualité, prix, clientèle}
\begin{center}
\begin{tikzpicture}
\node[inner sep=0pt, align=center] (tab){
\begin{tabularx}{\linewidth}{|X|X|X|X|}
\hline
\rowcolor{popgreenL} \textbf{Caractères} & \textbf{Pinyin} & \textbf{Nature} & \textbf{Traduction} \\ \hline
质量 & zhìliàng & n. & qualité \\ \hline
价格 & jiàgé & n. & prix \\ \hline
便宜 & piányi & adj. & bon marché, peu cher \\ \hline
贵 & guì & adj. & cher, coûteux \\ \hline
顾客 & gùkè & n. & client, cliente \\ \hline
欢迎 & huānyíng & v. & accueillir, souhaiter la bienvenue \\ \hline
光临 & guānglín & v. & honorer de sa visite (formule de politesse) \\ \hline
\end{tabularx}
};
\end{tikzpicture}
\end{center}
\end{popfigure}

\begin{attention}[Piège de prononciation : \textbf{便宜} se prononce \emph{piányi} (et non \emph{biànyí})]
C'est l'erreur n°1 des francophones à l'oral (HSK / Bac). Le caractère \zh{便} a deux lectures principales :
\begin{itemize}
    \item \textbf{pián} (2e ton) dans \zh{便宜} (bon marché), \zh{方便} (pratique, commode).
    \item \textbf{biàn} (4e ton) dans \zh{便利} (commodité), \zh{大便} (aller à la selle), \zh{便当} (bentō, boîte-repas).
\end{itemize}
\emph{Astuce mnémotechnique :} « Pour un bon \textbf{prix} (pián), c'est \textbf{pratique} (pián) ! » Les deux sens « bon marché » et « commode » partagent le ton montant \textbf{pián}.
\end{attention}

\begin{exemplebox}[Collocations prix / qualité / accueil]
\begin{itemize}
    \item \zh{质量好，价格不贵。} \emph{Zhìliàng hǎo, jiàgé bù guì.} — Bonne qualité, prix pas cher.
    \item \zh{这个手机很便宜。} \emph{Zhège shǒujī hěn piányi.} — Ce téléphone est très bon marché.
    \item \zh{顾客是上帝。} \emph{Gùkè shì Shàngdì.} — Le client est roi (litt. « Dieu »).
    \item \zh{欢迎光临！} \emph{Huānyíng guānglín!} — Bienvenue ! (formule d'entrée de magasin).
    \item \zh{欢迎下次再来。} \emph{Huānyíng xiàcì zài lái.} — Bienvenue, revenez la prochaine fois ! (formule de sortie).
\end{itemize}
\end{exemplebox}

\begin{textebox}[Formule rituelle : 欢迎光临！]
\textbf{欢迎光临！} \emph{Huānyíng guānglín!}

Cette phrase résonne à l'entrée de \textbf{tous} les magasins en Chine, des grands centres commerciaux de Pékin aux petites épiceries du Yunnan. Elle est l'équivalent exact du « Bienvenue ! » ou « Bonjour, bienvenue ! » que lancent les commerçants camerounais aux marchés de Mokolo (Yaoundé) ou de Nkololoun (Douala) quand un client approche.

\begin{itemize}
    \item \zh{欢迎} (huānyíng) = « accueillir avec joie » (bienvenue).
    \item \zh{光临} (guānglín) = « honorer de sa présence » (verbe humble, le client « illumine » le lieu de sa venue).
\end{itemize}
Ne répondez pas par « 谢谢 » (merci) mais par un sourire, un hochement de tête, ou simplement \zh{你好} (nǐ hǎo).
\end{textebox}

\courssub{Outils grammaticaux du thème : quantificateurs, coordination, adverbe de manière}

Ces trois entrées ne sont pas de simples mots de vocabulaire : ce sont des \textbf{outils syntaxiques} que vous retrouverez dans presque toute phrase publicitaire.

\begin{popfigure}
\legende{Outils syntaxiques : quantificateur, coordination, adverbe}
\begin{center}
\begin{tikzpicture}
\node[inner sep=0pt, align=center] (tab){
\begin{tabularx}{\linewidth}{|X|X|X|X|}
\hline
\rowcolor{poppurpleL} \textbf{Caractères} & \textbf{Pinyin} & \textbf{Nature} & \textbf{Traduction} \\ \hline
各种 & gè zhǒng & loc. (dét. + classif.) & toutes sortes de, divers \\ \hline
饮料 & yǐnliào & n. & boisson \\ \hline
又...又... & yòu...yòu... & conj. (corrélatif) & à la fois... et... \\ \hline
快 & kuài & adv./adj. & vite, rapide(ment) \\ \hline
\end{tabularx}
};
\end{tikzpicture}
\end{center}
\end{popfigure}

\begin{propriete}[Structure \textbf{又...又...} (yòu...yòu...) : coordination d'adjectifs]
\emph{Emploi :} Lie deux adjectifs \textbf{qualitatifs} (souvent à connotation positive ou négative) pour décrire un même sujet. En français : « à la fois... et... », « à la fois... et... ».
\emph{Structure :} \textbf{Sujet + 又 + Adj₁ + 又 + Adj₂}
\emph{Contrainte :} Les deux adjectifs doivent être de même registre (tous deux positifs ou tous deux négatifs) et généralement monosyllabiques ou bisyllabiques.

\begin{center}
\begin{tabularx}{\linewidth}{|X|X|X|}
\hline
\rowcolor{poppinkL} \textbf{Chinois (caractères + pinyin)} & \textbf{Traduction} & \textbf{Analyse} \\ \hline
这款饮料 \textbf{又甜又凉}。 & Cette boisson est à la fois sucrée et fraîche. & Adj° positifs \\ \hline
这个广告 \textbf{又长又无聊}。 & Cette pub est à la fois longue et ennuyeuse. & Adj° négatifs \\ \hline
我们的服务 \textbf{又快又好}。 & Notre service est à la fois rapide et bon. & Adj° positifs \\ \hline
\end{tabularx}
\end{center}

\emph{Attention :} On n'utilise \textbf{PAS} \zh{又...又...} pour coordonner un adjectif positif et un négatif (ex. *又贵又好). Pour cela, on utilisera \zh{虽然...但是...} (bien que... mais...) vu en grammaire.
\end{propriete}

\begin{exemplebox}[Exemples intégrés au thème publicitaire]
\begin{itemize}
    \item \zh{我们有各种饮料。} \emph{Wǒmen yǒu gè zhǒng yǐnliào.} — Nous avons toutes sortes de boissons.
    \item \zh{这瓶果汁又甜又好喝。} \emph{Zhè píng guǒzhī yòu tián yòu hǎohē.} — Ce jus est à la fois sucré et délicieux.
    \item \zh{送货很快。} \emph{Sònghuò hěn kuài.} — La livraison est très rapide.
    \item \zh{请快点儿来！} \emph{Qǐng kuài diǎnr lái!} — Veuillez venir vite ! (口语 \zh{快点儿}).
\end{itemize}
\end{exemplebox}

\courssub{Entraînement de prononciation : tons difficiles du glossaire}

Trois mots du glossaire concentrent les difficultés tonales pour un francophone. Travaillez-les en isolé, puis en phrase, avec un miroir ou un enregistreur.

\begin{definition}[Points de vigilance tonale]
\begin{itemize}
    \item \textbf{广告 \emph{guǎnggào} (3-4)} : Le premier ton est un \emph{ton bas descendant-levant} (214). Ne le prononcez pas plat (1er ton) ni seulement descendant (4e ton). La voix doit « creuser » puis remonter légèrement avant de chuter sur \emph{gào}.
    \item \textbf{便宜 \emph{piányi} (2-0)} : \emph{pián} est un 2e ton montant pur (35). \emph{yi} est un \emph{ton neutre} (léger, court, sans marque de ton), attaché au mot précédent. \textbf{Jamais} \emph{biànyí} (4-2) !
    \item \textbf{质量 \emph{zhìliàng} (4-4)} : Deux 4e tons (chutants 51) consécutifs. En chinois standard, le premier 4e ton ne se modifie pas (pas de sandhi comme pour \zh{不} ou \zh{一}), mais l'enchaînement demande de l'énergie : \emph{zhì-liàng}, net, sec, sans « traîner ».
\end{itemize}
\end{definition}

\begin{exoresolu}[Exercice de discrimination tonale]
\textbf{Écoutez (ou lisez à voix haute) et cochez la bonne prononciation.}

1. \zh{广告} : \textcircled{a} guānggào \quad \textcircled{b} \textbf{guǎnggào} \quad \textcircled{c} guǎnggǎo \\
2. \zh{便宜} : \textcircled{a} biànyí \quad \textcircled{b} piányí \quad \textcircled{c} \textbf{piányi} \\
3. \zh{质量} : \textcircled{a} zhíliáng \quad \textcircled{b} zhìliǎng \quad \textcircled{c} \textbf{zhìliàng}

\tcblower
\textbf{Solution détaillée :}
1. \textcircled{b} \textbf{guǎnggào} (3-4). \emph{guǎng} = ton creusé ; \emph{gào} = ton chutant.
2. \textcircled{c} \textbf{piányi} (2-neutre). \emph{pián} = ton montant ; \emph{yi} = ton neutre (court, léger).
3. \textcircled{c} \textbf{zhìliàng} (4-4). Deux tons chutants pleins. \emph{zhì} (qualité) et \emph{liàng} (quantité/mesure) gardent leur ton d'origine.
\end{exoresolu}

\courssub{Clés (radicaux) et ordre des traits : 7 caractères stratégiques}

Connaître la clé (radical) et l'ordre des traits permet de \textbf{retrouver le caractère dans un dictionnaire}, de l'écrire correctement (vitesse, lisibilité) et de deviner le sens des mots inconnus. Voici l'analyse pour 7 caractères « à haut rendement » du glossaire. La colonne « Nombre de traits » indique le nombre de traits de l'écriture courante (simplifiée) et, entre parenthèses, le nombre de traits Kangxi (utilisé pour l'indexation dans les dictionnaires traditionnels).

\begin{center}
\begin{tabularx}{\linewidth}{|c|X|X|X|}
\hline
\rowcolor{popgoldL} \textbf{Caractère} & \textbf{Clé (Radical) + Sens de la clé} & \textbf{Nombre de traits (Écriture / Kangxi)} & \textbf{Ordre des traits (règles principales)} \\ \hline
\zh{广} \emph{guǎng} & \zh{广} \textbf{guǎng} « abri, toit » (clé n°53) — \emph{Le caractère EST sa propre clé.} & 3 / 3 & 1. Point en haut à gauche 丶 2. Trait horizontal 長 3. Crochet vertical descendant 乛 \\ \hline
\zh{商} \emph{shāng} & \zh{口} \textbf{kǒu} « bouche » (clé n°30) — En bas. & 8 / 11 & Règle « enveloppe avant contenu » : écrire \zh{亠} (haut, 2 traits), puis \zh{八} (2 traits) + \zh{一} (1 trait) (milieu), enfin \zh{口} (bas, 3 traits, se ferme en dernier). \\ \hline
\zh{质} \emph{zhì} & \zh{贝} \textbf{bèi} « coquillage / monnaie » (clé n°154) — À gauche. & 8 / 15 & Gauche → droite. \zh{贝} d'abord (4 traits), puis \zh{斤} \emph{jīn} (hache, 4 traits) à droite. \\ \hline
\zh{价} \emph{jià} & \zh{人} \textbf{rén} « homme » (clé n°9) — À gauche (forme \zh{亻}). & 13 / 15 & \zh{亻} (2 traits) puis partie droite \zh{贾} : haut \zh{覀} (3 traits), milieu \zh{贝} (4 traits), bas \zh{贝} (4 traits). \\ \hline
\zh{顾} \emph{gù} & \zh{页} \textbf{yè} « page / tête » (clé n°181) — En bas. & 10 / 21 & Haut → bas. Partie supérieure (forme simplifiée de \zh{雇} / trad. \zh{僱}, 6 traits ici) puis \zh{页} en bas (6 traits). *Note : forme simplifiée très différente du traditionnel \zh{顧} (21 traits).* \\ \hline
\zh{饮} \emph{yǐn} & \zh{饣} \textbf{shí} « manger / nourriture » (clé n°184, forme de \zh{食}) — À gauche. & 7 / 12 & \zh{饣} (3 traits : point, horizontal, crochet bas) puis \zh{欠} \emph{qiàn} (4 traits : horizontal, vertical, horizontal, crochet). \\ \hline
\zh{折} \emph{zhé} & \zh{扌} \textbf{shǒu} « main / action manuelle » (clé n°64) — À gauche. & 7 / 7 & \zh{扌} (3 traits) puis \zh{斤} \emph{jīn} (4 traits : horizontal, vertical, horizontal, crochet). \\ \hline
\end{tabularx}
\end{center}

\courssub{Mini-dictée : 8 mots du glossaire}

Préparez une feuille. Votre professeur (ou l'audio de l'application OnBuch+) dictera les 8 mots suivants \textbf{en chinois uniquement}. Vous écrirez : \textbf{Caractères + Pinyin (tons) + Traduction française}. C'est un entraînement type « compréhension de l'oral + écriture de caractères » (épreuve du Bac).

\begin{exercice}[Mini-dictée — Liste pour l'enseignant / auto-évaluation]
\textbf{Mots à dicter (ordre aléatoire recommandé) :}
\begin{enumerate}
    \item 广告 (guǎnggào)
    \item 商品 (shāngpǐn)
    \item 减价 (jiǎnjià)
    \item 打折 (dǎzhé)
    \item 质量 (zhìliàng)
    \item 便宜 (piányi)
    \item 顾客 (gùkè)
    \item 欢迎 (huānyíng)
\end{enumerate}
\end{exercice}

\begin{corrige}[Mini-dictée n°30]
\begin{center}
\begin{tabularx}{\linewidth}{|c|X|X|X|}
\hline
\rowcolor{popblueL} \textbf{N°} & \textbf{Caractères} & \textbf{Pinyin} & \textbf{Traduction} \\ \hline
1 & 广告 & guǎnggào & publicité \\ \hline
2 & 商品 & shāngpǐn & produit, marchandise \\ \hline
3 & 减价 & jiǎnjià & baisser le prix / solde \\ \hline
4 & 打折 & dǎzhé & faire une réduction \\ \hline
5 & 质量 & zhìliàng & qualité \\ \hline
6 & 便宜 & piányi & bon marché \\ \hline
7 & 顾客 & gùkè & client \\ \hline
8 & 欢迎 & huānyíng & accueillir / bienvenue \\ \hline
\end{tabularx}
\end{center}

\emph{Auto-correction :} Vérifiez scrupuleusement les \textbf{tons} (surtout \emph{guǎnggào} 3-4, \emph{piányi} 2-0, \emph{zhìliàng} 4-4) et l'\textbf{ordre des traits} des caractères \zh{商}, \zh{质}, \zh{价}, \zh{顾}, \zh{饮}, \zh{折}. Chaque erreur de trait = 0.5pt en moins en examen officiel.
\end{corrige}

\textbf{Transition :} Ce glossaire maîtrisé, nous allons maintenant (Section 4) approfondir l'analyse des \textbf{clés et de l'ordre des traits} pour une sélection de caractères, puis (Section 5) tester votre compréhension fine du texte publicitaire et des collocations lexicales.

\coursec{Clés (radicaux) et ordre des traits}

Dans la section précédente, nous avons appris le glossaire de la publicité : \zh{广告} (guǎnggào, « publicité »), \zh{商品} (shāngpǐn, « marchandise »), \zh{便宜} (piányi, « bon marché »), \zh{贵} (guì, « cher »). Ces mots sont maintenant familiers à l'oreille et à l'œil. Mais un bon sinisant ne se contente pas de \emph{reconnaître} les caractères : il sait les \emph{écrire}, trait par trait, dans le bon ordre. C'est l'objet de cette section : comprendre la logique des \cle{clés} (radicaux) et maîtriser l'\cle{ordre des traits} de cinq caractères essentiels du thème : \zh{告}, \zh{商}, \zh{品}, \zh{贵}, \zh{便}.

\courssub{Qu'est-ce qu'une clé (radical) ?}

\begin{definition}[La clé ou radical]
Une \cle{clé} (部首, bùshǒu) est un élément graphique qui se répète dans de nombreux caractères et qui indique souvent le \textbf{domaine de sens} du caractère. Les dictionnaires chinois classent les caractères par clés. Connaître les clés, c'est posséder une carte mentale de l'écriture chinoise : on devine le sens d'un caractère inconnu et on mémorise plus vite.
\end{definition}

Observez ces caractères de notre glossaire :

\begin{exemplebox}[Observation : la clé 口 (kǒu, bouche)]
\zh{告} (gào, annoncer) contient \zh{口}.\\
\zh{品} (pǐn, article, produit) contient trois \zh{口}.\\
\zh{商} (shāng, commerce) contient aussi \zh{口}.\\
\zh{吃} (chī, manger), \zh{喝} (hē, boire), \zh{叫} (jiào, appeler) contiennent \zh{口}.
\end{exemplebox}

La clé \zh{口} (la bouche) apparaît dans les caractères liés à la parole, à la bouche, au goût. Logique : pour \emph{annoncer} une publicité, on utilise la bouche ; pour juger un \emph{produit}, on le goûte avec la bouche.

\begin{definition}[La clé 贝 (bèi, coquillage)]
La clé \cle{贝} (bèi) représente un \textbf{coquillage}. Dans la Chine ancienne, les coquillages servaient de \textbf{monnaie}. Aujourd'hui encore, la clé \zh{贝} signale les caractères liés à \textbf{l'argent, au prix et au commerce} : \zh{贵} (guì, cher), \zh{购} (gòu, acheter), \zh{赚} (zhuàn, gagner de l'argent), \zh{货} (huò, marchandise), \zh{费} (fèi, frais).
\end{definition}

\begin{savaistu}[Pourquoi trois bouches dans 品 ?]
\zh{品} (pǐn) signifie « article, produit, qualité ». L'étymologie traditionnelle l'explique ainsi : \textbf{trois bouches qui goûtent} un aliment permettent de \emph{juger sa qualité}. D'où les mots \zh{商品} (shāngpǐn, marchandise), \zh{食品} (shípǐn, produit alimentaire), \zh{品质} (pǐnzhì, qualité) et même \zh{品尝} (pǐncháng, déguster). Retenir l'image des trois bouches, c'est retenir le caractère pour toujours.
\end{savaistu}

\courssub{Les règles générales d'ordre des traits}

Avant d'étudier chaque caractère, il faut connaître les lois qui gouvernent l'écriture de \emph{tous} les caractères chinois.

\begin{methode}[Les quatre règles d'or de l'ordre des traits]
\begin{enumerate}
\item \textbf{De haut en bas} : on écrit d'abord la partie supérieure, puis la partie inférieure. Exemple : \zh{告} = le haut (\zh{牛} sans la pointe qui dépasse en bas) d'abord, la bouche \zh{口} ensuite.
\item \textbf{De gauche à droite} : on écrit d'abord la partie gauche, puis la partie droite. Exemple : \zh{便} = la clé personne \zh{亻} à gauche d'abord, puis \zh{更} à droite.
\item \textbf{L'extérieur avant l'intérieur} : pour les caractères encadrés, on trace d'abord le cadre, puis le contenu. Exemple : dans \zh{商}, on écrit le haut, puis la partie extérieure, puis l'intérieur.
\item \textbf{On ferme la boîte en dernier} : quand un caractère contient une enceinte fermée comme \zh{口}, le trait du bas qui « ferme » la boîte se trace en dernier, après le contenu.
\end{enumerate}
\end{methode}

\begin{attention}[Erreur fréquente des francophones]
En français, on écrit les lettres dans l'ordre de lecture, sans règle spatiale. En chinois, \textbf{l'ordre des traits n'est pas facultatif} : il conditionne la lisibilité, la vitesse et la beauté de l'écriture. Piège classique : écrire \zh{口} en un seul geste circulaire comme un « o » français. Non ! \zh{口} s'écrit en \textbf{trois traits} : trait vertical gauche (丨), puis le coin haut-droite ((trait brisé)), puis le trait horizontal du bas (一) qui ferme. Jamais de rond tracé d'un seul coup.
\end{attention}

\courssub{Étude détaillée des cinq caractères du thème}

\textbf{1. 告 (gào) — annoncer, informer}\par
Clé : \zh{口} (bouche). Nombre de traits : \textbf{7}. Le caractère se décompose en deux parties : en haut, une forme proche de \zh{牛} (niú, bœuf, sans la pointe qui dépasse en bas) ; en bas, la bouche \zh{口}. Image ancienne : un bœuf et une bouche — on annonçait un sacrifice. Ordre des traits : 丿 一 丨 一 (la partie du haut), puis 丨 (trait brisé) 一 (la bouche, fermée en dernier).

\textbf{2. 商 (shāng) — commerce, affaires}\par
Clé : \zh{亠} (tête) / \zh{口}. Nombre de traits : \textbf{11}. C'est un caractère à étages : on écrit \textbf{d'abord le haut} (le point et la ligne horizontale \zh{亠}, puis les deux petits traits), \textbf{puis l'extérieur} du cadre inférieur, \textbf{puis l'intérieur} (le petit \zh{八} et la bouche \zh{口}). On le retrouve dans \zh{商店} (shāngdiàn, magasin) et \zh{商人} (shāngrén, commerçant).

\textbf{3. 品 (pǐn) — article, produit, qualité}\par
Clé : \zh{口}. Nombre de traits : \textbf{9} (3 bouches × 3 traits). Trois « bouches » : on écrit \textbf{d'abord celle du haut}, \textbf{puis celle du bas à gauche}, \textbf{puis celle du bas à droite}. Chaque \zh{口} suit la règle : on ferme la boîte en dernier.

\textbf{4. 贵 (guì) — cher, précieux}\par
Clé : \zh{贝} (coquillage = argent). Nombre de traits : \textbf{9}. De haut en bas : d'abord la partie supérieure (\zh{中} surmonté d'une ligne), puis la clé \zh{贝} en bas. Rappel précieux : la présence de \zh{贝} vous signale immédiatement que le mot parle d'argent — \zh{贵} (cher), c'est ce qui coûte beaucoup de coquillages !

\textbf{5. 便 (pián) — pratique, bon marché}\par
Clé : \zh{亻} (homme, variante debout de \zh{人}). Nombre de traits : \textbf{9}. Règle de gauche à droite : on trace \textbf{d'abord la clé personne \zh{亻} à gauche} (deux traits : 丿 puis 丨), puis la partie droite \zh{更}. Attention à la prononciation : \zh{便} se lit \emph{pián} dans \zh{便宜} (piányi, bon marché) mais \emph{biàn} dans \zh{方便} (fāngbiàn, pratique).

\begin{center}
\begin{tabularx}{\linewidth}{|X|X|X|X|}
\hline
\rowcolor{popblueL} Caractère & Clé & Traits & Ordre à retenir \\
\hline
\zh{告} gào & \zh{口} bouche & 7 & haut (\zh{牛}) puis bouche en bas \\
\hline
\zh{商} shāng & \zh{亠}/\zh{口} & 11 & haut → extérieur → intérieur \\
\hline
\zh{品} pǐn & \zh{口} & 9 & bouche du haut, puis bas-gauche, puis bas-droite \\
\hline
\zh{贵} guì & \zh{贝} argent & 9 & de haut en bas, \zh{贝} en dernier \\
\hline
\zh{便} pián & \zh{亻} homme & 9 & \zh{亻} à gauche d'abord, puis \zh{更} \\
\hline
\end{tabularx}
\end{center}

\courssub{Schémas de décomposition}

\begin{popfigure}
\begin{tikzpicture}[scale=1, every node/.style={font=\large}]
% Décomposition de 品
\node[draw=popblue, thick, minimum width=1.1cm, minimum height=1.1cm, fill=popblueL] (p1) at (0,1.6) {\zh{口}};
\node[draw=poporange, thick, minimum width=1.1cm, minimum height=1.1cm, fill=poporangeL] (p2) at (-0.7,0) {\zh{口}};
\node[draw=popgreen, thick, minimum width=1.1cm, minimum height=1.1cm, fill=popgreenL] (p3) at (0.7,0) {\zh{口}};
\node[circle, draw=popdark, fill=popdark, text=white, inner sep=1.5pt, font=\small] at (-0.45,2.05) {1};
\node[circle, draw=popdark, fill=popdark, text=white, inner sep=1.5pt, font=\small] at (-1.15,0.45) {2};
\node[circle, draw=popdark, fill=popdark, text=white, inner sep=1.5pt, font=\small] at (1.15,0.45) {3};
\draw[-{Stealth}, thick, popmuted] (2.0,0.8) -- (3.0,0.8);
\node[draw=poppurple, very thick, minimum width=1.3cm, minimum height=1.3cm, fill=poppurpleL, font=\Huge] at (4.2,0.8) {\zh{品}};
\node[align=center, font=\small, text=popmuted] at (4.2,-0.4) {pǐn\\ produit};
\end{tikzpicture}
\legende{Décomposition de 品 : trois bouches 口, tracées dans l'ordre 1 (haut), 2 (bas-gauche), 3 (bas-droite).}
\end{popfigure}

\begin{popfigure}
\begin{tikzpicture}[scale=1, every node/.style={font=\large}]
% Décomposition de 告
\node[draw=popblue, thick, minimum width=1.6cm, minimum height=1.1cm, fill=popblueL, align=center] (g1) at (0,1.4) {\zh{牛}\\[-2pt]{\small (haut)}};
\node[draw=poporange, thick, minimum width=1.1cm, minimum height=1.1cm, fill=poporangeL, align=center] (g2) at (0,0) {\zh{口}};
\node[circle, draw=popdark, fill=popdark, text=white, inner sep=1.5pt, font=\small] at (-0.95,1.85) {1};
\node[circle, draw=popdark, fill=popdark, text=white, inner sep=1.5pt, font=\small] at (-0.7,0.45) {2};
\draw[-{Stealth}, thick, popmuted] (1.6,0.7) -- (2.6,0.7);
\node[draw=poppurple, very thick, minimum width=1.3cm, minimum height=1.3cm, fill=poppurpleL, font=\Huge] at (3.8,0.7) {\zh{告}};
\node[align=center, font=\small, text=popmuted] at (3.8,-0.5) {gào\\ annoncer};
\node[align=center, font=\small, text=popdark] at (0,-1.1) {de haut en bas :\\ la bouche en dernier};
\end{tikzpicture}
\legende{Décomposition de 告 : la partie haute (proche de 牛) d'abord, la bouche 口 ensuite.}
\end{popfigure}

\courssub{Pièges graphiques : caractères proches}

\begin{attention}[Ne pas confondre 品 et 晶]
\zh{品} (pǐn, produit) est composé de \textbf{trois bouches} \zh{口}. \zh{晶} (jīng, cristal, brillant) est composé de \textbf{trois soleils} \zh{日}. La différence tient à un petit trait horizontal au milieu de chaque élément : \zh{口} est une boîte vide, \zh{日} contient une ligne. Dans une publicité, écrire \zh{晶牌} au lieu de \zh{品牌} (pǐnpái, marque) serait une faute grave ! Vérifiez toujours : boîte vide = bouche = produit ; boîte avec une ligne = soleil = cristal.
\end{attention}

\begin{exoresolu}[Identifier la clé et compter les traits]
Énoncé : pour chaque caractère, donne la clé, le nombre de traits et le domaine de sens suggéré par la clé : \zh{贵}, \zh{便}, \zh{品}.
\tcblower
Solution détaillée :
\begin{itemize}
\item \zh{贵} (guì) : clé \zh{贝} (coquillage), 9 traits. Domaine : l'argent, les prix → « cher ».
\item \zh{便} (pián) : clé \zh{亻} (homme), 9 traits. Domaine : la personne, les qualités humaines → « pratique, bon marché ».
\item \zh{品} (pǐn) : clé \zh{口} (bouche), 9 traits (3 × 3). Domaine : la bouche, le goût → « article, qualité » (trois bouches qui goûtent).
\end{itemize}
\end{exoresolu}

\begin{exercice}[Tracé sur ardoise]
Sur ton ardoise ou ton cahier, écris les caractères \zh{告}, \zh{品} et \zh{贵}, \textbf{cinq fois chacun}, en respectant l'ordre des traits appris dans cette leçon. Pour chaque caractère, prononce le pinyin à voix haute à chaque répétition (gào, pǐn, guì) et dis la traduction française à la fin de chaque ligne. Vérifie ensuite : as-tu fermé chaque \zh{口} avec son trait du bas en dernier ? As-tu écrit \zh{贝} en bas de \zh{贵} ?
\end{exercice}

\begin{corrige}[Exercice de tracé]
Auto-vérification guidée :
\begin{itemize}
\item \zh{告} (7 traits) : 丿 一 丨 一 puis 丨 (trait brisé) 一. La bouche est bien en bas et fermée en dernier.
\item \zh{品} (9 traits) : bouche du haut, puis bas-gauche, puis bas-droite ; chaque bouche en 3 traits, jamais en un rond.
\item \zh{贵} (9 traits) : partie haute d'abord, clé \zh{贝} en bas ; le dernier trait est le point final de \zh{贝}.
\end{itemize}
Si un caractère te semble « déséquilibré », reprends-le en comptant les traits à voix haute : la régularité vient de la répétition.
\end{corrige}

\begin{aretenir}[L'essentiel de la section]
\begin{itemize}
\item Une \cle{clé} (radical) indique le domaine de sens : \zh{口} → bouche, parole, goût ; \zh{贝} → argent, commerce ; \zh{亻} → personne.
\item Quatre règles d'ordre des traits : de haut en bas, de gauche à droite, l'extérieur avant l'intérieur, on ferme la boîte en dernier.
\item \zh{告} (7 traits), \zh{商} (11 traits), \zh{品} (9 traits), \zh{贵} (9 traits), \zh{便} (9 traits).
\item Ne jamais confondre \zh{品} (trois bouches, produit) et \zh{晶} (trois soleils, cristal).
\end{itemize}
\end{aretenir}

\coursec{Compréhension du texte et collocations}

Dans la section précédente, nous avons appris à écrire les caractères clés de la publicité : \zh{广} (clé 广, « abri »), \zh{告} (clé 口, « bouche »), \zh{商} ou encore \zh{折}. Savoir écrire, c'est bien ; mais un bon sinisant doit aussi \cle{comprendre} ce qu'il lit et \cle{réutiliser} le vocabulaire dans ses propres phrases. C'est l'objectif de cette section : répondre à des questions de compréhension en phrases complètes, puis maîtriser les collocations (mots qui vont naturellement ensemble) du monde de la publicité.

\courssub{Rappel du texte support}

Relisons ensemble la publicité étudiée dans cette leçon. Lisez-la une première fois en silence, puis à voix haute en suivant le pinyin.

\begin{textebox}[Publicité : 大减价！Grande réduction !]
大减价！大减价！\\
Dà jiǎnjià! Dà jiǎnjià!\\
« Grande réduction ! Grande réduction ! »\\[0.4em]
欢迎光临友谊商店！我们商店出售各种商品：衣服、鞋子、水果、饮料，还有手机和电脑。\\
Huānyíng guānglín Yǒuyì Shāngdiàn! Wǒmen shāngdiàn chūshòu gè zhǒng shāngpǐn: yīfu, xiézi, shuǐguǒ, yǐnliào, háiyǒu shǒujī hé diànnǎo.\\
« Bienvenue au magasin de l'Amitié ! Notre magasin vend toutes sortes de produits : vêtements, chaussures, fruits, boissons, et aussi des téléphones portables et des ordinateurs. »\\[0.4em]
我们商店的商品质量很好，价格也很合理。\\
Wǒmen shāngdiàn de shāngpǐn zhìliàng hěn hǎo, jiàgé yě hěn hélǐ.\\
« Les produits de notre magasin sont de très bonne qualité, et les prix sont aussi très raisonnables. »\\[0.4em]
今天所有商品打八折！一件衬衫原来一百块，现在只要八十块！\\
Jīntiān suǒyǒu shāngpǐn dǎ bā zhé! Yí jiàn chènshān yuánlái yìbǎi kuài, xiànzài zhǐ yào bāshí kuài!\\
« Aujourd'hui, tous les produits sont à moins vingt pour cent ! Une chemise coûtait cent yuans, maintenant elle ne coûte que quatre-vingts yuans ! »\\[0.4em]
机会难得，快来买吧！我们等你！\\
Jīhuì nándé, kuài lái mǎi ba! Wǒmen děng nǐ!\\
« L'occasion est rare, venez vite acheter ! Nous vous attendons ! »
\end{textebox}

\textbf{Glossaire rapide du texte}

\begin{center}
\begin{tabularx}{\linewidth}{|X|X|X|X|}
\hline
\rowcolor{popblueL} Caractères & Pinyin & Nature & Traduction \\
\hline
减价 & jiǎnjià & verbe / nom & réduire les prix ; réduction \\
\hline
出售 & chūshòu & verbe & vendre, mettre en vente \\
\hline
商品 & shāngpǐn & nom & produit, marchandise \\
\hline
质量 & zhìliàng & nom & qualité \\
\hline
合理 & hélǐ & adjectif & raisonnable \\
\hline
打折 & dǎ zhé & verbe & faire une réduction \\
\hline
原来 & yuánlái & adverbe & à l'origine, avant \\
\hline
只要 & zhǐ yào & conjonction / adverbe & il suffit de ; seulement \\
\hline
机会难得 & jīhuì nándé & expression & l'occasion est rare \\
\hline
\end{tabularx}
\end{center}

\begin{attention}[只要 : deux emplois à ne pas confondre]
Dans le texte, \zh{只要八十块} signifie « seulement quatre-vingts yuans » : \zh{只要} a ici la valeur de l'adverbe \zh{只} (« seulement »), renforcé par \zh{要}. Mais \zh{只要……就……} est aussi une conjonction de condition : \zh{只要你来，我就很高兴。} (Zhǐyào nǐ lái, wǒ jiù hěn gāoxìng.) — « Pourvu que tu viennes, je serai très content. » Dans notre publicité, pas de \zh{就} : c'est donc le sens « seulement ».
\end{attention}

\courssub{Méthode : répondre à une question de compréhension}

En français comme en chinois, une bonne réponse de compréhension obéit à trois étapes : \cle{repérer} l'information dans le texte, \cle{citer} l'élément qui justifie la réponse, puis \cle{rédiger} une phrase complète qui reprend les mots de la question.

\begin{methode}[Répondre à une question de compréhension en chinois]
\begin{enumerate}
\item \textbf{Lire la question} et souligner le mot interrogatif : \zh{什么} (quoi), \zh{几} (combien), \zh{怎么样} (comment), \zh{谁} (qui), \zh{哪家} (quel, pour les magasins, écoles…).
\item \textbf{Repérer dans le texte} la phrase qui contient la réponse : cherche les mots de la question (ex. \zh{打几折} → cherche \zh{打折} dans le texte).
\item \textbf{Citer l'élément du texte} qui justifie la réponse : ex. \zh{今天所有商品打八折。}
\item \textbf{Rédiger une phrase complète} en reprenant le sujet et le verbe de la question : ne réponds pas « \zh{八折} », écris « \zh{今天打八折。} » (Jīntiān dǎ bā zhé. — Aujourd'hui, la réduction est de vingt pour cent).
\end{enumerate}
\end{methode}

\begin{attention}[Le piège de la réponse trop courte]
En français, on répond souvent par un mot : « Quelle réduction ? — Vingt pour cent. » En chinois, à l'écrit, l'examinateur attend une \cle{phrase complète} avec sujet et verbe. À la question \zh{今天打几折？}, la réponse « \zh{八折。} » seule est \textbf{incomplète}. Écris : \zh{今天所有商品打八折。} (Jīntiān suǒyǒu shāngpǐn dǎ bā zhé.) — « Aujourd'hui, tous les produits sont à moins vingt pour cent. » Reprendre les mots de la question dans la réponse est la marque d'un bon élève.
\end{attention}

\begin{popfigure}
\begin{tikzpicture}[
  node distance=0.55cm and 0.9cm,
  q/.style={rectangle, rounded corners, draw=popblue, fill=popblueL, align=center, minimum width=2.6cm, minimum height=1cm, font=\small},
  s/.style={rectangle, rounded corners, draw=poporange, fill=poporangeL, align=center, minimum width=2.6cm, minimum height=1cm, font=\small},
  c/.style={rectangle, rounded corners, draw=poppurple, fill=poppurpleL, align=center, minimum width=2.6cm, minimum height=1cm, font=\small},
  r/.style={rectangle, rounded corners, draw=popgreen, fill=popgreenL, align=center, minimum width=2.6cm, minimum height=1cm, font=\small},
  arr/.style={-{Stealth[length=2.5mm]}, thick, draw=popdark}
]
\node[q, align=center] (quest){1. Question\\今天打几折？};
\node[s, right=of quest, align=center] (spot){2. Repérage\\chercher 打折\\dans le texte};
\node[c, right=of spot, align=center] (cite){3. Citation\\今天所有商品\\打八折。};
\node[r, right=of cite, align=center] (rep){4. Réponse complète\\今天打八折。};
\draw[arr] (quest) -- (spot);
\draw[arr] (spot) -- (cite);
\draw[arr] (cite) -- (rep);
\end{tikzpicture}
\legende{Organigramme de la méthode : de la question à la réponse en phrase complète.}
\end{popfigure}

\courssub{Questions de compréhension : correction collective}

Appliquons la méthode aux cinq questions de la leçon. Pour chacune, nous donnons la citation du texte qui justifie la réponse, puis la réponse complète rédigée.

\begin{exoresolu}[Question 1 : 这是哪家商店的广告？]
Énoncé : \zh{这是哪家商店的广告？} (Zhè shì nǎ jiā shāngdiàn de guǎnggào?) — De quel magasin s'agit-il dans cette publicité ?
\tcblower
\textbf{Repérage} : la question contient \zh{哪家商店} (quel magasin) ; on cherche le nom propre du magasin dans le texte. Rappel : \zh{家} (jiā) est le classificateur des magasins, entreprises et familles.\\
\textbf{Citation} : \zh{欢迎光临友谊商店！} (Huānyíng guānglín Yǒuyì Shāngdiàn!)\\
\textbf{Réponse complète} : \zh{这是友谊商店的广告。} (Zhè shì Yǒuyì Shāngdiàn de guǎnggào.) — « C'est la publicité du magasin de l'Amitié. »
\end{exoresolu}

\begin{exoresolu}[Question 2 : 商店出售什么商品？]
Énoncé : \zh{商店出售什么商品？} (Shāngdiàn chūshòu shénme shāngpǐn?) — Quels produits le magasin vend-il ?
\tcblower
\textbf{Repérage} : le verbe \zh{出售} (vendre) est repris tel quel dans le texte.\\
\textbf{Citation} : \zh{我们商店出售各种商品：衣服、鞋子、水果、饮料，还有手机和电脑。}\\
\textbf{Réponse complète} : \zh{商店出售衣服、鞋子、水果、饮料，还有手机和电脑。} (Shāngdiàn chūshòu yīfu, xiézi, shuǐguǒ, yǐnliào, háiyǒu shǒujī hé diànnǎo.) — « Le magasin vend des vêtements, des chaussures, des fruits, des boissons, et aussi des téléphones portables et des ordinateurs. » Remarque : \zh{还有} (háiyǒu, « et aussi ») sert à ajouter un dernier élément à une liste ; dans une énumération, la virgule chinoise \zh{、} (dunhao) sépare les mots de même nature.
\end{exoresolu}

\begin{exoresolu}[Question 3 : 今天打几折？]
Énoncé : \zh{今天打几折？} (Jīntiān dǎ jǐ zhé?) — Quelle est la réduction aujourd'hui ?
\tcblower
\textbf{Repérage} : \zh{几} (combien) + \zh{折} annonce un chiffre.\\
\textbf{Citation} : \zh{今天所有商品打八折！}\\
\textbf{Réponse complète} : \zh{今天所有商品打八折。} (Jīntiān suǒyǒu shāngpǐn dǎ bā zhé.) — « Aujourd'hui, tous les produits sont à moins vingt pour cent. »\\
\textbf{Attention au calcul} : en chinois, \zh{打八折} signifie « payer 80 \% du prix », donc une réduction de \textbf{20 \%}. \zh{打七折} = payer 70 \% (réduction de 30 \%). C'est l'inverse de la logique française : le chiffre chinois indique ce qu'on \emph{paie}, pas ce qu'on \emph{économise}. Le texte le confirme : la chemise à 100 yuans coûte maintenant 80 yuans (\zh{原来一百块，现在只要八十块}).
\end{exoresolu}

\begin{exoresolu}[Question 4 : 商品怎么样？]
Énoncé : \zh{商品怎么样？} (Shāngpǐn zěnmeyàng?) — Comment sont les produits ?
\tcblower
\textbf{Repérage} : \zh{怎么样} demande un avis, une appréciation ; on cherche les adjectifs du texte.\\
\textbf{Citation} : \zh{我们商店的商品质量很好，价格也很合理。}\\
\textbf{Réponse complète} : \zh{商品质量很好，价格也很合理。} (Shāngpǐn zhìliàng hěn hǎo, jiàgé yě hěn hélǐ.) — « Les produits sont de très bonne qualité, et les prix sont aussi très raisonnables. »
\end{exoresolu}

\begin{exoresolu}[Question 5 : 广告最后请顾客做什么？]
Énoncé : \zh{广告最后请顾客做什么？} (Guǎnggào zuìhòu qǐng gùkè zuò shénme?) — Que demande la publicité aux clients à la fin ?
\tcblower
\textbf{Repérage} : \zh{最后} (à la fin) oriente vers la dernière phrase du texte.\\
\textbf{Citation} : \zh{机会难得，快来买吧！我们等你！}\\
\textbf{Réponse complète} : \zh{广告最后请顾客快来买。} (Guǎnggào zuìhòu qǐng gùkè kuài lái mǎi.) — « À la fin, la publicité demande aux clients de venir vite acheter. » Le petit mot \zh{吧} (ba) adoucit l'impératif : c'est une invitation chaleureuse, pas un ordre.
\end{exoresolu}

\courssub{Grammaire express : 怎么样, demander un avis}

\begin{propriete}[怎么样 zěnmeyàng : « comment est-ce ? »]
Le mot interrogatif \zh{怎么样} sert à \cle{demander un avis} sur une personne, une chose, un produit. Structure : \textbf{Sujet + 怎么样？}\\
\zh{这件衬衫怎么样？} (Zhè jiàn chènshān zěnmeyàng?) — « Comment est cette chemise ? »\\
\zh{喀麦隆的天气怎么样？} (Kāmàilóng de tiānqì zěnmeyàng?) — « Comment est le temps au Cameroun ? »\\
On y répond avec un adjectif précédé de \zh{很} ou \zh{非常} : \zh{质量很好。} (La qualité est très bonne.)\\
Variante pratique : \zh{……怎么样？} après une proposition signifie « et si on… ? » : \zh{我们去商店，怎么样？} (Wǒmen qù shāngdiàn, zěnmeyàng?) — « On va au magasin, qu'en dis-tu ? »
\end{propriete}

\begin{attention}[怎么样 ou 怎么 ?]
Ne confonds pas \zh{怎么} (zěnme, « comment », devant un verbe d'action) et \zh{怎么样} (zěnmeyàng, avis sur quelque chose).\\
\zh{怎么去商店？} (Zěnme qù shāngdiàn?) — « Comment aller au magasin ? » (le moyen, le chemin)\\
\zh{这个商店怎么样？} (Zhège shāngdiàn zěnmeyàng?) — « Comment est ce magasin ? » (l'avis, l'appréciation)
\end{attention}

\courssub{Les collocations essentielles de la publicité}

Une \cle{collocation} est un couple de mots qui vont naturellement ensemble : en français on dit « faire de la publicité » et non « faire publicité ». Le chinois fonctionne pareil : il faut apprendre les mots \emph{par paires}. Voici les sept collocations du thème.

\begin{center}
\begin{tabularx}{\linewidth}{|X|X|X|}
\hline
\rowcolor{popblueL} Collocation & Pinyin & Traduction \\
\hline
做广告 & zuò guǎnggào & faire de la publicité \\
\hline
大减价 & dà jiǎnjià & grande réduction, soldes \\
\hline
打八折 & dǎ bā zhé & réduction de 20 \% (payer 80 \%) \\
\hline
质量好 & zhìliàng hǎo & de bonne qualité \\
\hline
价格合理 & jiàgé hélǐ & prix raisonnable \\
\hline
欢迎光临 & huānyíng guānglín & bienvenue (formule de magasin) \\
\hline
快来买 & kuài lái mǎi & venez vite acheter \\
\hline
\end{tabularx}
\end{center}

\begin{exemplebox}[Chaque collocation dans une phrase]
\zh{这家店做广告。} \emph{Zhè jiā diàn zuò guǎnggào.} « Ce magasin fait de la publicité. »\\
\zh{周末超市大减价。} \emph{Zhōumò chāoshì dà jiǎnjià.} « Ce week-end, le supermarché fait de grandes réductions. »\\
\zh{今天所有商品打八折。} \emph{Jīntiān suǒyǒu shāngpǐn dǎ bā zhé.} « Aujourd'hui, tous les produits sont à moins vingt pour cent. »\\
\zh{这个商店的商品质量很好。} \emph{Zhège shāngdiàn de shāngpǐn zhìliàng hěn hǎo.} « Les produits de ce magasin sont de très bonne qualité. »\\
\zh{这个商店的价格很合理。} \emph{Zhège shāngdiàn de jiàgé hěn hélǐ.} « Les prix de ce magasin sont très raisonnables. »\\
\zh{欢迎光临！您要买什么？} \emph{Huānyíng guānglín! Nín yào mǎi shénme?} « Bienvenue ! Que souhaitez-vous acheter ? »\\
\zh{快来买吧，机会难得！} \emph{Kuài lái mǎi ba, jīhuì nándé!} « Venez vite acheter, l'occasion est rare ! »
\end{exemplebox}

\begin{attention}[Verbe + complément : ne sépare pas les collocations]
Les francophones traduisent souvent mot à mot et cassent les collocations. Retiens :\\
— On dit \zh{做广告} (zuò guǎnggào), jamais \zh{做宣传} pour « faire de la pub » au sens commercial courant (\zh{宣传} xuānchuán = propagande, communication officielle).\\
— \zh{打折} est inséparable dans l'idée : le chiffre se place \emph{entre} les deux : \zh{打八折}, \zh{打九折} (−10 \%), \zh{打五折} (−50 \%).\\
— \zh{欢迎光临} est une formule figée de politesse commerciale ; ne la traduis pas mot à mot (\zh{光临} = « honorer de sa présence », très formel et toujours positif).
\end{attention}

\begin{aretenir}[L'essentiel de la section]
\begin{itemize}
\item Une réponse de compréhension = phrase complète qui reprend les mots de la question + citation du texte comme justification.
\item \zh{打八折} = payer 80 \% du prix, donc \textbf{−20 \%} : le chiffre chinois dit ce qu'on paie.
\item \zh{怎么样} sert à demander un avis : \zh{商品怎么样？} — « Comment sont les produits ? »
\item Les collocations s'apprennent par paires : \zh{做广告}, \zh{大减价}, \zh{质量好}, \zh{价格合理}, \zh{欢迎光临}, \zh{快来买}.
\end{itemize}
\end{aretenir}

\courssub{Reformulation : du texte publicitaire au dialogue}

Transformer un texte en dialogue est un exercice classique de l'approche par les compétences : il oblige à \cle{sélectionner} les informations, à \cle{changer de point de vue} (le vendeur dit \zh{我们}, le client dit \zh{您}) et à \cle{réutiliser les collocations} à l'oral. Deux mots nouveaux apparaissent : \zh{售货员} (shòuhuòyuán, « vendeur, vendeuse ») et \zh{顾客} (gùkè, « client »). Voici le modèle en quatre répliques.

\begin{textebox}[Dialogue vendeur–client (modèle de reformulation)]
售货员：欢迎光临友谊商店！今天所有商品打八折！\\
Shòuhuòyuán: Huānyíng guānglín Yǒuyì Shāngdiàn! Jīntiān suǒyǒu shāngpǐn dǎ bā zhé!\\
Vendeur : « Bienvenue au magasin de l'Amitié ! Aujourd'hui, tous les produits sont à moins vingt pour cent ! »\\[0.4em]
顾客：太好了！这件衬衫怎么样？\\
Gùkè: Tài hǎo le! Zhè jiàn chènshān zěnmeyàng?\\
Cliente : « C'est formidable ! Et cette chemise, comment est-elle ? »\\[0.4em]
售货员：质量很好，价格也很合理：原来一百块，现在只要八十块。\\
Shòuhuòyuán: Zhìliàng hěn hǎo, jiàgé yě hěn hélǐ: yuánlái yìbǎi kuài, xiànzài zhǐ yào bāshí kuài.\\
Vendeur : « La qualité est très bonne et le prix est très raisonnable : elle coûtait cent yuans, maintenant elle ne coûte que quatre-vingts yuans. »\\[0.4em]
顾客：好，我买了！\\
Gùkè: Hǎo, wǒ mǎi le!\\
Cliente : « D'accord, je la prends ! »
\end{textebox}

Observe comment chaque réplique reprend une information du texte : le nom du magasin et la réduction (réplique 1), la question d'avis avec \zh{怎么样} (réplique 2), la qualité et le prix avec l'exemple de la chemise (réplique 3), et la décision d'achat (réplique 4). Rien n'est inventé : tout vient du texte, seule la \emph{forme} change. Note aussi l'emploi de \zh{了} dans \zh{我买了} : il marque ici la \cle{décision prise}, l'acte d'achat accompli — « je la prends ! ».

\begin{exercice}[À toi de jouer]
\begin{enumerate}
\item Réponds par une phrase complète en chinois : \zh{友谊商店出售手机和电脑吗？} (Yǒuyì Shāngdiàn chūshòu shǒujī hé diànnǎo ma?) — Le magasin de l'Amitié vend-il des téléphones et des ordinateurs ?
\item Traduis en chinois : « Ce magasin fait de la publicité : aujourd'hui, grande réduction, tout est à moins trente pour cent (\zh{打七折}) ! »
\item Avec un camarade, écris ton propre dialogue vendeur–client en 4 répliques pour un magasin de Yaoundé ou de Douala, en utilisant au moins quatre collocations de la leçon (\zh{欢迎光临}, \zh{打八折}, \zh{质量很好}, \zh{价格合理}, \zh{快来买}).
\end{enumerate}
\end{exercice}

\begin{corrige}[Exercice : Compréhension et collocations]
\textbf{1.} \zh{友谊商店出售手机和电脑。} (Yǒuyì Shāngdiàn chūshòu shǒujī hé diànnǎo.) — « Oui, le magasin de l'Amitié vend des téléphones portables et des ordinateurs. » Justification : le texte dit \zh{还有手机和电脑}. On peut aussi faire précéder la réponse de \zh{是的，} (« oui »), mais la phrase complète reste obligatoire.\\
\textbf{2.} \zh{这个商店做广告：今天大减价，所有商品打七折！} (Zhège shāngdiàn zuò guǎnggào: jīntiān dà jiǎnjià, suǒyǒu shāngpǐn dǎ qī zhé!) — Vérifie que tu as bien mis le chiffre \emph{entre} \zh{打} et \zh{折}, et que \zh{打七折} correspond bien à « moins trente pour cent » (on paie 70 \%).\\
\textbf{3.} Corrigé collectif en classe : le professeur vérifie la présence des quatre collocations, l'usage correct de \zh{怎么样}, le classificateur \zh{件} devant \zh{衬衫} et les tons du pinyin.
\end{corrige}

\coursec{Production guidée : ma publicité}

Après avoir lu, compris et analysé le texte support et ses collocations (\zh{打八折}, \zh{又便宜又好}, \zh{欢迎光临}…), il est temps de passer à la \cle{production} : chaque élève va rédiger sa propre mini-publicité en chinois. C'est l'aboutissement de la leçon : le vocabulaire et les structures ne sont vraiment acquis que lorsqu'on les réemploie soi-même. Cette tâche est aussi un excellent entraînement à l'épreuve d'expression écrite du Baccalauréat, où l'on demande souvent un court texte publicitaire ou informatif.

\courssub{La tâche : consigne et contraintes}

\begin{definition}[La tâche à réaliser]
Rédiger une \cle{mini-publicité} de \textbf{4 à 6 phrases} en chinois simplifié pour promouvoir un \textbf{produit camerounais} au choix : le jus de bissap \zh{洛神花汁} (luòshénhuā zhī), l'huile de palme \zh{棕榈油} (zōnglǘ yóu), les pagnes \zh{花布} (huā bù) ou le café de l'Ouest \zh{咖啡} (kāfēi). La publicité doit suivre la trame imposée en six étapes et n'utiliser que la banque de mots autorisés.
\end{definition}

\begin{propriete}[La trame imposée en six étapes]
Toute mini-publicité de cette leçon suit obligatoirement l'ordre suivant :
\begin{enumerate}
\item \textbf{Accroche} : \zh{大减价！} (dà jiǎnjià ! « Grande réduction ! ») ou \zh{好消息！} (hǎo xiāoxi ! « Bonne nouvelle ! ») ;
\item \textbf{Présentation du magasin} : \zh{本店出售……} (běn diàn chūshòu… « Notre magasin vend… ») ;
\item \textbf{Produits} : nommer un ou deux produits avec un classificateur ou une précision (origine, adjectif) ;
\item \textbf{Offre} : \zh{打……折} (dǎ … zhé « réduction de… ») ;
\item \textbf{Argument} : \zh{又……又……} (yòu … yòu … « à la fois… et… ») ;
\item \textbf{Appel à l'action} : \zh{快来买吧！} (kuài lái mǎi ba ! « Venez vite acheter ! »).
\end{enumerate}
\end{propriete}

\begin{popfigure}
\begin{tikzpicture}[node distance=0.35cm and 0.9cm]
\node[draw=popblue, fill=popblueL, rounded corners, minimum width=3.1cm, minimum height=0.8cm, align=center, font=\small] (e1) at (0,0) {\textbf{1. Accroche}\\ 大减价！/ 好消息！};
\node[draw=poporange, fill=poporangeL, rounded corners, minimum width=3.1cm, minimum height=0.8cm, align=center, font=\small, below right=of e1] (e2) {\textbf{2. Magasin}\\ 本店出售……};
\node[draw=popgreen, fill=popgreenL, rounded corners, minimum width=3.1cm, minimum height=0.8cm, align=center, font=\small, below right=of e2] (e3) {\textbf{3. Produits}\\ 咖啡、花布……};
\node[draw=poppurple, fill=poppurpleL, rounded corners, minimum width=3.1cm, minimum height=0.8cm, align=center, font=\small, below right=of e3] (e4) {\textbf{4. Offre}\\ 打九折！};
\node[draw=poppink, fill=poppinkL, rounded corners, minimum width=3.1cm, minimum height=0.8cm, align=center, font=\small, below right=of e4] (e5) {\textbf{5. Argument}\\ 又……又……};
\node[draw=popgold, fill=popgoldL, rounded corners, minimum width=3.1cm, minimum height=0.8cm, align=center, font=\small, below right=of e5] (e6) {\textbf{6. Appel}\\ 快来买吧！};
\draw[-{Stealth}, thick, popdark] (e1.east) -- (e2.north);
\draw[-{Stealth}, thick, popdark] (e2.east) -- (e3.north);
\draw[-{Stealth}, thick, popdark] (e3.east) -- (e4.north);
\draw[-{Stealth}, thick, popdark] (e4.east) -- (e5.north);
\draw[-{Stealth}, thick, popdark] (e5.east) -- (e6.north);
\end{tikzpicture}
\legende{Les six étapes de la mini-publicité, en escalier : on descend marche par marche, sans en sauter aucune.}
\end{popfigure}

\courssub{Observer avant d'écrire : un modèle complet}

Avant de rédiger, observons une publicité modèle rédigée selon la trame. Repérez chaque étape.

\begin{textebox}[Modèle de publicité : le café de l'Ouest]
好消息！欢迎光临！\\
Hǎo xiāoxi! Huānyíng guānglín!\\
« Bonne nouvelle ! Bienvenue ! »\\
本店出售西部省的好咖啡。\\
Běn diàn chūshòu Xībù shěng de hǎo kāfēi.\\
« Notre magasin vend du bon café de la province de l'Ouest. »\\
今天打九折，一公斤只要三千法郎！\\
Jīntiān dǎ jiǔ zhé, yì gōngjīn zhǐyào sān qiān fǎláng!\\
« Aujourd'hui moins dix pour cent, le kilo à seulement trois mille francs ! »\\
我们的咖啡又香又便宜。\\
Wǒmen de kāfēi yòu xiāng yòu piányi.\\
« Notre café est à la fois parfumé et bon marché. »\\
快来买吧！\\
Kuài lái mǎi ba!\\
« Venez vite l'acheter ! »
\end{textebox}

\begin{exemplebox}[Exemples traduits à imiter]
\begin{itemize}
\item \zh{好消息！本店的咖啡又好又便宜，今天打九折！}\\
\emph{Hǎo xiāoxi! Běn diàn de kāfēi yòu hǎo yòu piányi, jīntiān dǎ jiǔ zhé!}\\
« Bonne nouvelle ! Notre café est bon et bon marché, moins dix pour cent aujourd'hui ! »
\item \zh{大减价！本店的花布又漂亮又便宜，快来看看吧！}\\
\emph{Dà jiǎnjià! Běn diàn de huābù yòu piàoliang yòu piányi, kuài lái kànkan ba!}\\
« Grande réduction ! Nos pagnes sont beaux et bon marché, venez vite voir ! »
\item \zh{本店出售新鲜的棕榈油，今天打八折！}\\
\emph{Běn diàn chūshòu xīnxiān de zōnglǘyóu, jīntiān dǎ bā zhé!}\\
« Notre magasin vend de l'huile de palme fraîche, moins vingt pour cent aujourd'hui ! »
\end{itemize}
\end{exemplebox}

\courssub{La banque de mots autorisés}

Au Baccalauréat comme ici, on n'écrit bien qu'avec des mots que l'on maîtrise. Voici la banque de mots autorisés : le glossaire de la leçon, complété par le vocabulaire des prix vu en Première et les noms des produits ciblés.

\begin{center}
\begin{tabularx}{\linewidth}{|X|X|X|X|}
\hline
\rowcolor{popblueL} Caractères & Pinyin & Nature & Traduction \\
\hline
好消息 & hǎo xiāoxi & expression & bonne nouvelle \\
\hline
大减价 & dà jiǎnjià & expression & grande réduction \\
\hline
欢迎光临 & huānyíng guānglín & expression & bienvenue (aux clients) \\
\hline
本店 & běn diàn & nom & notre magasin \\
\hline
出售 & chūshòu & verbe & vendre, mettre en vente \\
\hline
打折 & dǎ zhé & verbe + compl. & faire une réduction \\
\hline
又……又…… & yòu… yòu… & structure & à la fois… et… \\
\hline
便宜 & piányi & adjectif & bon marché \\
\hline
新鲜 & xīnxiān & adjectif & frais \\
\hline
漂亮 & piàoliang & adjectif & beau, joli \\
\hline
香 & xiāng & adjectif & parfumé, savoureux \\
\hline
只要 & zhǐyào & conjonction & seulement (prix), pourvu que \\
\hline
块 & kuài & classif. & yuan (unité monétaire orale) \\
\hline
法郎 & fǎláng & nom & franc CFA \\
\hline
公斤 & gōngjīn & nom & kilogramme \\
\hline
快来买吧 & kuài lái mǎi ba & expression & venez vite acheter \\
\hline
买一送一 & mǎi yī sòng yī & expression & un acheté, un offert \\
\hline
洛神花汁 & luòshénhuā zhī & nom & jus de bissap (hibiscus) \\
\hline
棕榈油 & zōnglǘ yóu & nom & huile de palme \\
\hline
花布 & huā bù & nom & pagne, tissu wax \\
\hline
咖啡 & kāfēi & nom & café \\
\hline
西部省 & Xībù shěng & nom propre & province de l'Ouest (Cameroun) \\
\hline
喀麦隆 & Kāmàilóng & nom propre & Cameroun \\
\hline
杜阿拉 & Dùālā & nom propre & Douala \\
\hline
\end{tabularx}
\end{center}

\begin{attention}[Piège : le sens de 打九折]
Les francophones traduisent souvent \zh{打九折} par « moins 90 \% » : c'est l'inverse ! En chinois, \zh{打九折} signifie « on paie 90 \% du prix », donc une réduction de \textbf{10 \%}. De même, \zh{打八折} = on paie 80 \% = réduction de 20 \%. Plus le chiffre est petit, plus la réduction est grande : \zh{打五折} = moitié prix. Retenez : chiffre chinois élevé = petite réduction.
\end{attention}

\courssub{Le modèle à trous : première ébauche}

Pour les élèves qui hésitent, on commence par compléter ce modèle à trous, puis on le personnalise.

\begin{textebox}[Modèle à compléter]
大\_\_\_\_！欢迎光临！\\
Dà \_\_\_\_! Huānyíng guānglín!\\
本店出售\_\_\_\_。\\
Běn diàn chūshòu \_\_\_\_.\\
今天打\_\_\_\_折！\\
Jīntiān dǎ \_\_\_\_ zhé!\\
我们的东西又\_\_\_\_又\_\_\_\_。\\
Wǒmen de dōngxi yòu \_\_\_\_ yòu \_\_\_\_.\\
快来买吧！\\
Kuài lái mǎi ba!
\end{textebox}

\begin{exoresolu}[Compléter le modèle pour le jus de bissap]
Énoncé : complétez le modèle à trous pour annoncer une promotion sur le jus de bissap (\zh{洛神花汁}), avec une réduction de 20 \%, en utilisant les adjectifs \zh{新鲜} et \zh{便宜}.
\tcblower
Solution détaillée :\\
1. Accroche : \zh{大减价！} (dà jiǎnjià) convient parfaitement pour une promotion.\\
2. Produit : \zh{本店出售新鲜的洛神花汁。} (Běn diàn chūshòu xīnxiān de luòshénhuā zhī.) « Notre magasin vend du jus de bissap frais. »\\
3. Offre : réduction de 20 \% = on paie 80 \% = \zh{打八折} : \zh{今天打八折！} (Jīntiān dǎ bā zhé !).\\
4. Argument : \zh{我们的东西又新鲜又便宜。} (Wǒmen de dōngxi yòu xīnxiān yòu piányi.) « Nos produits sont frais et bon marché. »\\
5. Appel : \zh{快来买吧！}\\
Texte complet : \zh{大减价！欢迎光临！本店出售新鲜的洛神花汁。今天打八折！我们的东西又新鲜又便宜。快来买吧！} — six phrases courtes, toutes les étapes sont présentes, aucune structure non maîtrisée.
\end{exoresolu}

\courssub{Méthode : réussir la production écrite au Bac}

\begin{methode}[Rédiger une publicité en cinq gestes]
\begin{enumerate}
\item \textbf{Choisir} un seul produit et noter en marge les 5 ou 6 mots chinois utiles (produit, deux adjectifs, chiffre de la réduction).
\item \textbf{Suivre la trame} en six étapes, une phrase par étape ; ne jamais sauter l'accroche ni l'appel final.
\item \textbf{N'utiliser que des structures maîtrisées} : \zh{本店出售…}, \zh{打…折}, \zh{又…又…}, \zh{快来买吧}. Au Bac, une phrase simple et correcte rapporte plus qu'une phrase ambitieuse et fautive.
\item \textbf{Vérifier chaque ton} en relisant à voix basse : \zh{piányi} (2\textsuperscript{e} ton puis ton léger), \zh{zhé} (2\textsuperscript{e} ton), \zh{mǎi} (3\textsuperscript{e} ton), \zh{luòshénhuā} (4\textsuperscript{e}, 2\textsuperscript{e}, 1\textsuperscript{er}).
\item \textbf{Soigner les caractères étudiés} : écrire lentement \zh{店} (clé \zh{广} + \zh{占}), \zh{折} (clé de la main \zh{扌}), \zh{便} (clé de l'homme \zh{亻}), en respectant l'ordre des traits : gauche avant droite, haut avant bas, horizontal avant vertical.
\end{enumerate}
\end{methode}

\begin{attention}[Erreurs fréquentes des francophones dans cette production]
\begin{itemize}
\item Oublier le second \zh{又} : on écrit \zh{又好便宜} au lieu de \zh{又好又便宜}. La structure exige les deux \zh{又}.
\item Traduire mot à mot « venez vite acheter » en plaçant mal l'adverbe : \zh{快} se place \textbf{avant} le verbe : \zh{快来买吧}, jamais après.
\item Écrire \zh{打十折} en pensant « moins 10 \% » : \zh{打十折} voudrait dire « prix plein » ; « moins 10 \% » se dit \zh{打九折}.
\item Mettre une virgule française dans les prix ou oublier la ponctuation pleine chasse ：，。！
\item Confondre \zh{木槿} (hibiscus syriacus) et \zh{洛神花} (hibiscus sabdariffa = bissap) : pour le bissap camerounais, on écrit \zh{洛神花汁}.
\end{itemize}
\end{attention}

\courssub{Relecture en binôme et présentation orale}

\begin{experience}[Atelier publicité en binôme]
\begin{enumerate}
\item Chaque élève rédige sa mini-publicité (10 minutes), seul, dictionnaire fermé, avec la banque de mots.
\item \textbf{Relecture croisée} : échangez votre texte avec votre voisin. Le relecteur coche trois points : (a) chaque syllabe du pinyin porte-t-elle sa marque de ton ? (b) les caractères \zh{店}, \zh{折}, \zh{便}, \zh{宜} sont-ils complets et corrects ? (c) la structure \zh{又……又……} a-t-elle bien ses deux \zh{又} ?
\item Corrigez ensemble, puis lisez votre publicité à voix haute à votre binôme en soignant les tons.
\item \textbf{Présentation orale} : deux ou trois volontaires présentent leur publicité devant la classe, avec le geste et le sourire du vendeur ; la classe identifie le produit, le montant de la réduction et les deux adjectifs de l'argument.
\end{enumerate}
\end{experience}

\begin{savaistu}[Un acheté, un offert : 买一送一]
En Chine, la promotion la plus célèbre est \zh{买一送一} (mǎi yī sòng yī), littéralement « acheter un, offrir un » : c'est l'équivalent exact des promotions « 2 = 3 » ou « un acheté, un offert » des supermarchés de Douala et de Yaoundé. On la voit partout : sur les boissons, les nouilles, le thé. Les Chinois adorent aussi les chiffres porte-bonheur dans les prix : le 8 (八, bā) évoque la prospérité car il sonne comme 发 (fā, « s'enrichir »), tandis que le 4 (四, sì), proche de 死 (sǐ, « mourir »), est évité. Un commerçant chinois préférera donc afficher 8,8 plutôt qu'un prix avec un 4 !
\end{savaistu}

\courssub{Pour aller plus loin : complément utile à l'examen}

\begin{aretenir}[Formules d'accroche à connaître pour le Bac]
Ces formules reviennent souvent dans les textes publicitaires des sujets d'examen ; sachez au moins les reconnaître :
\begin{itemize}
\item \zh{好消息！} (hǎo xiāoxi !) « Bonne nouvelle ! »
\item \zh{大减价！} (dà jiǎnjià !) « Grande réduction ! »
\item \zh{买一送一！} (mǎi yī sòng yī !) « Un acheté, un offert ! »
\item \zh{欢迎光临！} (huānyíng guānglín !) « Bienvenue ! »
\item \zh{物美价廉} (wù měi jià lián) « belle marchandise, prix doux » (chengyu à quatre caractères, niveau Bac).
\end{itemize}
\end{aretenir}

\begin{exercice}[À vous de jouer : la publicité des pagnes]
Rédigez une mini-publicité de 4 à 6 phrases pour un magasin de pagnes (\zh{花布}) du marché de Douala, en suivant la trame en six étapes : accroche avec \zh{好消息}, présentation du magasin, produit, offre \zh{打七折}, argument avec \zh{又……又……} (adjectifs \zh{漂亮} et \zh{便宜}), appel à l'action. Écrivez les caractères, le pinyin avec tons et la traduction française.
\end{exercice}

\begin{corrige}[Exercice : la publicité des pagnes]
Proposition de corrigé (d'autres formulations correctes sont possibles) :\\
\zh{好消息！欢迎光临！本店出售喀麦隆的花布。今天打七折！我们的花布又漂亮又便宜。快来买吧！}\\
\emph{Hǎo xiāoxi! Huānyíng guānglín! Běn diàn chūshòu Kāmàilóng de huābù. Jīntiān dǎ qī zhé! Wǒmen de huābù yòu piàoliang yòu piányi. Kuài lái mǎi ba!}\\
« Bonne nouvelle ! Bienvenue ! Notre magasin vend des pagnes du Cameroun. Aujourd'hui moins trente pour cent ! Nos pagnes sont beaux et bon marché. Venez vite les acheter ! »\\
Points de vérification : accroche présente ; \zh{打七折} = on paie 70 \% = moins 30 \% ; les deux \zh{又} encadrent bien les deux adjectifs ; l'appel final \zh{快来买吧} est en place.
\end{corrige}

\newpage

\coursec{Activité d'intégration}

\coursec{Leçon 30 — Vocabulaire et compréhension de texte : publicité}

\courssub{Mise en route et découverte du thème}

\begin{experience}[Échange initial : la publicité autour de nous]
\begin{enumerate}
\item Observez les affiches présentes dans la salle de classe, le couloir ou le quartier Mvog-Ada / Mokolo. Relevez trois mots ou expressions qui reviennent souvent (ex. : \zh{促销}, \zh{特价}, \zh{买一送一}).
\item En binôme, discutez en français : qu'est-ce qui rend une publicité efficace pour vous ? (Prix bas, image, slogan, réduction, confiance).
\item Le professeur note au tableau les idées-clés en chinois : \zh{便宜} (pas cher), \zh{好看} (beau), \zh{打折} (réduction), \zh{质量好} (bonne qualité).
\end{enumerate}
\end{experience}

\begin{savaistu}[Publicité en Chine et au Cameroun]
En Chine, les caractères \zh{广告} (guǎnggào, « large annonce ») désignent la publicité. On voit partout des banderoles rouges \zh{横幅} (héngfú) avec des slogans en caractères géants : \zh{质量第一，信誉至上} (Qualité d'abord, réputation suprême). Au Cameroun, les panneaux peints à la main sur les murs des boutiques (à Yaoundé, Douala, Bafoussam) utilisent souvent le français, l'anglais ou le camfranglais : « Promo », « Prix cassé », « Achetez 2, gagnez 1 ». L'activité d'aujourd'hui mêle ces deux cultures : vendre un produit camerounais avec les codes publicitaires chinois.
\end{savaistu}

\courssub{Texte support : lecture et écoute}

\begin{textebox}[Modèle d'affiche : le magasin 好味道]
好味道食品店 \\
Hǎo Wèidào shípǐn diàn \\
« Magasin d'alimentation Bon Goût » \\[0.4em]
欢迎光临！本店卖喀麦隆的特产。 \\
Huānyíng guānglín! Běn diàn mài Kāmàilóng de tèchǎn. \\
Bienvenue ! Notre magasin vend des produits spéciaux du Cameroun. \\[0.4em]
我们的咖啡又香又便宜，一包只要十五块。 \\
Wǒmen de kāfēi yòu xiāng yòu piányi, yì bāo zhǐ yào shíwǔ kuài. \\
Notre café est à la fois parfumé et bon marché : seulement 15 yuans le paquet. \\[0.4em]
今天打八折！买两包送一包。机会难得，快来看看吧！ \\
Jīntiān dǎ bā zhé! Mǎi liǎng bāo sòng yì bāo. Jīhuì nándé, kuài lái kànkan ba! \\
Aujourd'hui, 20 \% de réduction ! Deux paquets achetés, un paquet offert. Occasion rare, venez vite voir !
\end{textebox}

\begin{methode}[Comment lire une affiche publicitaire chinoise]
\begin{enumerate}
\item \textbf{Identifiez le nom du magasin} (souvent 2--3 caractères + \zh{店} / \zh{超市} / \zh{商行}).
\item \textbf{Repérez la formule de bienvenue} : \zh{欢迎光临} (standard), \zh{欢迎惠顾} (plus formel).
\item \textbf{Trouvez le produit} (après \zh{卖} / \zh{有} / \zh{主营}) et ses \textbf{qualités} (structure \zh{又……又……} ou adjectifs simples).
\item \textbf{Localisez le prix} : mot-clé \zh{只要} + nombre + \zh{块} (oral) / \zh{元} (écrit formel) + classificateur (\zh{包}, \zh{瓶}, \zh{斤}, \zh{个}).
\item \textbf{Décodez la promotion} : \zh{打} + chiffre + \zh{折} (réduction) ; \zh{买} + chiffre + \zh{送} + chiffre (offre « achetés / offerts »).
\item \textbf{Terminez par l'appel à l'action} : \zh{快来看看吧}, \zh{欢迎选购}, \zh{别错过}.
\end{enumerate}
\end{methode}

\courssub{Glossaire du thème}

\begin{center}
\begin{tabularx}{\linewidth}{|X|X|X|X|}
\hline
\rowcolor{popblueL} Caractères & Pinyin & Nature & Traduction \\
\hline
广告 & guǎnggào & nom & publicité, annonce \\
\hline
本店 & běn diàn & expr. & notre magasin (humble) \\
\hline
特产 & tèchǎn & nom & produit spécial, spécialité locale \\
\hline
咖啡 & kāfēi & nom & café \\
\hline
香 & xiāng & adj. & parfumé, bon odeur \\
\hline
便宜 & piányi & adj. & bon marché, pas cher \\
\hline
只要 & zhǐ yào & adv. / conj. & il suffit de, seulement (marque le prix bas) \\
\hline
块 & kuài & classif. & yuan (unité monétaire orale) \\
\hline
包 & bāo & classif. & paquet, sachet \\
\hline
打折 & dǎ zhé & v. séparable & faire une réduction (lit. « frapper le rabais ») \\
\hline
送 & sòng & verbe & offrir, donner en cadeau \\
\hline
欢迎光临 & huānyíng guānglín & expr. & bienvenue (aux clients) \\
\hline
机会难得 & jīhuì nándé & expr. & occasion rare, à ne pas manquer \\
\hline
选购 & xuǎngòu & verbe & choisir et acheter \\
\hline
阳光 & yángguāng & nom & soleil, lumière du soleil \\
\hline
菠萝干 & bōluó gān & nom & ananas séché \\
\hline
蜂蜜 & fēngmì & nom & miel \\
\hline
纯 & chún & adj. & pur, authentique \\
\hline
甜 & tián & adj. & sucré \\
\hline
好吃 & hǎochī & adj. & bon, délicieux (litt. « bon à manger ») \\
\hline
瓶 & píng & classif. & bouteille, flacon, pot \\
\hline
\end{tabularx}
\end{center}

\courssub{Clés (radicaux) et ordre des traits}

\begin{propriete}[Analyse de caractères clés du thème]
\begin{center}
\begin{tabularx}{\linewidth}{|X|X|X|X|}
\hline
\rowcolor{popblueL} Caractère & Radical (clé) & Nombre de traits & Ordre des traits (description) \\
\hline
\zh{广} (guǎng) & \zh{广} (yǎnzìtóu, « toit large ») & 3 & 1. Point (丶) 2. Horizontal (一) 3. Pié gauche (ノ) — Le radical couvre le haut. \\
\hline
\zh{本} (běn) & \zh{木} (mù, « bois, arbre ») & 5 & 1. Horizontal 2. Vertical 3. Pié gauche 4. Pié droit 5. Horizontal (base). Le trait horizontal final marque la racine. \\
\hline
\zh{特} (tè) & \zh{牛} (niú, « bœuf ») & 10 & Écrire \zh{牜} (niú à gauche) : horizontal, vertical, pié gauche, pié droit, horizontal (5 traits). Puis \zh{寺} (sì) à droite : horizontal, vertical, horizontal, horizontal, horizontal, vertical, horizontal (7 traits) — total 10. \\
\hline
\zh{便} (pián) & \zh{亻} (rén, « homme ») & 9 & 1. Pié gauche (亻) 2. Vertical 3. Horizontal 4. Pié gauche 5. Point 6. Horizontal 7. Vertical 8. Horizontal 9. Horizontal (partie droite \zh{更}). \\
\hline
\zh{折} (zhé) & \zh{扌} (shǒu, « main ») & 7 & 1. Horizontal 2. Crochet vertical 3. Relevé (提) — radical main. Puis \zh{斤} (jīn) : horizontal, vertical, horizontal, horizontal (4 traits). \\
\hline
\zh{送} (sòng) & \zh{辶} (chuò, « marche ») & 9 & Partie gauche \zh{关} (guān) : 6 traits (haut en bas, gauche à droite). Partie droite \zh{辶} : 3 traits (point, pié gauche,捺 long). Règle : l'intérieur avant la fermeture \zh{辶}. \\
\hline
\end{tabularx}
\end{center}
\end{propriete}

\begin{aretenir}[Rappel : Règles générales d'ordre des traits]
\begin{itemize}
\item De haut en haut (上 → 下) : \zh{三}, \zh{本}.
\item De gauche à droite (左 → 右) : \zh{好}, \zh{特}.
\item Horizontal avant vertical (横 → 竖) : \zh{十}, \zh{本}.
\item Pié gauche avant pié droit (撇 → 捺) : \zh{人}, \zh{八}.
\item Extérieur avant intérieur, puis fermeture : \zh{国}, \zh{送} (écrire \zh{关} puis \zh{辶}).
\item Point en haut à gauche en premier ; point en haut à droite ou au milieu en dernier.
\end{itemize}
\end{aretenir}

\courssub{Compréhension du texte et collocations}

\begin{exercice}[Tâche 1 — Compréhension du modèle (15 min)]
Lisez l'affiche \zh{好味道食品店} et répondez par écrit (français ou chinois) :
\begin{enumerate}
\item Que vend le magasin \zh{好味道} ?
\item Quel est le prix normal d'un paquet de café ?
\item Que signifie \zh{打八折} ? Calculez le prix du paquet aujourd'hui.
\item Quelle offre spéciale le magasin propose-t-il en plus de la réduction ?
\item Relevez dans le texte une phrase construite avec \zh{又……又……} et traduisez-la.
\item \textbf{Collocations :} Associez les mots de la colonne A avec ceux de la colonne B pour former des expressions du texte.
\begin{center}
\begin{tabular}{|c|c|c|c|}
\hline
A & & B & \\
\hline
1. 欢迎 & & a. 折 & \\
2. 打 & & b. 光临 & \\
3. 买两包 & & c. 只要 & \\
4. 一包 & & d. 送一包 & \\
5. 机会 & & e. 难得 & \\
\hline
\end{tabular}
\end{center}
\end{enumerate}
\end{exercice}

\begin{corrige}[Exercice 1]
\begin{enumerate}
\item Le magasin vend des \zh{特产} (spécialités) du Cameroun, dont du \zh{咖啡} (café).
\item 15 \zh{块} (kuài) le paquet (\zh{一包只要十五块}).
\item \zh{打八折} = payer 80 \% du prix initial. 20 \% de réduction. Prix aujourd'hui : 15 $\times$ 0,8 = 12 \zh{块}.
\item \zh{买两包送一包} (Achetez 2 paquets, 1 offert).
\item \zh{我们的咖啡又香又便宜} — « Notre café est à la fois parfumé et bon marché. »
\item Correspondances : 1-b, 2-a, 3-d, 4-c, 5-e.
\end{enumerate}
\end{corrige}

\begin{propriete}[Structure : \zh{又……又……} (à la fois… et…)]
\begin{itemize}
\item \textbf{Schéma} : Sujet + \zh{又} + Adj.1 + \zh{又} + Adj.2.
\item \textbf{Emploi} : Énumère deux qualités (souvent positives) simultanées. Les adjectifs sont généralement monosyllabiques ou disyllabiques de même longueur.
\item \textbf{Exemples} :
\begin{itemize}
\item \zh{这个苹果又大又甜。} (Zhège píngguǒ yòu dà yòu tián.) — Cette pomme est à la fois grosse et sucrée.
\item \zh{他又高又帅。} (Tā yòu gāo yòu shuài.) — Il est à la fois grand et beau gosse.
\end{itemize}
\item \textbf{Interdiction} : Ne pas ajouter \zh{很} (hěn) devant. \textbf{*很又大又甜} (incorrect).
\end{itemize}
\end{propriete}

\begin{propriete}[Structure : Prix avec \zh{只要} (seulement / il suffit de)]
\begin{itemize}
\item \textbf{Schéma} : Produit + Classificateur + \zh{只要} + Nombre + \zh{块 / 元}.
\item \textbf{Exemples} :
\begin{itemize}
\item \zh{一斤只要十块。} (Yì jīn zhǐ yào shí kuài.) — Un jin (500 g) coûte seulement 10 yuans.
\item \zh{这件衣服只要五十元。} (Zhè jiàn yīfu zhǐ yào wǔshí yuán.) — Ce vêtement ne coûte que 50 yuans.
\end{itemize}
\item \textbf{Nuance} : \zh{只要} souligne le caractère avantageux / bas du prix. \zh{是} (shì) est neutre : \zh{一包是十五块}.
\end{itemize}
\end{propriete}

\begin{propriete}[Structure : Réduction avec \zh{打……折}]
\begin{itemize}
\item \textbf{Schéma} : \zh{打} + Chiffre (1--9) + \zh{折}.
\item \textbf{Règle de calcul} : Le chiffre indique le \textbf{pourcentage du prix initial à payer}.
\begin{itemize}
\item \zh{打八折} (dǎ bā zhé) = payer 80 \% = \textbf{--20 \%}.
\item \zh{打五折} (dǎ wǔ zhé) = payer 50 \% = \textbf{--50 \%} (moitié prix).
\item \zh{打九折} (dǎ jiǔ zhé) = payer 90 \% = \textbf{--10 \%}.
\end{itemize}
\item \textbf{Attention} : En français « 80 \% de réduction » = payer 20 \% = \zh{打二折} (rare). Ne confondez pas le chiffre chinois et le pourcentage de réduction français.
\end{itemize}
\end{propriete}

\begin{propriete}[Structure : Offre « Achetés / Offerts » \zh{买……送……}]
\begin{itemize}
\item \textbf{Schéma} : \zh{买} + Nombre A + (Classif.) + \zh{送} + Nombre B + (Classif.).
\item \textbf{Exemples} :
\begin{itemize}
\item \zh{买两送一} (Mǎi liǎng sòng yī) — 2 achetés, 1 offert.
\item \zh{买三送一} (Mǎi sān sòng yī) — 3 achetés, 1 offert.
\item \zh{买一送一} (Mǎi yī sòng yī) — 1 acheté, 1 offert (50 \% de réduction effective sur le 2e).
\end{itemize}
\item \textbf{Écriture publicitaire} : Le classificateur est souvent omis pour la brièveté (\zh{买二送一}).
\end{itemize}
\end{propriete}

\begin{attention}[Pièges fréquents pour francophones]
\begin{itemize}
\item \textbf{\zh{打八折} $\neq$ « 80 \% de réduction »}. C'est 20 \% de réduction. Le chiffre = \% payé.
\item \textbf{\zh{又……又……} sans \zh{很}}. \zh{很} est incompatible avec cette structure.
\item \textbf{\zh{买} (mǎi, 3\textsuperscript{e} ton) vs \zh{卖} (mài, 4\textsuperscript{e} ton)}. \zh{本店卖} (Notre magasin vend) ; \zh{我买} (J'achète). Confusion fatale à l'oral !
\item \textbf{Classificateurs obligatoires} avec les nombres : \zh{一包} café, \zh{一瓶} miel, \zh{一斤} fruits, \zh{一条} poisson / pagne. Pas de \zh{一个} café.
\item \textbf{Ponctuation} : Utilisez la ponctuation chinoise pleine chasse (，。！？) dans les phrases chinoises. Pas d'espace entre caractères.
\end{itemize}
\end{attention}

\courssub{Production guidée : ma publicité}

\begin{methode}[Étapes de rédaction de votre affiche (Tâches 2 \& 3)]
\begin{enumerate}
\item \textbf{Choisir le produit} (café, cacao, ndolé en conserve, ananas séché, huile de palme, pagne, miel, bananes plantain séchées…).
\item \textbf{Remplir la fiche vocabulaire} (produit, 2 adjectifs qualité, prix normal \zh{块}, réduction \zh{打……折}, offre \zh{买……送……}).
\item \textbf{Rédiger selon la trame en 7 points} (voir ci-dessous).
\item \textbf{Vérifier} : ordre des mots (SVO), classificateurs, tons (\zh{mài/mǎi}, \zh{zhé}, \zh{piányi}), ponctuation chinoise.
\item \textbf{Écrire au propre} sur format A4/A3, caractères soignés, pinyin au dos pour l'oral.
\end{enumerate}
\end{methode}

\begin{textebox}[Trame obligatoire de l'affiche (7 points)]
\begin{enumerate}
\item Nom du magasin (2--3 caractères) + \zh{店} \quad (ex. \zh{喀麦隆阳光店})
\item \zh{欢迎光临！}
\item Présentation : \zh{本店卖……} ou \zh{我们的……}
\item Qualité : \zh{又……又……} (deux adjectifs)
\item Prix : \zh{……只要……块。} (avec classificateur)
\item Réduction : \zh{今天打……折！}
\item Appel à l'action : \zh{快来看看吧！} ou \zh{欢迎选购！}
\end{enumerate}
\end{textebox}

\begin{experience}[Tâche 2 — Choix du produit et fiche vocabulaire (15 min)]
En groupe de 4, complétez ce tableau au brouillon :
\begin{center}
\begin{tabularx}{\linewidth}{|X|X|}
\hline
\rowcolor{poporangeL} Élément & Votre choix (caractères + pinyin) \\
\hline
Produit camerounais & \\
\hline
Adjectif qualité 1 & \\
\hline
Adjectif qualité 2 & \\
\hline
Prix normal (块) + classif. & \\
\hline
Réduction (打…折) & \\
\hline
Offre (买…送…) & \\
\hline
Nom du magasin (+ 店) & \\
\hline
\end{tabularx}
\end{center}
\end{experience}

\begin{experience}[Tâche 3 — Rédaction de l'affiche (40 min)]
Rédigez l'affiche finale sur une grande feuille. Respectez la trame. Le professeur circule pour corriger les caractères (ordre des traits) et la syntaxe.
\end{experience}

\begin{experience}[Tâche 4 — Présentation orale et questions (30 min)]
\begin{enumerate}
\item Chaque groupe lit son affiche à voix haute (prononciation, tons, fluidité).
\item Le groupe pose \textbf{deux questions de compréhension} à la classe (ex. \zh{我们的产品是什么？}, \zh{今天打几折？}, \zh{一包多少钱？}).
\item La classe répond en chinois (phrase complète si possible).
\end{enumerate}
\end{experience}

\begin{experience}[Tâche 5 — Vote et justification (20 min)]
\begin{enumerate}
\item Vote à bulletin secret pour la publicité la plus convaincante (critères : exactitude langue, respect trame, originalité, persuasion).
\item Chaque élève justifie son vote en chinois : \zh{我觉得……组的广告最好，因为……} (Wǒ juéde … zǔ de guǎnggào zuì hǎo, yīnwèi…).
\item Exemples de justifications : \zh{因为价格便宜} / \zh{因为又香又甜} / \zh{因为打五折，很划算} (huásuàn, rentable).
\end{enumerate}
\end{experience}

\courssub{Proposition de production et corrigé commenté}

\begin{exoresolu}[Affiche modèle : le magasin 喀麦隆阳光]
Énoncé : Rédigez l'affiche publicitaire complète du groupe qui vend des ananas séchés et du miel du Cameroun, en respectant la trame en sept points.
\tcblower
\textbf{Production modèle :} \\[0.3em]
\zh{喀麦隆阳光食品店} \\
Kāmàilóng Yángguāng shípǐn diàn \\
« Magasin d'alimentation Soleil du Cameroun » \\[0.3em]
\zh{欢迎光临！本店卖喀麦隆的特产：菠萝干和蜂蜜。} \\
Huānyíng guānglín! Běn diàn mài Kāmàilóng de tèchǎn: bōluó gān hé fēngmì. \\
Bienvenue ! Notre magasin vend des spécialités du Cameroun : ananas séché et miel. \\[0.3em]
\zh{我们的菠萝干又甜又好吃，一包只要十块。} \\
Wǒmen de bōluó gān yòu tián yòu hǎochī, yì bāo zhǐ yào shí kuài. \\
Notre ananas séché est à la fois sucré et délicieux : seulement 10 yuans le paquet. \\[0.3em]
\zh{我们的蜂蜜又纯又香，一瓶二十块。} \\
Wǒmen de fēngmì yòu chún yòu xiāng, yì píng èrshí kuài. \\
Notre miel est à la fois pur et parfumé : 20 yuans le flacon. \\[0.3em]
\zh{今天打九折！买三送一。机会难得，快来看看吧！} \\
Jīntiān dǎ jiǔ zhé! Mǎi sān sòng yī. Jīhuì nándé, kuài lái kànkan ba! \\
Aujourd'hui, 10 \% de réduction ! Trois achetés, un offert. Occasion rare, venez vite ! \\[0.5em]
\textbf{Commentaire des choix de langue :}
\begin{itemize}
\item Le nom \zh{阳光} (« soleil ») évoque le climat camerounais et une image positive (chaleur, naturel) : stratégie de \cle{naming} courante en Chine.
\item Classificateurs adaptés : \zh{包} (bāo, paquet) pour l'ananas séché ; \zh{瓶} (píng, bouteille/pot) pour le miel.
\item \zh{又……又……} utilisé deux fois avec des paires d'adjectifs distincts (\zh{甜/好吃}, \zh{纯/香}) pour varier le lexique.
\item \zh{打九折} = 10 \% de réduction (paye 90 \%), cohérent pour une petite promotion.
\item \zh{买三送一} = abréviation publicitaire standard de \zh{买三个送一个}.
\item Particule finale \zh{吧} dans \zh{快来看看吧} : adoucit l'invitation, ton engageant.
\end{itemize}
\end{exoresolu}

\courssub{Bilan de l'activité et auto-évaluation}

\begin{aretenir}[Bilan des compétences travaillées]
\begin{itemize}
\item \cle{Compréhension de l'écrit} : Lecture d'une affiche authentique, extraction d'informations (prix, réduction, offre, produit).
\item \cle{Vocabulaire} : Lexique de la publicité (\zh{本店, 特产, 便宜, 打折, 送, 欢迎光临, 机会难得, 选购}) et classificateurs (\zh{包, 瓶, 斤}).
\item \cle{Grammaire} : Structures \zh{又……又……}, \zh{只要……块}, \zh{打……折}, \zh{买……送……}, ordre des mots SVO.
\item \cle{Expression écrite} : Rédaction d'un texte court fonctionnel, ponctuation chinoise, ordre des traits soigné.
\item \cle{Expression orale} : Lecture à voix haute, tons (\zh{mǎi/mài}, \zh{zhé}, \zh{piányi}), interaction questions/réponses.
\item \cle{Interculturalité} : Adaptation de produits camerounais aux codes publicitaires chinois (nom évocateur, formules fixes, offres groupées).
\end{itemize}
\end{aretenir}

\begin{methode}[Auto-évaluation : mon affiche est-elle réussie ?]
Cochez les cases :
\begin{itemize}
\item[ ] Les 7 éléments de la trame sont présents et dans l'ordre.
\item[ ] Chaque phrase suit l'ordre Sujet + Verbe + Complément.
\item[ ] Les classificateurs sont corrects (\zh{包, 瓶, 斤, 条, 个}).
\item[ ] Les tons des mots-clés sont maîtrisés (\zh{mài} vendre, \zh{mǎi} acheter, \zh{zhé} réduction, \zh{piányi} bon marché).
\item[ ] La ponctuation est chinoise (，。！？).
\item[ ] L'affiche est lisible, caractères bien formés.
\item[ ] J'ai su répondre aux questions des camarades en chinois.
\end{itemize}
Si tout est coché : \zh{你做得很好！} (Nǐ zuò de hěn hǎo ! — Tu as très bien travaillé !)
\end{methode}

\newpage

\coursec{Méthodes et exercices résolus}

\courssub{Méthodes et exercices résolus}

\begin{textebox}[Texte support : Publicité pour le smartphone « \zh{星耀 X20} »]
\zh{全新 星耀 X20，专为 学生 设计！} \\
\zh{超大 内存，运行 游戏 不 卡顿；} \\
\zh{长续航 电池，看 视频 两 天 不 充电。} \\
\zh{拍照 超 清晰，夜景 也 像 白天 一样 亮。} \\
\zh{价格 实惠，仅需 129 900 FCFA！} \\
\zh{现在 购买 赠送 耳机 一 副，数量 有限，先到 先得。} \\
\zh{地址：雅温得 华商 大厦 一楼；} \\
\zh{电话：6XX XX XX XX。} \\
\zh{星耀 X20，让 学习 娱乐 两 不 误，快来 体验 吧！}
\end{textebox}

\begin{methode}[Lire et comprendre une publicité : repérer l'information essentielle]
\textbf{Objectif :} Identifier rapidement le produit, ses arguments clés, le prix et l'appel à l'action dans un texte court et dense.
\begin{enumerate}
    \item \textbf{Balayer le titre et la première phrase :} Repérer le nom du produit (\zh{星耀 X20}) et la cible (\zh{专为学生设计}).
    \item \textbf{Repérer les mots-clés « arguments » :} Chercher les adjectifs de qualité (\zh{超大}, \zh{长续航}, \zh{超清晰}, \zh{实惠}) et les structures \zh{不...} (négation de résultat) ou \zh{也...} (concession).
    \item \textbf{Localiser les informations pratiques :} Prix (\zh{仅需...}), offre promotionnelle (\zh{赠送...}), lieu (\zh{地址...}), contact (\zh{电话...}).
    \item \textbf{Analyser la chute (slogan/appel) :} La dernière phrase résume le bénéfice global (\zh{两不误}) et utilise une particule incitative (\zh{吧}).
\end{enumerate}
\cle{Astuce} : Dans une publicité chinoise, la structure est souvent \textbf{Produit $\rightarrow$ Caractéristiques (Parallélisme) $\rightarrow$ Prix/Offre $\rightarrow$ Coordonnées $\rightarrow$ Slogan}.
\end{methode}

\begin{exoresolu}[Compréhension globale : Questions ouvertes]
Répondez en français aux questions suivantes sur le texte support.
\begin{enumerate}
    \item Quel est le produit annoncé et quel est son public cible principal ?
    \item Citez trois arguments techniques avancés pour séduire l'acheteur.
    \item Quelle est l'offre promotionnelle limitée dans le temps/quantité ?
    \item Où peut-on acheter ce téléphone ? Donnez l'adresse exacte.
    \item Traduisez la dernière phrase : \zh{星耀 X20，让 学习 娱乐 两 不 误，快来 体验 吧！}
\end{enumerate}
\tcblower
\textbf{Solution détaillée :}
\begin{enumerate}
    \item \textbf{Produit :} Smartphone \zh{星耀 X20} (Xingyao X20). \textbf{Cible :} \zh{学生} (étudiants/lycéens) — indiqué par \zh{专为学生设计} (conçu spécialement pour les étudiants).
    \item \textbf{Arguments techniques :}
        \begin{itemize}
            \item \zh{超大内存，运行游戏不卡顿} : Mémoire vive énorme, jeux fluides (pas de lag).
            \item \zh{长续航电池，看视频两天不充电} : Batterie longue durée, 2 jours vidéo sans recharge.
            \item \zh{拍照超清晰，夜景也像白天一样亮} : Photos ultra-nettes, vision nocturne claire comme le jour.
        \end{itemize}
    \item \textbf{Offre :} \zh{赠送耳机一副} (un casque/écouteurs offert) pour tout achat, \zh{数量有限，先到先得} (quantités limitées, premier arrivé premier servi).
    \item \textbf{Adresse :} \zh{雅温得华商大厦一楼} (1er étage de l'immeuble Huashang à Yaoundé). \zh{雅温得} = Yaoundé ; \zh{华商大厦} = Immeuble Huashang (nom propre) ; \zh{一楼} = 1er étage (RDC en Chine, souvent 1er étage au Cameroun).
    \item \textbf{Traduction :} « Le Xingyao X20, pour que les études et les loisirs ne se gênent pas mutuellement [soient compatibles], venez vite l'essayer ! »
        \begin{itemize}
            \item \zh{让} (ràng) : faire en sorte que / laisser.
            \item \zh{两不误} (liǎng bù wù) : expression idiomatique « les deux ne s'empêchent pas l'un l'autre » = concilier les deux.
            \item \zh{快来} (kuài lái) : venez vite.
            \item \zh{吧} (ba) : particule de suggestion/incitation douce.
        \end{itemize}
\end{enumerate}
\cle{Conclusion} : La lecture efficace d'une pub repose sur l'extraction des \textbf{informations factuelles} (produit, prix, lieu) et des \textbf{arguments persuasifs} (adjectifs + structures résultat/négation).
\end{exoresolu}

\begin{methode}[Maîtriser le vocabulaire thématique : radicaux, collocations et classificateurs]
\textbf{Objectif :} Mémoriser le lexique de la publicité/économie en comprenant la logique des caractères (clés phonétiques/sémantiques) et en apprenant les \textbf{collocations} (associations mots-mots) et \textbf{classificateurs} obligatoires.
\begin{enumerate}
    \item \textbf{Analyser le radical (clé) :} Il donne l'indice sémantique. Ex: \zh{电} (électricité) dans \zh{电池} (batterie), \zh{视频} (vidéo), \zh{电话} (téléphone).
    \item \textbf{Mémoriser les collocations verbales :} En chinois, les verbes s'associent à des objets spécifiques. Ex: \zh{运行} (faire tourner) + \zh{游戏/软件} ; \zh{拍照} (prendre photo) + \zh{夜景/风景} ; \zh{赠送} (offrir) + \zh{礼品/耳机}.
    \item \textbf{Intégrer le classificateur :} Jamais de nombre nu. \zh{一副} 耳机 (une paire d'écouteurs), \zh{一台} 手机 (un téléphone), \zh{一款} 产品 (un modèle de produit).
    \item \textbf{Travailler les paires synonymes/antonymes :} \zh{实惠} (bon rapport qualité/prix) $\neq$ \zh{贵} ; \zh{卡顿} (laguer) $\neq$ \zh{流畅} (fluide).
\end{enumerate}
\end{methode}

\begin{exoresolu}[Vocabulaire : Radicaux, classificateurs et collocations]
\begin{enumerate}
    \item Complétez le tableau ci-dessous pour les mots extraits du texte.
    \item Choisissez le bon classificateur parmi \zh{个, 台, 副, 套, 款, 条} pour : \zh{手机}, \zh{耳机}, \zh{电池}, \zh{广告}, \zh{产品}, \zh{裤子}.
    \item Reliez le verbe à son objet logique (collocation) :
    \begin{center}
    \begin{tabular}{cc}
    \zh{运行} & \zh{耳机} \\
    \zh{拍摄} & \zh{游戏} \\
    \zh{赠送} & \zh{视频} \\
    \zh{观看} & \zh{夜景} \\
    \zh{佩戴} & \zh{礼品}
    \end{tabular}
    \end{center}
\end{enumerate}
\tcblower
\textbf{Solution détaillée :}
\begin{enumerate}
    \item \textbf{Tableau vocabulaire + radicaux :}
    \begin{center}
    \begin{tabularx}{\linewidth}{|X|X|X|X|X|}
    \hline \rowcolor{popblueL} \textbf{Caractères} & \textbf{Pinyin} & \textbf{Radical (Clé)} & \textbf{Sens du radical} & \textbf{Traduction} \\ \hline
    \zh{电池} & diànchí & \zh{电} (diàn) & Électricité / Éclair & Batterie \\ \hline
    \zh{视频} & shìpín & \zh{见} (jiàn) \textit{(dans 视)} & Voir / Regarder & Vidéo \\ \hline
    \zh{续航} & xùháng & \zh{舟} (zhōu) \textit{(dans 航)} & Bateau / Navigation & Autonomie (véhicule/tél) \\ \hline
    \zh{清晰} & qīngxī & \zh{氵} (shuǐ) \textit{(dans 清)} & Eau / Clarté & Clair, net \\ \hline
    \zh{实惠} & shíhuì & \zh{宀} (mián) \textit{(dans 实)} & Toit / Réel & Avantageux, bon prix \\ \hline
    \zh{赠送} & zèngsòng & \zh{贝} (bèi) \textit{(dans 赠)} & Coquillage / Valeur & Offrir (cadeau) \\ \hline
    \zh{数量} & shùliàng & \zh{攵} (pū) \textit{(dans 数)} & Action / Compter & Quantité \\ \hline
    \end{tabularx}
    \end{center}
    \item \textbf{Classificateurs corrects :}
    \begin{itemize}
        \item \zh{一台手机} (appareil électronique) — \zh{台} pour machines/appareils.
        \item \zh{一副耳机} (paire) — \zh{副} pour paires (lunettes, gants, écouteurs).
        \item \zh{一块/节电池} — \zh{块} (morceau) ou \zh{节} (section) pour piles/batteries.
        \item \zh{一条广告} — \zh{条} pour textes longs, annonces, nouvelles.
        \item \zh{一款产品} — \zh{款} pour modèles, styles, articles de commerce.
        \item \zh{一条裤子} — \zh{条} pour objets longs et fins (pantalons, routes, rivières).
    \end{itemize}
    \item \textbf{Collocations (Verbe + Objet) :}
    \begin{itemize}
        \item \zh{运行游戏} / \zh{运行软件} (faire tourner un jeu/logiciel).
        \item \zh{拍摄夜景} / \zh{拍摄视频} (filmer/photographier une vue de nuit/vidéo).
        \item \zh{赠送礼品} / \zh{赠送耳机} (offrir un cadeau/écouteurs).
        \item \zh{观看视频} / \zh{观看比赛} (regarder une vidéo/un match).
        \item \zh{佩戴耳机} / \zh{佩戴眼镜} (porter/mettre des écouteurs/lunettes).
    \end{itemize}
    \cle{Point de grammaire} : \zh{运行} s'utilise pour logiciels/systèmes ; \zh{播放} (bōfàng) pour vidéos/musique ; \zh{观看} pour regarder activement (écran, spectacle).
\end{enumerate}
\end{exoresolu}

\begin{methode}[Compréhension détaillée : Vrai/Faux et Justification par le texte]
\textbf{Objectif :} Valider une compréhension fine en distinguant l'explicite de l'implicite, et en citant la preuve textuelle (caractères + pinyin).
\begin{enumerate}
    \item \textbf{Lire l'énoncé} et souligner les mots-clés (noms, chiffres, adjectifs, négations).
    \item \textbf{Scanner le texte} pour localiser la zone correspondante (repérage visuel des caractères).
    \item \textbf{Comparer strictement :} Le texte dit-il \textbf{exactement} cela ? Attention aux pièges : \zh{不卡顿} (pas de lag) $\neq$ \zh{很快} (très vite) ; \zh{两天} (deux jours) $\neq$ \zh{一周}.
    \item \textbf{Rédiger la justification :} « Faux. Le texte dit : \zh{看视频两天不充电} (kàn shìpín liǎng tiān bù chōngdiàn). »
    \item \textbf{Gérer l'inférence culturelle :} \zh{实惠} implique « pas cher pour la qualité », pas « gratuit ».
\end{enumerate}
\end{methode}

\begin{exoresolu}[Compréhension détaillée : Vrai ou Faux ? Justifiez par une citation du texte.]
\begin{enumerate}
    \item Le téléphone est conçu pour les personnes âgées.
    \item La batterie permet de regarder des vidéos pendant trois jours sans recharger.
    \item Les photos de nuit sont aussi claires que celles de jour.
    \item Le prix est de 129 900 Francs CFA.
    \item Les écouteurs offerts sont sans fil (Bluetooth).
    \item On peut acheter le téléphone au deuxième étage de l'immeuble Huashang.
\end{enumerate}
\tcblower
\textbf{Solution détaillée :}
\begin{enumerate}
    \item \textbf{Faux.} Texte : \zh{专为学生设计} (zhuān wéi xuéshēng shèjì) — « Conçu spécialement pour les étudiants ».
    \item \textbf{Faux.} Texte : \zh{看视频两天不充电} (kàn shìpín \textbf{liǎng} tiān bù chōngdiàn) — « Deux jours », pas trois.
    \item \textbf{Vrai.} Texte : \zh{夜景也像白天一样亮} (yèjǐng yě xiàng báitiān yīyàng liàng) — « La nuit est aussi brillante que le jour ». Structure \zh{像...一样...}.
    \item \textbf{Vrai.} Texte : \zh{仅需129900 FCFA} (jǐn xū 129900 FCFA) — « Il ne faut que 129 900 FCFA ». \zh{仅需} = « seulement besoin de ».
    \item \textbf{Faux (Non mentionné).} Texte : \zh{赠送耳机一副} — Type non précisé (filaire ou Bluetooth). Ne pas inférer.
    \item \textbf{Faux.} Texte : \zh{雅温得华商大厦一楼} (Yǎwēndé Huáshāng Dàshà \textbf{yīlóu}) — « 1er étage », pas deuxième (\zh{二楼}).
\end{enumerate}
\cle{Conseil Bac} : Pour « Non mentionné », ne pas cocher « Faux » si l'info n'est pas contredite explicitement, mais ici l'exercice impose V/F. Si « Non mentionné » n'est pas option, cocher « Faux » car l'affirmation n'est pas prouvée par le texte.
\end{exoresolu}

\begin{methode}[Analyser la langue de la persuasion : structures grammaticales publicitaires]
\textbf{Objectif :} Reconnaître et reproduire les structures syntaxiques types de l'argumentation commerciale (comparatif, superlatif, conditionnel, impératif doux).
\begin{enumerate}
    \item \textbf{Comparatif \zh{比} (bǐ) / \zh{更} (gèng) / \zh{最} (zuì) :} \zh{A 比 B + Adj} (A plus que B) ; \zh{更 + Adj} (encore plus) ; \zh{最 + Adj} (le plus). \textit{Ex: \zh{比同价位手机更流畅}.}
    \item \textbf{Structure résultat / degré : \zh{Adj + 不 + Verbe} / \zh{Adj + 极了} :} \zh{不卡顿} (pas de lag = résultat négatif absent) ; \zh{便宜极了} (super pas cher).
    \item \textbf{Conditionnel commercial \zh{如果...就...} (rúguǒ... jiù...) :} \zh{如果现在买，就送耳机}. Souvent elliptique : \zh{现在买，送耳机}.
    \item \textbf{Particules modales finales :} \zh{吧} (suggestion/incitation : \zh{快来体验吧}) ; \zh{啊/呀} (exclamation/émotion : \zh{真便宜啊!}) ; \zh{呢} (relance : \zh{你要不要呢?}).
    \item \textbf{Parallélisme (排比) :} Séries de phrases courtes, même structure, rythme percutant. \zh{内存大，电池强，拍照美，价格低.}
\end{enumerate}
\end{methode}

\begin{exoresolu}[Grammaire du texte : Identification et transformation]
\begin{enumerate}
    \item Identifiez dans le texte support une phrase illustrant : a) Le comparatif implicite ; b) La structure résultat (négation) ; c) L'ellipse conditionnelle ; d) La particule d'incitation.
    \item Transformez les phrases suivantes en utilisant la structure indiquée :
    \begin{itemize}
        \item « Ce téléphone est très rapide. » $\rightarrow$ \zh{这个手机很快。}
\end{itemize}
\end{enumerate}
\end{exoresolu}


\newpage

\coursec{Exercices}

\courssub{Je vérifie mes connaissances}

\begin{exercice}[QCM : le vocabulaire de la publicité]
Pour chaque question, entoure la bonne réponse.
\begin{enumerate}
\item \zh{广告} signifie : a) le magasin \quad b) la publicité \quad c) le prix
\item \zh{打折} signifie : a) faire une réduction \quad b) fermer le magasin \quad c) acheter beaucoup
\item \zh{便宜} est le contraire de : a) \zh{好} \quad b) \zh{贵} \quad c) \zh{新}
\item « Venez vite ! » se dit : a) \zh{快来吧！} \quad b) \zh{慢走！} \quad c) \zh{再见！}
\item \zh{商品} désigne : a) les clients \quad b) les marchandises \quad c) les vendeurs
\end{enumerate}
\end{exercice}

\begin{exercice}[Vrai ou faux]
Lis le texte support de la leçon, puis indique si chaque phrase est vraie ou fausse. Justifie ta réponse en citant une phrase du texte.
\begin{enumerate}
\item \zh{今天所有商品打五折。} (Jīntiān suǒyǒu shāngpǐn dǎ wǔ zhé.)
\item \zh{这个商店的水果很新鲜。} (Zhège shāngdiàn de shuǐguǒ hěn xīnxiān.)
\item \zh{商店晚上七点关门。} (Shāngdiàn wǎnshang qī diǎn guān mén.)
\item \zh{买两件衣服，第三件不要钱。} (Mǎi liǎng jiàn yīfu, dì-sān jiàn bú yào qián.)
\end{enumerate}
\end{exercice}

\begin{exercice}[Complète avec le bon mot]
Complète chaque phrase avec le mot qui convient : \zh{广告} — \zh{打折} — \zh{质量} — \zh{顾客} — \zh{特价}.
\begin{enumerate}
\item 我在电视上看到了一个新手机的 \zh{＿＿＿＿} 。
\item 这家商店的衣服 \zh{＿＿＿＿} 很好，也不贵。
\item 今天牛奶 \zh{＿＿＿＿} ，一升只要五百法郎。
\item 周末商店里 \zh{＿＿＿＿} 很多。
\item 过新年的时候，很多商店都 \zh{＿＿＿＿} 。
\end{enumerate}
\end{exercice}

\begin{exercice}[Associe le pinyin aux caractères]
Relie chaque caractère ou mot à son pinyin.
\begin{center}
\begin{tabularx}{\linewidth}{|X|X|}
\hline
\rowcolor{popblueL} Caractères & Pinyin \\
\hline
1. 商店 & a. guǎnggào \\
\hline
2. 广告 & b. dǎzhé \\
\hline
3. 打折 & c. shāngdiàn \\
\hline
4. 新鲜 & d. piányi \\
\hline
5. 便宜 & e. xīnxiān \\
\hline
\end{tabularx}
\end{center}
\end{exercice}

\courssub{Je m'entraîne}

\begin{exercice}[Traduis en chinois]
Traduis les phrases suivantes en chinois simplifié, avec le pinyin.
\begin{enumerate}
\item Notre magasin vend des fruits et des boissons ; aujourd'hui moins vingt pour cent, venez vite !
\item Ces chaussures sont à la fois jolies et bon marché.
\item Tous les produits sont en promotion ce week-end.
\item La qualité de ce téléphone est très bonne, mais il est un peu cher.
\end{enumerate}
\end{exercice}

\begin{exercice}[Transforme avec 又……又……]
Transforme chaque paire de phrases en une seule phrase avec la structure \zh{又……又……} (yòu... yòu...).
\begin{enumerate}
\item 这个商店很大。这个商店很干净。
\item 这里的苹果很便宜。这里的苹果很新鲜。
\item 这件衣服很好看。这件衣服很舒服。
\item 那个广告很有意思。那个广告很有用。
\end{enumerate}
\end{exercice}

\begin{exercice}[Texte à trous]
Complète la publicité suivante avec les mots de la liste : \zh{欢迎} — \zh{打} — \zh{又} — \zh{从} — \zh{了} — \zh{吧}.

\begin{textebox}[Publicité à compléter]
好消息！＿＿＿＿ 今天起，本店所有水果 ＿＿＿＿ 八折。我们的水果 ＿＿＿＿ 新鲜 ＿＿＿＿ 便宜。买 ＿＿＿＿ 还送饮料！＿＿＿＿ 光临！快来 ＿＿＿＿ ！\\
(Hǎo xiāoxi ! Cóng jīntiān qǐ, běn diàn suǒyǒu shuǐguǒ dǎ bā zhé. Wǒmen de shuǐguǒ yòu xīnxiān yòu piányi. Mǎi le hái sòng yǐnliào ! Huānyíng guānglín ! Kuài lái ba !)\\
« Bonne nouvelle ! À partir d'aujourd'hui, tous les fruits de notre magasin sont à moins vingt pour cent. Nos fruits sont à la fois frais et bon marché. Achetez et recevez une boisson en cadeau ! Bienvenue ! Venez vite ! »
\end{textebox}
\end{exercice}

\begin{exercice}[Remets la publicité dans l'ordre]
Les six phrases de cette publicité sont mélangées. Remets-les dans l'ordre logique : accroche, présentation du produit, avantages, prix, réduction, appel à l'action.
\begin{enumerate}[label=\Alph*.]
\item 又好吃，又有营养，老人孩子都喜欢。 (Yòu hǎochī, yòu yǒu yíngyǎng, lǎorén háizi dōu xǐhuan.)
\item 快来买吧，机会难得！ (Kuài lái mǎi ba, jīhuì nándé !)
\item 好消息！新面包到了！ (Hǎo xiāoxi ! Xīn miànbāo dào le !)
\item 一个只要三百法郎。 (Yí ge zhǐ yào sānbǎi fǎláng.)
\item 我们的面包每天早上现做。 (Wǒmen de miànbāo měi tiān zǎoshang xiàn zuò.)
\item 今天买两个打九折。 (Jīntiān mǎi liǎng ge dǎ jiǔ zhé.)
\end{enumerate}
\end{exercice}

\courssub{Je me prépare au Bac}

\begin{exercice}[Compréhension écrite : une annonce de supermarché]
Lis le texte suivant, puis réponds aux questions en chinois, par des phrases complètes.

\begin{textebox}[好又多超市 (Hǎoyòuduō Chāoshì)]
好又多超市大促销！这个星期六和星期天，本超市所有商品打八折。水果又新鲜又便宜：一公斤芒果只要一千法郎。饮料买二送一。我们还有各种衣服和鞋子，质量好，价格低。超市早上八点开门，晚上九点关门。地址：雅温得市中心，邮局旁边。欢迎光临！\\
(Hǎoyòuduō chāoshì dà cùxiāo ! Zhège xīngqīliù hé xīngqītiān, běn chāoshì suǒyǒu shāngpǐn dǎ bā zhé. Shuǐguǒ yòu xīnxiān yòu piányi : yì gōngjīn mángguǒ zhǐ yào yìqiān fǎláng. Yǐnliào mǎi èr sòng yī. Wǒmen hái yǒu gè zhǒng yīfu hé xiézi, zhìliàng hǎo, jiàgé dī. Chāoshì zǎoshang bā diǎn kāi mén, wǎnshang jiǔ diǎn guān mén. Dìzhǐ : Yǎwēndé shì zhōngxīn, yóujú pángbiān. Huānyíng guānglín !)\\
« Grande promotion au supermarché Hǎoyòuduō ! Ce samedi et ce dimanche, tous les articles du supermarché sont à moins vingt pour cent. Les fruits sont à la fois frais et bon marché : un kilo de mangues à mille francs seulement. Deux boissons achetées, une offerte. Nous avons aussi toutes sortes de vêtements et de chaussures, de bonne qualité et à bas prix. Le supermarché ouvre à 8 h et ferme à 21 h. Adresse : centre-ville de Yaoundé, à côté de la poste. Bienvenue ! »
\end{textebox}

\begin{enumerate}
\item 超市什么时候有促销？ (Quand le supermarché fait-il une promotion ?)
\item 所有商品打几折？ (Quelle est la réduction sur tous les articles ?)
\item 一公斤芒果多少钱？ (Combien coûte un kilo de mangues ?)
\item 超市几点开门，几点关门？ (À quelle heure le supermarché ouvre-t-il et ferme-t-il ?)
\item 超市在什么地方？ (Où se trouve le supermarché ?)
\end{enumerate}
\textbf{Questions de langue :}
\begin{enumerate}
\item Relève dans le texte une phrase avec la structure \zh{又……又……} et traduis-la en français.
\item Explique en français ce que signifie \zh{买二送一}.
\end{enumerate}
\end{exercice}

\begin{exercice}[Thème et version]
\textbf{Thème.} Traduis en chinois simplifié (avec le pinyin) :
\begin{enumerate}
\item Ce magasin vend des vêtements de très bonne qualité.
\item Aujourd'hui, tous les fruits sont à moins trente pour cent.
\item Nos produits sont à la fois bons et bon marché : venez vite !
\end{enumerate}
\textbf{Version.} Traduis en français :
\begin{enumerate}
\item \zh{这个广告说新手机又便宜又好用。} (Zhège guǎnggào shuō xīn shǒujī yòu piányi yòu hǎo yòng.)
\item \zh{星期天商店里的顾客很多，因为所有商品都打折。} (Xīngqītiān shāngdiàn li de gùkè hěn duō, yīnwèi suǒyǒu shāngpǐn dōu dǎzhé.)
\end{enumerate}
\end{exercice}

\begin{exercice}[Expression écrite : ma publicité]
Tu es le directeur (la directrice) d'un petit magasin à Douala. Rédige en chinois une publicité d'au moins \textbf{60 caractères} pour un produit de ton choix (fruits, vêtements, téléphone, boissons...).
Consignes obligatoires :
\begin{itemize}
\item une accroche (par exemple \zh{好消息！}) ;
\item la structure \zh{又……又……} pour présenter les avantages du produit ;
\item le mot \zh{打折} pour annoncer une réduction ;
\item un appel à l'action final (par exemple \zh{快来吧！} ou \zh{欢迎光临！}).
\end{itemize}
\end{exercice}

\newpage

\coursec{Corrigés des exercices}

\begin{corrige}[Exercice 1 — QCM : le vocabulaire de la publicité]
\begin{enumerate}
\item \textbf{b) la publicité.} \zh{广告} (guǎnggào) signifie « publicité, annonce ». « Le magasin » se dit \zh{商店} (shāngdiàn) et « le prix » \zh{价格} (jiàgé).
\item \textbf{a) faire une réduction.} \zh{打折} (dǎzhé) = « faire une remise ». Attention : \zh{打八折} signifie « vendre à 80 \% du prix », donc « moins vingt pour cent ».
\item \textbf{b) \zh{贵}.} \zh{便宜} (piányi) = « bon marché » ; son contraire est \zh{贵} (guì) = « cher ».
\item \textbf{a) \zh{快来吧！}} (Kuài lái ba !) = « Venez vite ! ». \zh{慢走！} est une formule de politesse (« Allez doucement ! », dit au client qui part) ; \zh{再见！} = « Au revoir ! ».
\item \textbf{b) les marchandises.} \zh{商品} (shāngpǐn) = « marchandises, articles ». « Client » = \zh{顾客} (gùkè) ; « vendeur » = \zh{售货员} (shòuhuòyuán).
\end{enumerate}
\end{corrige}

\begin{corrige}[Exercice 2 — Vrai ou faux]
\begin{enumerate}
\item \textbf{Faux.} Le texte annonce une réduction, mais pas \zh{打五折} (moins cinquante pour cent) sur tous les produits ; il faut citer la phrase exacte du texte, par exemple \zh{所有水果打八折} (« tous les fruits à moins vingt pour cent »).
\item \textbf{Vrai.} Le texte dit \zh{我们的水果又新鲜又便宜} (Wǒmen de shuǐguǒ yòu xīnxiān yòu piányi) : « nos fruits sont à la fois frais et bon marché ».
\item \textbf{Faux.} L'heure de fermeture indiquée dans le texte n'est pas 19 h ; il faut citer la phrase du texte donnant les horaires, par exemple \zh{晚上九点关门} (« fermeture à 21 h »).
\item \textbf{Vrai.} C'est le sens de la formule du texte de type \zh{买二送一} (mǎi èr sòng yī) : « deux achetés, un offert », donc le troisième article ne coûte rien.
\end{enumerate}
\textbf{Méthode :} pour justifier, recopie toujours la phrase exacte du texte entre guillemets, puis traduis-la en français.
\end{corrige}

\begin{corrige}[Exercice 3 — Complète avec le bon mot]
\begin{enumerate}
\item 我在电视上看到了一个新手机的\zh{广告}。 (Wǒ zài diànshì shang kàndào le yí ge xīn shǒujī de guǎnggào.) « J'ai vu à la télévision une publicité pour un nouveau téléphone. »
\item 这家商店的衣服\zh{质量}很好，也不贵。 (Zhè jiā shāngdiàn de yīfu zhìliàng hěn hǎo, yě bú guì.) « Les vêtements de ce magasin sont de très bonne qualité et pas chers. »
\item 今天牛奶\zh{特价}，一升只要五百法郎。 (Jīntiān niúnǎi tèjià, yì shēng zhǐ yào wǔbǎi fǎláng.) « Aujourd'hui le lait est en promotion : cinq cents francs le litre seulement. »
\item 周末商店里\zh{顾客}很多。 (Zhōumò shāngdiàn li gùkè hěn duō.) « Le week-end, il y a beaucoup de clients dans le magasin. »
\item 过新年的时候，很多商店都\zh{打折}。 (Guò Xīnnián de shíhou, hěn duō shāngdiàn dōu dǎzhé.) « Au Nouvel An, beaucoup de magasins font des réductions. »
\end{enumerate}
\end{corrige}

\begin{corrige}[Exercice 4 — Associe le pinyin aux caractères]
\begin{center}
\begin{tabularx}{\linewidth}{|X|X|X|}
\hline
\rowcolor{popblueL} Caractères & Pinyin & Traduction \\
\hline
1. 商店 & c. shāngdiàn & le magasin \\
\hline
2. 广告 & a. guǎnggào & la publicité \\
\hline
3. 打折 & b. dǎzhé & faire une réduction \\
\hline
4. 新鲜 & e. xīnxiān & frais \\
\hline
5. 便宜 & d. piányi & bon marché \\
\hline
\end{tabularx}
\end{center}
\textbf{Point de vigilance :} attention aux tons de \zh{shāngdiàn} (1\textsuperscript{er} ton + 4\textsuperscript{e} ton) et à \zh{piányi}, dont le second caractère se lit au ton neutre.
\end{corrige}

\begin{corrige}[Exercice 5 — Traduis en chinois]
\begin{enumerate}
\item \zh{我们商店卖水果和饮料，今天打八折，快来吧！} (Wǒmen shāngdiàn mài shuǐguǒ hé yǐnliào, jīntiān dǎ bā zhé, kuài lái ba !) — « Moins vingt pour cent » se dit \zh{打八折} (on paie 80 \% du prix).
\item \zh{这双鞋又好看又便宜。} (Zhè shuāng xié yòu hǎokàn yòu piányi.) — Classificateur des chaussures : \zh{双} (shuāng, « la paire »).
\item \zh{这个周末所有商品都特价。} (Zhège zhōumò suǒyǒu shāngpǐn dōu tèjià.) — Variante possible : \zh{所有商品都打折}.
\item \zh{这个手机的质量很好，但是有点儿贵。} (Zhège shǒujī de zhìliàng hěn hǎo, dànshì yǒudiǎnr guì.) — « Un peu cher » = \zh{有点儿贵}.
\end{enumerate}
\end{corrige}

\begin{corrige}[Exercice 6 — Transforme avec 又……又……]
Rappel de la structure : Sujet + \zh{又} + adjectif 1 + \zh{又} + adjectif 2 (les deux adjectifs qualifient le même sujet).
\begin{enumerate}
\item \zh{这个商店又大又干净。} (Zhège shāngdiàn yòu dà yòu gānjìng.) « Ce magasin est à la fois grand et propre. »
\item \zh{这里的苹果又便宜又新鲜。} (Zhèli de píngguǒ yòu piányi yòu xīnxiān.) « Les pommes d'ici sont à la fois bon marché et fraîches. »
\item \zh{这件衣服又好看又舒服。} (Zhè jiàn yīfu yòu hǎokàn yòu shūfu.) « Ce vêtement est à la fois beau et confortable. »
\item \zh{那个广告又有意思又有用。} (Nàge guǎnggào yòu yǒu yìsi yòu yǒuyòng.) « Cette publicité est à la fois intéressante et utile. »
\end{enumerate}
\textbf{Point de vigilance :} ne mets pas \zh{很} entre \zh{又} et l'adjectif ; on dit \zh{又大又干净}, jamais \zh{又很大又很干净} dans cette structure de base.
\end{corrige}

\begin{corrige}[Exercice 7 — Texte à trous]
Publicité complétée :
\begin{enumerate}
\item \zh{从}今天起 (cóng jīntiān qǐ) : « à partir d'aujourd'hui » — structure \zh{从……起}.
\item 所有水果\zh{打}八折 (dǎ bā zhé) : « tous les fruits à moins vingt pour cent » — le verbe \zh{打} est inséparable de \zh{折}.
\item 水果\zh{又}新鲜\zh{又}便宜 : structure \zh{又……又……} (les deux trous).
\item 买\zh{了}还送饮料 (mǎi le hái sòng yǐnliào) : « quand on a acheté, on offre en plus une boisson » — \zh{了} marque l'action achevée.
\item \zh{欢迎}光临 (huānyíng guānglín) : « soyez les bienvenus » — formule de politesse commerciale.
\item 快来\zh{吧} (kuài lái ba) : « venez vite ! » — particule finale d'invitation.
\end{enumerate}
\end{corrige}

\begin{corrige}[Exercice 8 — Remets la publicité dans l'ordre]
Ordre logique : \textbf{C — E — A — D — F — B}.
\begin{enumerate}
\item \textbf{C.} \zh{好消息！新面包到了！} : l'accroche annonce la nouveauté.
\item \textbf{E.} \zh{我们的面包每天早上现做。} : présentation du produit (pain frais fait chaque matin).
\item \textbf{A.} \zh{又好吃，又有营养，老人孩子都喜欢。} : les avantages du produit.
\item \textbf{D.} \zh{一个只要三百法郎。} : le prix.
\item \textbf{F.} \zh{今天买两个打九折。} : la réduction.
\item \textbf{B.} \zh{快来买吧，机会难得！} : l'appel à l'action final.
\end{enumerate}
\textbf{Méthode :} une publicité suit presque toujours le schéma accroche → produit → avantages → prix → promotion → appel à l'action. Repère ce schéma pour réussir ce type d'exercice au Bac.
\end{corrige}

\begin{corrige}[Exercice 9 — Compréhension écrite : une annonce de supermarché]
\textbf{Questions de compréhension (réponses en phrases complètes) :}
\begin{enumerate}
\item \zh{这个星期六和星期天超市有促销。} (Zhège xīngqīliù hé xīngqītiān chāoshì yǒu cùxiāo.) « La promotion a lieu ce samedi et ce dimanche. »
\item \zh{所有商品打八折。} (Suǒyǒu shāngpǐn dǎ bā zhé.) « Tous les articles sont à moins vingt pour cent. »
\item \zh{一公斤芒果只要一千法郎。} (Yì gōngjīn mángguǒ zhǐ yào yìqiān fǎláng.) « Un kilo de mangues coûte mille francs seulement. »
\item \zh{超市早上八点开门，晚上九点关门。} (Chāoshì zǎoshang bā diǎn kāi mén, wǎnshang jiǔ diǎn guān mén.) « Il ouvre à 8 h et ferme à 21 h. »
\item \zh{超市在雅温得市中心，邮局旁边。} (Chāoshì zài Yǎwēndé shì zhōngxīn, yóujú pángbiān.) « Il se trouve au centre-ville de Yaoundé, à côté de la poste. »
\end{enumerate}
\textbf{Questions de langue :}
\begin{enumerate}
\item \zh{水果又新鲜又便宜。} (Shuǐguǒ yòu xīnxiān yòu piányi.) Traduction : « Les fruits sont à la fois frais et bon marché. » (On accepte aussi \zh{质量好，价格低} comme exemple de juxtaposition, mais la structure \zh{又……又……} attendue est celle-ci.)
\item \zh{买二送一} (mǎi èr sòng yī) signifie littéralement « achète deux, on en offre un » : pour deux boissons achetées, la troisième est gratuite.
\end{enumerate}
\textbf{Barème indicatif (sur 10) :} 1 point par réponse de compréhension correcte (5 pts), 2 pts pour la phrase \zh{又……又……} relevée et traduite, 2 pts pour l'explication de \zh{买二送一}, 1 pt pour la correction de la langue.
\textbf{Points de vigilance :} réponds par des phrases complètes avec sujet et verbe ; recopie les chiffres et les horaires sans erreur ; ne traduis pas \zh{打八折} par « moins quatre-vingts pour cent » (c'est « moins vingt pour cent »).
\end{corrige}

\begin{corrige}[Exercice 10 — Thème et version]
\textbf{Thème :}
\begin{enumerate}
\item \zh{这家商店卖质量很好的衣服。} (Zhè jiā shāngdiàn mài zhìliàng hěn hǎo de yīfu.) — Le complément \zh{质量很好} précède le nom \zh{衣服} grâce à \zh{的}.
\item \zh{今天所有水果打七折。} (Jīntiān suǒyǒu shuǐguǒ dǎ qī zhé.) — « Moins trente pour cent » = \zh{打七折} (on paie 70 \%).
\item \zh{我们的商品又好又便宜，快来吧！} (Wǒmen de shāngpǐn yòu hǎo yòu piányi, kuài lái ba !)
\end{enumerate}
\textbf{Version :}
\begin{enumerate}
\item \zh{这个广告说新手机又便宜又好用。} : « Cette publicité dit que le nouveau téléphone est à la fois bon marché et pratique (agréable à utiliser). »
\item \zh{星期天商店里的顾客很多，因为所有商品都打折。} : « Le dimanche, il y a beaucoup de clients dans le magasin, parce que tous les articles sont en réduction. »
\end{enumerate}
\textbf{Barème indicatif (sur 10) :} 2 pts par phrase (lexique 1 pt, grammaire et ordre des mots 1 pt).
\textbf{Points de vigilance :} en thème, attention au calcul de la réduction (\zh{打七折} = moins 30 \%) et à la place de \zh{都} (après le sujet, avant le verbe) ; en version, \zh{好用} se traduit par « pratique, facile à utiliser », pas par « bon à utiliser » mot à mot.
\end{corrige}

\begin{corrige}[Exercice 11 — Expression écrite : ma publicité]
\textbf{Proposition de rédaction modèle (environ 75 caractères) :}

\zh{好消息！杜阿拉阳光商店大促销！我们的芒果又大又甜，又新鲜又便宜。今天所有水果打七折，一公斤只要八百法郎。买两公斤还送一瓶饮料！机会难得，欢迎光临，快来买吧！}\\
(Hǎo xiāoxi ! Dù'ālā Yángguāng shāngdiàn dà cùxiāo ! Wǒmen de mángguǒ yòu dà yòu tián, yòu xīnxiān yòu piányi. Jīntiān suǒyǒu shuǐguǒ dǎ qī zhé, yì gōngjīn zhǐ yào bābǎi fǎláng. Mǎi liǎng gōngjīn hái sòng yì píng yǐnliào ! Jīhuì nándé, huānyíng guānglín, kuài lái mǎi ba !)\\
« Bonne nouvelle ! Grande promotion au magasin Soleil de Douala ! Nos mangues sont à la fois grosses et sucrées, fraîches et bon marché. Aujourd'hui tous les fruits à moins trente pour cent : huit cents francs le kilo seulement. Pour deux kilos achetés, une bouteille de boisson offerte ! Occasion rare, soyez les bienvenus, venez vite acheter ! »

\textbf{Barème indicatif (sur 10) :} accroche présente (1 pt) ; structure \zh{又……又……} correcte (2 pts) ; \zh{打折} bien employé avec le bon chiffre (2 pts) ; appel à l'action final (1 pt) ; longueur minimale de 60 caractères respectée (1 pt) ; correction générale de la langue et des caractères (3 pts).

\textbf{Points de vigilance :}
\begin{itemize}
\item Vérifie le calcul de la réduction : \zh{打八折} = moins 20 \%, \zh{打七折} = moins 30 \%, \zh{打九折} = moins 10 \%.
\item Pas de \zh{很} entre \zh{又} et l'adjectif : \zh{又大又甜}.
\item Utilise le bon classificateur : \zh{一公斤芒果}, \zh{一瓶饮料}, \zh{一件衣服}, \zh{一部手机}.
\item Termine par une formule commerciale : \zh{欢迎光临！} ou \zh{快来吧！}.
\end{itemize}
\end{corrige}

\newpage

\coursec{Fiche bilan}

\begin{popfigure}
\begin{center}\tcbox[colback=popblue,colframe=popblue,coltext=white,arc=9pt,boxsep=4pt,left=6pt,right=6pt]{\bfseries\small 广告 \emph{guǎnggào} La publicité}\end{center}
\vspace{-2pt}
\begin{tcbraster}[raster columns=3,raster equal height=rows,raster column skip=6pt,raster row skip=6pt,enhanced,blankest]
\begin{tcolorbox}[enhanced,colback=poporange,colframe=poporange,coltext=white,boxrule=0.9pt,arc=3pt,left=4pt,right=4pt,top=3pt,bottom=3pt,boxsep=1pt,halign=left]\bfseries\scriptsize Lexique 产品 顾客 价格 打折\end{tcolorbox}
\begin{tcolorbox}[enhanced,colback=popgreen,colframe=popgreen,coltext=white,boxrule=0.9pt,arc=3pt,left=4pt,right=4pt,top=3pt,bottom=3pt,boxsep=1pt,halign=left]\bfseries\scriptsize Texte support 欢迎光临！ Lire et comprendre\end{tcolorbox}
\begin{tcolorbox}[enhanced,colback=poppurple,colframe=poppurple,coltext=white,boxrule=0.9pt,arc=3pt,left=4pt,right=4pt,top=3pt,bottom=3pt,boxsep=1pt,halign=left]\bfseries\scriptsize Structures 为了 + but 只要…就…\end{tcolorbox}
\begin{tcolorbox}[enhanced,colback=poppink,colframe=poppink,coltext=white,boxrule=0.9pt,arc=3pt,left=4pt,right=4pt,top=3pt,bottom=3pt,boxsep=1pt,halign=left]\bfseries\scriptsize Radicaux 讠(parole) 宀(toit) 亻(homme)\end{tcolorbox}
\begin{tcolorbox}[enhanced,colback=popgold,colframe=popgold,coltext=white,boxrule=0.9pt,arc=3pt,left=4pt,right=4pt,top=3pt,bottom=3pt,boxsep=1pt,halign=left]\bfseries\scriptsize Compréhension Questions 谁？什么？为什么？\end{tcolorbox}
\begin{tcolorbox}[enhanced,colback=popmuted,colframe=popmuted,coltext=white,boxrule=0.9pt,arc=3pt,left=4pt,right=4pt,top=3pt,bottom=3pt,boxsep=1pt,halign=left]\bfseries\scriptsize Production Ma publicité 4 à 6 phrases\end{tcolorbox}
\begin{tcolorbox}[enhanced,colback=popblue!60,colframe=popblue!60,coltext=white,boxrule=0.9pt,arc=3pt,left=4pt,right=4pt,top=3pt,bottom=3pt,boxsep=1pt,halign=left]\bfseries\scriptsize Culture Marchés du Cameroun et publicité en Chine\end{tcolorbox}
\end{tcbraster}

\legende{Carte mentale de la leçon 30 : la publicité en chinois.}
\end{popfigure}

\begin{bilan}
\textbf{1. Lexique clé de la publicité}
\begin{itemize}
  \item \zh{广告} \emph{guǎnggào} : publicité ; \zh{宣传} \emph{xuānchuán} : faire de la publicité, promouvoir
  \item \zh{产品} \emph{chǎnpǐn} : produit ; \zh{质量} \emph{zhìliàng} : qualité ; \zh{品牌} \emph{pǐnpái} : marque
  \item \zh{顾客} \emph{gùkè} : client ; \zh{价格} \emph{jiàgé} : prix ; \zh{便宜} \emph{piányi} : bon marché
  \item \zh{打折} \emph{dǎzhé} : faire une réduction ; \zh{优惠} \emph{yōuhuì} : avantage, promotion
  \item \zh{商店} \emph{shāngdiàn} : magasin ; \zh{欢迎光临} \emph{huānyíng guānglín} : soyez les bienvenus !
\end{itemize}

\textbf{2. Structures à connaître par cœur}
\begin{center}
\begin{tabularx}{\linewidth}{|X|X|X|}
\hline
\rowcolor{popblueL} Structure & Exemple & Traduction \\
\hline
为了 + but & \zh{为了吸引顾客，我们打折。} \emph{Wèile xīyǐn gùkè, wǒmen dǎzhé.} & Pour attirer les clients, nous faisons des réductions. \\
\hline
只要 … 就 … & \zh{只要质量好，顾客就会来。} \emph{Zhǐyào zhìliàng hǎo, gùkè jiù huì lái.} & Pourvu que la qualité soit bonne, les clients viendront. \\
\hline
越 … 越 … & \zh{价格越便宜，买的人越多。} \emph{Jiàgé yuè piányi, mǎi de rén yuè duō.} & Plus le prix est bas, plus il y a d'acheteurs. \\
\hline
不 但 … 而且 … & \zh{这个产品不但便宜，而且质量好。} \emph{Zhège chǎnpǐn búdàn piányi, érqiě zhìliàng hǎo.} & Ce produit est non seulement bon marché, mais aussi de bonne qualité. \\
\hline
\end{tabularx}
\end{center}

\textbf{3. Radicaux et écriture}
\begin{itemize}
  \item \zh{广} (clé 广 « abri ») : le point d'abord, puis l'horizontale, puis la verticale oblique de gauche à droite.
  \item \zh{告} : la clé 牛 « bœuf » en haut, la clé 口 « bouche » en bas ; on écrit de haut en bas.
  \item \zh{价} : clé 亻 « homme » à gauche d'abord, puis 介 à droite (gauche avant droite).
  \item \zh{客} : clé 宀 « toit » en haut, puis 夂, puis 口 ; toujours de haut en bas.
\end{itemize}

\textbf{4. Méthodes express}
\begin{itemize}
  \item \cle{Comprendre une publicité} : repérer le produit (\zh{产品}), le prix (\zh{价格}), l'argument de vente (\zh{质量好、价格便宜}) et l'appel à l'action (\zh{欢迎光临！快来买！}).
  \item \cle{Rédiger une publicité} : 1) salutation d'accroche ; 2) nom du produit ; 3) deux qualités avec \zh{不但…而且…} ; 4) prix ou réduction avec \zh{打折} ; 5) appel final \zh{欢迎大家来！}.
  \item \cle{Répondre aux questions} : recopier la structure de la question et remplacer l'information, sans changer l'ordre des mots.
\end{itemize}
\end{bilan}

\begin{aretenir}[Checklist avant le Bac]
\begin{itemize}
  \item $\square$ Je sais écrire et prononcer \zh{广告、产品、顾客、价格、打折、质量} avec les bons tons.
  \item $\square$ Je connais la formule d'accueil \zh{欢迎光临！} \emph{huānyíng guānglín}.
  \item $\square$ Je sais exprimer le but avec \zh{为了} + proposition.
  \item $\square$ Je sais construire une condition avec \zh{只要…就…} sans oublier \zh{就}.
  \item $\square$ Je sais comparer avec \zh{越…越…} (plus… plus…).
  \item $\square$ Je sais enchaîner deux qualités avec \zh{不但…而且…}.
  \item $\square$ Je place correctement le complément de prix : \zh{这个产品价格很便宜} et non l'ordre français.
  \item $\square$ Je connais les radicaux 广 (abri), 亻 (homme), 宀 (toit), 讠 (parole) et leur sens.
  \item $\square$ Je respecte l'ordre des traits : haut avant bas, gauche avant droite.
  \item $\square$ Je sais rédiger une publicité de 4 à 6 phrases avec accroche, produit, qualités, prix et appel final.
  \item $\square$ Je réponds aux questions de compréhension en reprenant la structure de la phrase du texte.
  \item $\square$ Je vérifie la ponctuation chinoise pleine chasse ：，。！ dans ma production.
\end{itemize}
\end{aretenir}

\findecours

\end{document}
